
\documentclass[
  10pt,
  a4paper,
  twoside,
  open=any,
  parskip=half,
  headings=big
]{scrreprt}

\usepackage[a4paper,inner=27mm,outer=23mm,top=24mm,bottom=26mm,headsep=8mm,footskip=12mm]{geometry}
\usepackage{amsmath}
\usepackage{libertinus-otf}
\usepackage{fontspec}
\setmonofont[Scale=0.82]{DejaVu Sans Mono}
\usepackage{microtype}
\usepackage{xcolor}
\usepackage{graphicx}
\usepackage{booktabs}
\usepackage{longtable}
\usepackage{tabularx}
\usepackage{array}
\usepackage{ragged2e}
\usepackage{enumitem}
\usepackage{listings}
\usepackage[most]{tcolorbox}
\usepackage{tikz}
\usetikzlibrary{arrows.meta,positioning,fit,calc,backgrounds,shapes.geometric}
\usepackage{caption}
\usepackage{subcaption}
\usepackage{scrlayer-scrpage}
\usepackage{bookmark}
\usepackage{xurl}
\usepackage{hyperref}
\usepackage[capitalise,noabbrev]{cleveref}
\usepackage{csquotes}
\usepackage{etoolbox}
\usepackage{needspace}
\usepackage{multicol}
\usepackage{pdflscape}

\definecolor{navy}{HTML}{172A46}
\definecolor{blue}{HTML}{285F8F}
\definecolor{sky}{HTML}{EAF3F8}
\definecolor{teal}{HTML}{2D766A}
\definecolor{mint}{HTML}{E9F4F0}
\definecolor{gold}{HTML}{9A6A10}
\definecolor{sand}{HTML}{F8F1E2}
\definecolor{brick}{HTML}{9C443E}
\definecolor{rose}{HTML}{F8EDEC}
\definecolor{ink}{HTML}{20252C}
\definecolor{muted}{HTML}{5E6874}
\definecolor{rulegray}{HTML}{C8D0D8}
\definecolor{codebg}{HTML}{F5F7F9}
\definecolor{lightgray}{HTML}{F1F3F5}
\definecolor{darkgreen}{HTML}{255C4E}

\hypersetup{
  colorlinks=true,
  linkcolor=navy,
  citecolor=teal,
  urlcolor=blue,
  pdfauthor={Repository walkthrough},
  pdftitle={ProveIt: A Detailed Codebase Walkthrough},
  pdfsubject={Architecture, proof boundaries, build topology, and project guide},
  pdfkeywords={ProveIt, Lean 4, Rocq, Coq, formal mathematics, codebase walkthrough}
}
\urlstyle{same}

\addtokomafont{disposition}{\sffamily\color{navy}}
\addtokomafont{chapter}{\bfseries}
\addtokomafont{section}{\bfseries}
\addtokomafont{subsection}{\bfseries}
\setkomafont{descriptionlabel}{\sffamily\bfseries\color{navy}}
\setkomafont{captionlabel}{\sffamily\bfseries\color{navy}}
\setkomafont{caption}{\small}
\setkomafont{pageheadfoot}{\small\sffamily\color{muted}}
\setkomafont{pagenumber}{\sffamily\color{navy}}
\RedeclareSectionCommand[beforeskip=0pt,afterskip=1.2\baselineskip]{chapter}
\RedeclareSectionCommand[beforeskip=1.8\baselineskip,afterskip=.65\baselineskip]{section}
\RedeclareSectionCommand[beforeskip=1.35\baselineskip,afterskip=.45\baselineskip]{subsection}

\clearpairofpagestyles
\automark[section]{chapter}
\ihead{\headmark}
\ohead{ProveIt walkthrough}
\ifoot{\texttt{675d0a2f}}
\ofoot{\pagemark}
\KOMAoptions{headsepline=0.35pt}
\pagestyle{scrheadings}

\setlist[itemize]{topsep=.35\baselineskip,itemsep=.18\baselineskip,parsep=0pt,leftmargin=1.7em}
\setlist[enumerate]{topsep=.35\baselineskip,itemsep=.22\baselineskip,parsep=0pt,leftmargin=1.9em}
\setlist[description]{topsep=.4\baselineskip,itemsep=.25\baselineskip,style=nextline,leftmargin=1.8em,labelindent=0pt}

\newcolumntype{L}[1]{>{\RaggedRight\arraybackslash}p{#1}}
\newcolumntype{Y}{>{\RaggedRight\arraybackslash}X}
\newcolumntype{C}[1]{>{\Centering\arraybackslash}p{#1}}

\lstdefinelanguage{Lean}{
  morekeywords={import,namespace,end,open,section,noncomputable,def,abbrev,structure,inductive,theorem,lemma,example,where,by,match,with,if,then,else,let,in,have,show,exact,refine,constructor,cases,rcases,fun,forall,exists,deriving,instance,class,private,protected,set_option},
  sensitive=true,
  morecomment=[l]{--},
  morecomment=[s]{/-}{-/},
  morestring=[b]"
}
\lstdefinelanguage{Rocq}{
  morekeywords={Require,Import,Export,Module,End,Section,Variable,Variables,Hypothesis,Hypotheses,Definition,Fixpoint,Inductive,Theorem,Lemma,Proof,Qed,Defined,Admitted,Print,Assumptions,Check,forall,exists,match,with,end,fun,let,in},
  sensitive=true,
  morecomment=[s]{(*}{*)},
  morestring=[b]"
}
\lstdefinelanguage{PowerShell}{
  morekeywords={param,function,begin,process,end,if,else,elseif,foreach,ForEach-Object,Where-Object,throw,return,Set-Location,Get-Content},
  sensitive=false,
  morecomment=[l]{\#},
  morestring=[b]",
  morestring=[b]'
}
\lstset{
  basicstyle=\ttfamily\footnotesize,
  keywordstyle=\bfseries\color{navy},
  commentstyle=\itshape\color{muted},
  stringstyle=\color{darkgreen},
  columns=fullflexible,
  keepspaces=true,
  breaklines=true,
  breakatwhitespace=false,
  showstringspaces=false,
  upquote=true,
  frame=single,
  framerule=.35pt,
  rulecolor=\color{rulegray},
  backgroundcolor=\color{codebg},
  xleftmargin=.7em,
  xrightmargin=.7em,
  framexleftmargin=.45em,
  framexrightmargin=.45em,
  aboveskip=.7\baselineskip,
  belowskip=.7\baselineskip
}

\tcbset{
  enhanced,
  boxrule=.65pt,
  arc=1.2mm,
  left=2.2mm,
  right=2.2mm,
  top=1.5mm,
  bottom=1.5mm,
  before skip=.75\baselineskip,
  after skip=.75\baselineskip,
  fonttitle=\sffamily\bfseries
}
\newtcolorbox{overviewbox}[1][]{colback=sky,colframe=blue,title={Orientation},#1}
\newtcolorbox{boundarybox}[1][]{colback=sand,colframe=gold,title={Proof boundary},#1}
\newtcolorbox{statusbox}[1][]{colback=mint,colframe=teal,title={Status},#1}
\newtcolorbox{gotchabox}[1][]{colback=rose,colframe=brick,title={Gotcha},#1}
\newtcolorbox{readingbox}[1][]{colback=lightgray,colframe=muted,title={Reading route},#1}
\newtcolorbox{sourcebox}[1][]{colback=white,colframe=rulegray,title={Primary files},#1}

\newcommand{\CommitFull}{675d0a2ff0fc624fba199e006049016ba806f40a}
\newcommand{\CommitShort}{675d0a2f}
\newcommand{\SnapshotDate}{13 August 2026}
\newcommand{\RepoWeb}{https://github.com/VladimirReshetnikov/ProveIt}
\newcommand{\RepoBlobBase}{https://github.com/VladimirReshetnikov/ProveIt/blob/\CommitFull/}
\newcommand{\RepoTreeBase}{https://github.com/VladimirReshetnikov/ProveIt/tree/\CommitFull/}
\newcommand{\RepoFile}[1]{%
  \href{\RepoBlobBase\detokenize{#1}}{\nolinkurl{#1}}}
\newcommand{\RepoDir}[1]{%
  \href{\RepoTreeBase\detokenize{#1}}{\nolinkurl{#1}}}
\newcommand{\code}[1]{\texttt{\detokenize{#1}}}
\newcommand{\bcode}[1]{{\ttfamily\footnotesize\nolinkurl{#1}}}
\newcommand{\Lean}{Lean 4}
\newcommand{\Rocq}{Rocq/Coq}
\newcommand{\Mathlib}{Mathlib}
\newcommand{\term}[1]{\textsf{#1}}
\newcommand{\project}[1]{\textsf{\bfseries #1}}
\newcommand{\yes}{\textcolor{teal}{\textbf{yes}}}
\newcommand{\partialstatus}{\textcolor{gold}{\textbf{partial}}}
\newcommand{\nostatus}{\textcolor{brick}{\textbf{no}}}
\newcommand{\conditional}{\textcolor{gold}{\textbf{conditional}}}
\newcommand{\kernelchecked}{\textcolor{teal}{\textbf{kernel-checked}}}
\newcommand{\generated}{\textcolor{blue}{\textbf{generated}}}

\newcommand{\blankpage}{%
  \newpage
  \thispagestyle{empty}
  \mbox{}
  \newpage
}

\makeatletter\renewcommand*{\@pnumwidth}{2.2em}\makeatother
\setcounter{tocdepth}{2}
\setcounter{secnumdepth}{3}
\emergencystretch=3em
\clubpenalty=10000
\widowpenalty=10000
\displaywidowpenalty=10000

\begin{document}

\begin{titlepage}
\thispagestyle{empty}
\begin{tikzpicture}[remember picture,overlay]
  \fill[navy] (current page.north west) rectangle ([yshift=-64mm]current page.north east);
  \fill[blue] ([yshift=-64mm]current page.north west) rectangle ([yshift=-68mm]current page.north east);
  \draw[rulegray,line width=.6pt] ([xshift=22mm,yshift=22mm]current page.south west)
    -- ([xshift=-22mm,yshift=22mm]current page.south east);
  \node[anchor=north west,text width=164mm] at ([xshift=22mm,yshift=-22mm]current page.north west) {
    {\sffamily\bfseries\fontsize{38}{42}\selectfont\color{white} ProveIt}\\[3mm]
    {\sffamily\fontsize{19}{24}\selectfont\color{white} A detailed codebase walkthrough}\\[4mm]
    {\sffamily\large\color{white!88} Architecture, proof boundaries, build topology,\\
      and a subject-by-subject guide}
  };
  \node[anchor=north west,text width=52mm,draw=blue!35,fill=sky,rounded corners=1.5mm,
        inner sep=4mm] at ([xshift=22mm,yshift=-91mm]current page.north west) {
    {\sffamily\bfseries\color{navy} Lean workspace}\\[1.2mm]
    Facades, semantic cores, executable decision procedures, and kernel-reduced certificates.
  };
  \node[anchor=north west,text width=52mm,draw=teal!35,fill=mint,rounded corners=1.5mm,
        inner sep=4mm] at ([xshift=80mm,yshift=-91mm]current page.north west) {
    {\sffamily\bfseries\color{navy} Rocq federation}\\[1.2mm]
    Independent ports, complementary developments, vendored proof libraries, and assumption audits.
  };
  \node[anchor=north west,text width=52mm,draw=gold!45,fill=sand,rounded corners=1.5mm,
        inner sep=4mm] at ([xshift=138mm,yshift=-91mm]current page.north west) {
    {\sffamily\bfseries\color{navy} Certificate pipelines}\\[1.2mm]
    Generated candidates, literal data modules, small checkers, and explicit trust boundaries.
  };
  \node[anchor=south west,text width=164mm] at ([xshift=22mm,yshift=31mm]current page.south west) {
    {\sffamily\large\color{navy} Snapshot}\\[1mm]
    \texttt{\CommitFull}\\
    \SnapshotDate\\[4mm]
    {\sffamily\small\color{muted} All repository hyperlinks in this document are pinned to this commit.}
  };
\end{tikzpicture}
\end{titlepage}

\cleardoublepage
\thispagestyle{plain}
\begin{center}
{\sffamily\bfseries\Large About this walkthrough}
\end{center}

This document explains \emph{how the ProveIt repository is put together}: what
each major subtree proves, where its public surface lives, how its files depend
on one another, where computation enters a proof, how the two proof assistants
relate, and how to build, audit, debug, or extend a development without losing
track of its logical status.

It is pinned to commit \texttt{\CommitFull}.  That is important for a repository
whose active branches can change both mathematical status and physical layout
quickly.  Paths, module inventories, and statements below refer to that exact
tree; every path rendered in blue is a clickable commit-pinned source link.

The inspection used GitHub's repository and Git-object APIs to read the exact
commit, its trees, manifests, selected implementation modules, audit modules,
support programs, and project READMEs.  A complete checkout and full
multi-hour/multi-toolchain build of ProveIt was \emph{not} performed in the
environment used to prepare this text.  Consequently, this is an architectural
and source walkthrough, not an independent re-certification of every theorem.
The PDF itself was compiled and visually checked from the accompanying
\LaTeX{} source.

\begin{overviewbox}
The repository is easiest to understand along three independent axes.

\begin{enumerate}
  \item \textbf{Mathematical subject}: algebra, analysis, combinatorics,
        computability, logic, number theory, and set theory.
  \item \textbf{Proof shape}: a short direct theorem, a from-scratch semantic
        development, a finite certificate, a generated algebraic certificate
        stack, or a port/coverage project.
  \item \textbf{Evidence boundary}: kernel-checked proof terms, kernel
        computation, VM/native computation, generated candidate data, external
        source transcription, or an explicitly retained hypothesis.
\end{enumerate}

Folder names mostly encode the first axis.  Facades, audit modules, status
tables, and file naming conventions expose the other two.
\end{overviewbox}

\section*{How to use the document}

Readers new to the repository should read \cref{chap:orientation,chap:build-topology,chap:trust}
before choosing a subject chapter.  A contributor modifying one project can
jump directly to its subject section and then use
\cref{chap:working,chap:extension,chap:debugging}.  Readers evaluating theorem
status should begin with \cref{chap:trust} and consult the status/frontier
catalogue in \cref{chap:frontiers} rather than inferring status from a theorem
name or file size.

The terms \emph{Coq} and \emph{Rocq} both appear because the repository is in a
transition period reflected in filenames, logical roots, packages, and command
names.  This document uses \emph{Rocq/Coq} for the prover family, but preserves
literal names such as \code{_CoqProject}, \code{coqc}, and \code{Coq/} when
describing repository artifacts.

\cleardoublepage
\pdfbookmark[0]{Contents}{contents}
\tableofcontents
\cleardoublepage
\listoffigures
\addcontentsline{toc}{chapter}{List of figures}
\listoftables
\addcontentsline{toc}{chapter}{List of tables}

\part{The repository as a system}

\chapter{Orientation: what is ProveIt?}
\label{chap:orientation}

\section{A formal-mathematics monorepo, not a single library}

ProveIt is a monorepo of formalized mathematics.  It has a shared Lean
workspace and a repository-wide Rocq build description, but it does not have
one domain model, one executable entry point, or one uniform proof strategy.
A two-line propositional theorem, a from-scratch completeness proof, a verified
decision procedure, a finite geometric certificate, and a generated
multi-thousand-module invariant-theory calculation can all be first-class
projects in the same tree.

The top-level organization is mathematical:

\begin{center}
\begin{tikzpicture}[
  node distance=6mm and 7mm,
  dombox/.style={draw=blue!45,fill=sky,rounded corners=1.2mm,
                 minimum width=33mm,minimum height=10mm,
                 align=center,font=\sffamily\small\bfseries},
  root/.style={draw=navy,fill=navy,text=white,rounded corners=1.5mm,
               minimum width=42mm,minimum height=12mm,
               align=center,font=\sffamily\bfseries},
  >=Latex
]
\node[root] (root) {ProveIt};
\node[dombox,below left=13mm and 55mm of root] (alg) {Algebra};
\node[dombox,right=of alg] (ana) {Analysis};
\node[dombox,right=of ana] (comb) {Combinatorics};
\node[dombox,right=of comb] (comp) {Computability};
\node[dombox,below=of alg] (logic) {Logic};
\node[dombox,right=of logic] (nt) {Number theory};
\node[dombox,right=of nt] (st) {Set theory};
\node[dombox,right=of st] (infra) {Infrastructure};
\foreach \n in {alg,ana,comb,comp,logic,nt,st,infra}
  \draw[->,rulegray] (root) -- (\n);
\end{tikzpicture}
\captionsetup{hypcap=false}%
\captionof{figure}{The physical top level is a subject taxonomy.  Proof shape and trust boundary cut across it.}
\label{fig:subject-taxonomy}
\end{center}

The root Lean facade, \RepoFile{ProveIt.lean}, creates a convenient broad import
surface over selected projects.  The root Lake file,
\RepoFile{lakefile.toml}, registers many more libraries and gives some
heavyweight and audit libraries explicit default-target status.  On the Rocq
side, \RepoFile{_CoqProject} is both a logical-path map and an ordered source
registry.  These three files answer different questions:

\begin{longtable}{L{.23\textwidth}L{.69\textwidth}}
\caption{Three different notions of ``what belongs to ProveIt.''}
\label{tab:three-surfaces}\\
\toprule
Artifact & Meaning \\
\midrule
\endfirsthead
\toprule
Artifact & Meaning \\
\midrule
\endhead
\RepoFile{ProveIt.lean} &
The intentionally broad \emph{user-facing Lean import}.  It omits several
expensive or status-sensitive modules even when their libraries are registered
and buildable.\\
\RepoFile{lakefile.toml} &
The Lean workspace registry: logical library names, physical source roots,
facade defaults, and additional default targets.\\
\RepoFile{_CoqProject} &
Rocq logical roots, vendored roots, and the source set whose dependency graph
is handed to \code{rocq makefile}.\\
\bottomrule
\end{longtable}

A frequent mistake is to treat these surfaces as synonyms.  For example,
\code{import ProveIt} does not import the expensive Busy Beaver classification
certificates, while the corresponding project is fully present and has
separate build targets.  Conversely, a project-local Lean directory may have a
standalone \code{lakefile.toml} even though the root workspace also maps the
same logical library.

\section{Snapshot and scale}

The inspected snapshot is commit
\texttt{\CommitFull}, dated \SnapshotDate.  Its GitHub language accounting is
shown in \cref{tab:languages}.  Byte counts are useful for estimating build and
storage shape, but not for measuring mathematical originality: generated
literal certificates and vendored sources dominate raw volume.

\begin{table}[htbp]
\centering
\caption{GitHub language accounting at the pinned commit.}
\label{tab:languages}
\begin{tabular}{lr r}
\toprule
Language & Bytes & Share \\
\midrule
Lean & 139,510,311 & 77.60\% \\
Rocq Prover & 39,617,616 & 22.04\% \\
Python & 380,455 & 0.21\% \\
\TeX & 198,247 & 0.11\% \\
Wolfram Language & 51,635 & 0.03\% \\
PowerShell & 19,257 & 0.01\% \\
C & 2,195 & $<0.01$\% \\
Shell & 2,193 & $<0.01$\% \\
\midrule
Total & 179,781,909 & 100\% \\
\bottomrule
\end{tabular}
\end{table}

GitHub code search reports roughly 30,464 Lean files and 2,968 Rocq
\code{.v} files in this snapshot.  Those numbers are startling until one
separates the following populations:

\begin{itemize}
  \item handwritten semantic and theorem modules;
  \item generated data and certificate shards, especially under
        \project{PolynomialFormulas};
  \item copied or vendored proof libraries under \RepoDir{lib};
  \item audit facades and compatibility facades;
  \item independent Lean and Rocq versions of the same mathematical project.
\end{itemize}

The repository therefore has a \emph{small conceptual core relative to its
physical size}.  A productive walkthrough identifies the semantic core and the
certificate checker before opening the largest files.

\section{Top-level map}

\begin{longtable}{L{.18\textwidth}L{.52\textwidth}L{.22\textwidth}}
\caption{Top-level directories at the pinned commit.}
\label{tab:top-level}\\
\toprule
Directory & Main contents & Typical proof shape \\
\midrule
\endfirsthead
\toprule
Directory & Main contents & Typical proof shape \\
\midrule
\endhead
\RepoDir{Algebra} &
Jacobian counterexamples; low-degree polynomial formulas; Abel--Ruffini,
quintic and sextic radical decision procedures; Lazard/Dummit certificate
stacks. &
Direct algebra, imported theory, and very large generated certificates.\\
\RepoDir{Analysis} &
Exact trigonometric and exponential identities, including the tiny exponent
tower floor calculation. &
Short analytic developments with exact rational bounds.\\
\RepoDir{Combinatorics} &
Expression enumeration, polyomino growth, and squared-square dissections. &
Finite certificates, explicit injections, and asymptotic arguments.\\
\RepoDir{Computability} &
Busy Beaver, combinatory logic, and Turing degrees. &
Executable semantics, simulation theorems, classification certificates,
and cross-prover coverage.\\
\RepoDir{Logic} &
Propositional/equational microprojects, first-order completeness,
Foundation ports, modal logic, interpretability, Peano arithmetic, and
Presburger arithmetic. &
From-scratch calculi, arithmetization, ports, decision procedures, and
status-heavy metamathematics.\\
\RepoDir{NumberTheory} &
Fermat exponent four, integer sums, Calkin--Wilf enumeration, and a PA
statement of the Riemann hypothesis. &
Small wrappers plus independent paired proofs.\\
\RepoDir{SetTheory} &
A lightweight ZF development, closure axiomatization, and bounded-complexity
consistency for ZFC. &
Internal model mathematics and partial-satisfaction constructions.\\
\RepoDir{Oeis} &
An auxiliary/legacy home for sequence material such as A290268. &
Sequence-specific formalization.\\
\RepoDir{Shenanigans} &
A redirect to a separate repository rather than an active development. &
Documentation only.\\
\RepoDir{Tools} &
The Rocq 9.2 compatibility overlay. &
Toolchain engineering, not theorem content.\\
\RepoDir{lib} &
Vendored and submodule dependencies, with provenance and licensing
information. &
Read-only/reference or locally hardened third-party proof code.\\
\bottomrule
\end{longtable}

The domain READMEs---for example \RepoFile{Algebra/README.md},
\RepoFile{Computability/README.md}, and \RepoFile{Logic/README.md}---are
useful tables of contents, but project READMEs often contain the more
important information: exact theorem status, intentional nonclaims, build
commands, trust boundaries, and module-by-module roles.

\section{Five recurring species of project}

\subsection{The microproject}

A microproject has a facade, one or a few proof modules, and an audit.  Examples
include \project{QuantifierCommutation}, \project{ShefferStroke}, and several
analysis and number-theory topics.  These are the best places to learn local
naming and audit conventions because their dependency graph fits in one view.

\subsection{The hand-written semantic stack}

Here the repository defines syntax, semantics, a proof system, metatheory, a
public textbook theorem, and an assumption audit.  The independent
first-order completeness project is the canonical example:
\code{Fol} $\to$ \code{Calculus} $\to$ \code{Completeness}
$\to$ \code{ClassicalCompleteness}/\code{Compactness} $\to$ \code{Audit}.

\subsection{The finite certificate project}

A finite object is supplied as data and checked by a small executable
predicate whose correctness is proved.  Duijvestijn's squared-square placement
and the Busy Beaver small-state classifications use this shape.  The crucial
question is not ``was a computer used?'' but ``which computation is replayed
inside the proof assistant, and by which evaluator?''

\subsection{The generated algebraic certificate pipeline}

The largest Lazard and Dummit developments go further.  External programs
construct and shard enormous exact polynomial identities; generated Lean or
Rocq modules contain literal normal forms; small certificate modules prove
that a formally defined substitution, multiplication, or normalization
operation accepts them.  The generator is replaceable.  The formal checker and
the semantic theorem connecting acceptance to the intended mathematics are
the proof boundary.

\subsection{The port and coverage project}

Some projects independently recreate a theorem family in both proof
assistants.  Others, notably \project{FoundationPort} and \project{Modal},
track a pinned source library and record which declarations have been ported,
partially ported, or remain outside scope.  Here a status ledger is part of the
architecture: a green module build is not automatically a parity claim.

\section{The fastest first inspection}

\begin{readingbox}
For an unfamiliar project \code{X}, use this sequence.

\begin{enumerate}
  \item Read the nearest project \code{README.md}, especially headings named
        \emph{Status}, \emph{Trust boundary}, \emph{Deliberate boundary},
        \emph{Layout}, and \emph{Building}.
  \item Open the facade \code{Lean/X.lean} or its equivalent.  Its import list
        identifies the public dependency cut.
  \item Open \code{Audit.lean} or \code{Audit.v}.  The declarations audited
        there are usually a better public-API inventory than filename size.
  \item Follow one headline theorem backward until definitions stop being
        project-specific.
  \item If names contain \code{Data}, \code{Certificate}, \code{Shard},
        \code{Tail}, or \code{Support}, identify the checker before reading the
        literal data.
  \item Compare the root build target with any project-local Lake or Rocq
        instructions.  They can differ intentionally.
\end{enumerate}
\end{readingbox}

\begin{sourcebox}
\RepoFile{README.md};
\RepoFile{ProveIt.lean};
\RepoFile{lakefile.toml};
\RepoFile{_CoqProject};
\RepoFile{lean-toolchain};
\RepoFile{lake-manifest.json};
\RepoFile{.gitmodules};
\RepoFile{.gitattributes}.
\end{sourcebox}

\chapter{Import surfaces and build topology}
\label{chap:build-topology}

\section{The root Lean facade}

The root module is intentionally declarative: it contains imports, not a
repository-wide namespace or executable.  At the pinned commit its public
surface is equivalent to the following list.

\begin{lstlisting}[language=Lean,caption={Root facade import surface, normalized for presentation.},label={lst:root-imports}]
import JacobianConjecture
import PolynomialFormulas
import TrigonometricIdentities
import ExponentialIdentities
import DiophantineEquations
import IntegerSums
import RationalEnumeration
import ShefferStroke
import NaturalDeduction
import FiniteMatrixNoncharacterizability
import MonotonicityOfEntailment
import PrincipleOfExplosion
import EquationalLogic
import BooleanAlgebra
import PowerTowers
import A158415
import A290268
import SquaredSquare
import KlarnerConstant
import FirstOrder
import ArbitraryLanguageCompactness
import QuantifierCommutation
import PAHF
import NoFiniteModel
import PAFiniteBasisReduction
import PresburgerArithmetic
import RiemannHypothesis
import ZF
import ClosureAxiomatization
import CombinatoryLogic
import BusyBeaver
import TuringDegrees
\end{lstlisting}

The comment in \RepoFile{ProveIt.lean} is a design contract: the facade is
broad, but the expensive BB2/BB3 certificate modules are not imported by the
\project{BusyBeaver} facade.  Similar choices elsewhere prevent a convenient
top-level import from becoming a request to elaborate every generated
certificate in the repository.

\begin{overviewbox}
A facade answers ``which modules should a consumer see by default?''  It does
not answer ``which modules exist?'', ``which modules are default build
targets?'', or ``which theorem is unconditional?''  Those questions belong to
the Lake registry, audits, and project status documents respectively.
\end{overviewbox}

\section{The root Lake workspace}

The root \RepoFile{lakefile.toml} maps logical library names to physical source
directories.  This lets a project keep a mathematically natural location such
as

\begin{center}
\code{Logic/PeanoArithmetic/BoundedConsistency/Lean/}
\end{center}

while exposing a short import root such as \code{BoundedPAConsistency}.  The
workspace also maps the vendored Foundation tree as a Lean library, so
repository-authored developments can import the pinned source without copying
it.

\subsection{Default targets versus registered libraries}

The default targets include the root facade and several libraries that are
deliberately not imported by it:

\begin{lstlisting}[language={},caption={Conceptual form of the root default-target set.}]
defaultTargets := [
  `ProveIt,
  `PAListCoding,
  `PAListCoding.Audit,
  `PAUndecidable,
  `PAUndecidable.Audit,
  `BoundedPAConsistency,
  `BoundedPAConsistency.Audit
]
\end{lstlisting}

This is a subtle but valuable separation.

\begin{itemize}
  \item \textbf{Facade inclusion} controls downstream import cost and namespace
        exposure.
  \item \textbf{Default-target inclusion} controls what a broad workspace
        build is expected to validate.
  \item \textbf{Library registration} controls what can be requested
        explicitly from the root workspace.
\end{itemize}

Thus \project{PAUndecidable} can be part of the ordinary validation policy
without becoming part of every \code{import ProveIt} environment.

\subsection{Registered Lean libraries}

\begin{longtable}{L{.27\textwidth}L{.60\textwidth}}
\caption{Root Lean libraries and their physical roles.}
\label{tab:lean-libraries}\\
\toprule
Logical library & Physical project or role \\
\midrule
\endfirsthead
\toprule
Logical library & Physical project or role \\
\midrule
\endhead
\code{JacobianConjecture} & Algebraic counterexample project.\\
\code{PolynomialFormulas} & Formula, solvability, and Lazard/Dummit project.\\
\code{ProveIt} & Root facade.\\
\code{TrigonometricIdentities} & Exact trigonometric identities.\\
\code{ExponentialIdentities} & Exact exponential identities and bounds.\\
\code{PowerTowers} & Parenthesized power-tower enumeration.\\
\code{A158415} & Radical-expression enumeration.\\
\code{A290268} & Derivative-expansion sequence development.\\
\code{SquaredSquare} & Perfect squared-square existence and lower bound.\\
\code{KlarnerConstant} & Polyomino growth upper bound.\\
\code{BusyBeaver} & Machine semantics and opt-in classifications.\\
\code{CombinatoryLogic} & Lambda/SK/SKI/Iota simulations and universality.\\
\code{TuringDegrees} & Oracle reducibility and degree structure.\\
\code{FirstOrder} & Dependency-free fixed-language first-order metatheory.\\
\bcode{ArbitraryLanguageCompactness} & Mathlib-backed arbitrary-language compactness facade.\\
\code{QuantifierCommutation} & Constructive commutation and counterexamples.\\
\code{ShefferStroke} & Propositional Sheffer-stroke development.\\
\code{NaturalDeduction} & Propositional natural deduction.\\
\bcode{FiniteMatrixNoncharacterizability} & Propositional noncharacterizability result.\\
\code{MonotonicityOfEntailment} & Entailment monotonicity microproject.\\
\code{PrincipleOfExplosion} & Explosion microproject.\\
\code{EquationalLogic} & Shared equational logic.\\
\code{BooleanAlgebra} & Boolean-algebra development.\\
\code{PAHF} & Peano arithmetic / hereditarily finite set interpretation.\\
\code{NoFiniteModel} & No finite model of PA.\\
\code{BoundedPAConsistency} & Bounded-complexity consistency for PA.\\
\code{PAListCoding} & PA-definable list and number-theory predicates.\\
\code{PAFiniteBasisReduction} & Finite-axiomatizability reduction.\\
\code{Foundation} & Pinned vendored Foundation Lean source.\\
\code{PAUndecidable} & Undecidability of PA.\\
\code{PresburgerArithmetic} & Cooper elimination and decision.\\
\code{DiophantineEquations} & Diophantine-equation results, including FLT exponent four.\\
\code{IntegerSums} & Integer-sum identities.\\
\code{RationalEnumeration} & Calkin--Wilf enumeration.\\
\code{RiemannHypothesis} & Arithmetic statement of RH.\\
\code{ZF} & Repository-local lightweight ZF base.\\
\code{BoundedZFCConsistency} & Numeralwise bounded consistency for ZFC.\\
\code{ZFCinPA} & Uniform PA-provability project, with a conditional endpoint.\\
\code{ClosureAxiomatization} & Alternative closure axiom and equivalence with ZF.\\
\bottomrule
\end{longtable}

Several subprojects also carry their own \code{lakefile.toml}.  The local
project is useful when it has a smaller dependency closure, is intentionally
Mathlib-free, or needs different target discovery.  The root workspace is
useful for integration and shared cached dependencies.  When documentation
offers both commands, that is generally an intentional choice, not duplicate
boilerplate.

\section{Pinned Lean dependency graph}

\RepoFile{lean-toolchain} selects
\code{leanprover/lean4:v4.32.0}.  The root manifest pins Mathlib and its
transitive Lake dependencies.  At the inspected snapshot those include
\project{Plausible}, \project{LeanSearchClient}, \project{importGraph},
\project{ProofWidgets}, \project{Aesop}, \project{Qq},
\project{Batteries}, and \project{Cli}.  The manifest is the reproducibility
boundary for the Lean package graph; a source module's imports are only the
human-readable view of that graph.

The ordinary cold-checkout bootstrap is:

\begin{lstlisting}[language=PowerShell,caption={Root Lean bootstrap and broad build.}]
lake exe cache get
lake build
\end{lstlisting}

The broad build is intentionally expensive.  Many project READMEs recommend
setting \code{LEAN_NUM_THREADS=0} to disable Lake's worker pool when individual
modules consume several gigabytes, while some projects that rely on a worker
require the ambient value to be at least one.  Treat the project-local command
as authoritative.

\section{Rocq as a federation of logical roots}

The root \RepoFile{_CoqProject} begins with logical-path declarations, for
example:

\begin{lstlisting}[language={},caption={Representative Rocq logical roots.}]
-Q Algebra/JacobianConjecture/Coq JacobianConjecture
-Q Algebra/PolynomialFormulas/Coq PolynomialFormulas
-R lib/MathComp-Abel/theories Abel
-Q Computability/CombinatoryLogic/Coq CombinatoryLogic
-Q Logic/FirstOrder/Coq FirstOrder
-Q Logic/Modal/Coq FoundationModal
-Q NumberTheory/RationalEnumeration/Coq RationalEnumeration
-Q SetTheory/ZF/Coq ZF
-Q lib/Coq-Synthetic-Computability/theories SyntheticComputability
-Q Computability/TuringDegrees/Coq TuringDegrees
\end{lstlisting}

It then enumerates the source files to be handed to the generated dependency
graph.  The \project{PolynomialFormulas} portion alone includes a long sequence
of core definitions, invariant-theory bridges, generated data shards, and
certificate modules.  The order is not meant to be read as a hand-maintained
linear execution script; \code{rocq makefile} resolves declared imports, while
the registry states the intended source universe.

A root build follows the repository's pinned vendored dependencies:

\begin{lstlisting}[language=PowerShell,caption={Repository-wide Rocq setup and build.}]
git submodule update --init `
  lib/Coq-Synthetic-Computability `
  lib/MathComp-Abel

opam install --yes --deps-only `
  ./lib/MathComp-Abel/coq-mathcomp-abel.opam

pwsh -NoProfile -File `
  Computability/TuringDegrees/Coq/BuildSyntheticComputability.ps1

rocq makefile -f _CoqProject -o Makefile.coq
make -f Makefile.coq
\end{lstlisting}

Focused, sequential \code{coqc} commands are common in project READMEs.  They
serve three purposes: smaller feedback loops, explicit import order for
dependency-light developments, and avoidance of platform-specific parallel
build problems.  Many projects finish with \code{coqchk} or \code{rocqchk} on
the public modules.

\section{Submodules and vendored trees}

The root \RepoFile{.gitmodules} records four submodules:

\begin{enumerate}
  \item the current repository fork/snapshot of the Coq Library of
        Undecidability;
  \item \project{FormalizedFormalLogic/Foundation}, used as a pinned Lean
        dependency and a source reference for independent ports;
  \item \project{MathComp-Abel}, used by the Rocq Abel--Ruffini development;
  \item \project{coq-synthetic-computability}, used by Turing-degree wrappers.
\end{enumerate}

In addition, selected Busy Beaver Rocq certificates are vendored as ordinary
repository content.  \RepoFile{lib/README.md} is the provenance index and states
a strong locality rule: \RepoDir{lib} is the only home for vendored code;
repository-authored wrappers, patches, bridges, and models live in the subject
tree that owns them.

\begin{figure}[htbp]
\centering
\begin{tikzpicture}[
  box/.style={draw=blue!45,fill=sky,rounded corners=1mm,align=center,
              minimum width=39mm,minimum height=12mm,font=\sffamily\small},
  ext/.style={draw=gold!50,fill=sand,rounded corners=1mm,align=center,
              minimum width=39mm,minimum height=12mm,font=\sffamily\small},
  arrow/.style={-{Latex[length=2mm]},thick,draw=muted},
  node distance=7mm and 13mm
]
\node[box] (root) {root Lake workspace};
\node[box,below left=of root] (facade) {\code{ProveIt.lean}\\consumer facade};
\node[box,below=of root] (libs) {subject \code{lean_lib}s\\and local Lake projects};
\node[ext,below right=of root] (mathlib) {pinned Mathlib\\package graph};
\node[box,below=17mm of libs] (rocq) {\code{_CoqProject}\\Rocq federation};
\node[ext,below left=of rocq] (abel) {MathComp-Abel\\and Undecidability};
\node[ext,below right=of rocq] (synth) {Synthetic Computability\\and BB certificates};
\draw[arrow] (root) -- (facade);
\draw[arrow] (root) -- (libs);
\draw[arrow] (mathlib) -- (root);
\draw[arrow] (rocq) -- (abel);
\draw[arrow] (rocq) -- (synth);
\draw[arrow,dashed] (libs) -- node[right,font=\scriptsize\sffamily]{parallel mathematics} (rocq);
\end{tikzpicture}
\caption{Build topology.  Lean and Rocq developments share a repository and often a theorem surface, but not proof artifacts or one common package graph.}
\label{fig:build-topology}
\end{figure}

\section{Why no single build command tells the whole story}

A successful \code{lake build} establishes that the selected Lean targets
elaborate under the pinned package graph.  It does not run Rocq.  A successful
root Rocq make does not validate support scripts, Wolfram replays, article
PDFs, or every local executable generator unless they are registered sources.
A focused project command may intentionally validate more audit modules than a
facade import, or less generated data than a repository-wide build.

The repository handles this by putting validation instructions next to the
project's status claims.  For a serious review, the README command, the facade,
the audit, and the root registry should be read as a four-part specification.

\begin{sourcebox}
\RepoFile{ProveIt.lean};
\RepoFile{lakefile.toml};
\RepoFile{lake-manifest.json};
\RepoFile{lean-toolchain};
\RepoFile{_CoqProject};
\RepoFile{.gitmodules};
\RepoFile{lib/README.md}.
\end{sourcebox}

\chapter{Trust, computation, and theorem status}
\label{chap:trust}

\section{The repository's central engineering question}

Formal mathematics repositories often say that a result is
``machine-checked.''  ProveIt usually goes further and records
\emph{which machine-checked object closes the theorem} and \emph{what remains
outside that object}.  This matters because the repository contains all of
the following:

\begin{itemize}
  \item hand-written proof terms;
  \item elaborator tactics whose output is rechecked by the kernel;
  \item Boolean decision procedures reduced by the kernel;
  \item virtual-machine and native evaluators;
  \item external Python or Wolfram programs that generate exact candidates;
  \item manually transcribed finite data from papers;
  \item imported theorems from Mathlib, MathComp-Abel, Foundation, or other
        pinned libraries;
  \item conditional endpoints with explicit mathematical hypotheses.
\end{itemize}

These mechanisms are not interchangeable.  A useful code walkthrough follows
the evidence path from a public theorem down to the smallest trusted checker
or imported theorem, rather than stopping at the tactic used in the last line.

\section{A layered evidence model}

\begin{figure}[htbp]
\centering
\begin{tikzpicture}[
  layer/.style={draw=blue!45,rounded corners=1.2mm,minimum width=132mm,
                minimum height=12mm,align=left,text width=124mm,
                font=\small},
  trusted/.style={layer,fill=mint,draw=teal!55},
  checked/.style={layer,fill=sky},
  support/.style={layer,fill=sand,draw=gold!55},
  status/.style={layer,fill=rose,draw=brick!55},
  node distance=3.5mm,
  arrow/.style={-{Latex[length=2mm]},thick,draw=muted}
]
\node[trusted] (kernel) {\textbf{Kernel acceptance:} Lean or Rocq checks a theorem term against the type theory and imported constants.};
\node[checked,below=of kernel] (checker) {\textbf{Verified checker or semantics:} a recursive evaluator, substitution function, proof calculus, model semantics, or certificate predicate whose correctness is itself proved.};
\node[checked,below=of checker] (data) {\textbf{Concrete proof data:} explicit polynomial normal forms, machine classifications, placements, traces, finite maps, or derivation codes.};
\node[support,below=of data] (gen) {\textbf{Candidate production and transcription:} Python, Wolfram, search programs, paper diagrams, or upstream tables.  These need not be trusted when the checker is complete for the claim it consumes.};
\node[status,below=of gen] (hyp) {\textbf{Retained hypothesis or open claim:} an explicit premise, named search claim, source-transcription boundary, or documented parity gap.};
\draw[arrow] (kernel) -- (checker);
\draw[arrow] (checker) -- (data);
\draw[arrow] (data) -- (gen);
\draw[arrow,dashed] (hyp) -- node[right,font=\scriptsize\sffamily]{must remain visible} (kernel);
\end{tikzpicture}
\caption{Evidence layers.  External generation can sit below a verified checker; an unresolved mathematical premise must instead remain in the theorem type or status document.}
\label{fig:evidence-layers}
\end{figure}

\subsection{Kernel-checked proof terms}

The ordinary case is a theorem whose elaborated term is accepted by the Lean
or Rocq kernel.  Tactics such as \code{ring}, \code{norm_num}, and
\code{linarith} are proof producers: their implementation can be complicated,
but a bug normally causes them to produce no term or an invalid term, not a
silently accepted theorem.

Imported constants are part of this boundary.  A Lean wrapper around a Mathlib
theorem is fully kernel-checked, but its proof provenance is ``specialization
of imported theorem,'' not ``independent local proof.''  The repository usually
states that distinction explicitly, as in arbitrary-language compactness and
the Lean proof of Fermat's Last Theorem for exponent four.

\subsection{Kernel reduction with \code{decide}}

When a proposition has a decidable instance, a theorem such as

\begin{lstlisting}[language=Lean]
theorem finite_certificate_is_valid :
    checkCertificate certificate = true := by
  decide
\end{lstlisting}

can be reduced inside Lean's kernel.  Some modules use \code{decide +kernel}
to make this intent especially explicit.  This is appropriate for finite
integer grids, sparse-polynomial list equality, and bounded machine searches
when the reduction remains feasible.

The cost is elaboration time, recursion depth, and memory.  Large certificate
modules often raise \code{maxRecDepth} and \code{maxHeartbeats}, or split one
calculation into term, row, tail, and merge certificates so each kernel
reduction stays bounded.

\subsection{VM and native computation}

Rocq projects often use \code{vm_compute}.  Lean can use
\code{native_decide}.  Both accelerate executable reflection by moving
computation away from the slowest kernel reduction path, but they change the
operational trust story.  ProveIt project READMEs commonly state which
evaluator is used.

This is not a binary ``formal/not formal'' distinction.  The theorem is still
accepted by the proof assistant, and the checker can still be proved correct.
The distinction concerns the evaluator and compiler/runtime relied upon to
produce the proof witness.  For example:

\begin{itemize}
  \item \project{SquaredSquare} advertises Lean \code{decide +kernel} and Rocq
        \code{vm_compute};
  \item several projects explicitly say that no \code{native_decide} boundary
        is present;
  \item the root README separates kernel proofs from
        \code{native_decide}/VM computation rather than hiding the difference
        behind a generic ``computed'' label.
\end{itemize}

\subsection{External generators}

A generator can enumerate a finite search space, derive a Groebner-style
normal form, split a giant expression, or emit source modules.  The generator
is outside the proof boundary when all of the following hold:

\begin{enumerate}
  \item its output is represented as ordinary formal data;
  \item a formal checker recomputes the relevant relation;
  \item a theorem proves that checker acceptance implies the mathematical
        statement;
  \item the generated module is actually imported into the endpoint.
\end{enumerate}

A Python script that merely prints ``verified'' does not satisfy this pattern.
A script that emits an exact sparse polynomial subsequently normalized by
formally defined Lean operations does.

\begin{boundarybox}
The right review question is: \emph{Could a malicious or buggy generator make
the proof assistant accept a false theorem?}  If the generator can only choose
candidate data and the formal checker is sound for arbitrary data, the answer
is no.  If the script computes a fact that is merely asserted as an axiom or
unverified constant, the script belongs to the trusted base.
\end{boundarybox}

\section{Audit modules as executable documentation}

Lean audit modules typically contain two kinds of command:

\begin{lstlisting}[language=Lean,caption={Typical Lean audit pattern.}]
#check headline_theorem
#print axioms headline_theorem
\end{lstlisting}

The first catches accidental API drift and documents the intended declaration.
The second prints the constants that Lean treats axiomatically in the theorem's
transitive proof term.  For Mathlib-backed classical mathematics, the common
footprint is \code{propext}, \code{Classical.choice}, and \code{Quot.sound}.
A smaller constructive project may print no assumptions.

Rocq audits use \code{Print Assumptions}; project build recipes may additionally
run \code{coqchk} or \code{rocqchk} over the public modules.  An audit is most
valuable when it names the \emph{headline endpoints}, not merely internal
lemmas.  The Jacobian audit, for example, checks determinant, collision,
noninjectivity, no-polynomial-inverse, dimension-specific, stabilization,
scaling, and simplified-counterexample declarations.

Audits also expose a subtle failure mode: a theorem can have no axiom
dependency yet state the wrong proposition.  The \project{ZFCinPA} README
records an instance where an out-of-scope identifier became an accidental
implicit variable.  The project responded by placing
\code{set_option autoImplicit false} in statement-bearing modules and auditing
the literal endpoint.  Assumption audits verify proof dependencies; they do
not replace specification review.

\section{Conditional endpoints and named open claims}

A conditional theorem is not a defect when the condition is the intended
formal boundary and is documented accurately.  ProveIt uses several patterns.

\subsection{Hypothesis in the theorem type}

The reusable \project{KlarnerConstant} theorem in \code{Main.lean} accepts
abstract recurrence and domination hypotheses.  The concrete endpoint in
\code{GeometricComplete.lean} supplies those hypotheses for actual polyomino
counts, so the user-facing bound is unconditional.  Here the conditional
intermediate theorem is an abstraction boundary.

By contrast, the current \project{ZFCinPA} endpoint keeps
\bcode{ZFCSuccessorImplicationsInAllModels} as a premise.  Its README states
that nonstandard-index successor implications remain open.  The reduction is
proved to be an equivalence, so the premise is not disguised as a harmless
implementation detail.

\subsection{Named proposition for an unformalized search}

\project{SquaredSquare} defines \code{DuijvestijnSearchClaim} for the exact
statement that every perfect squared square has at least 21 pieces.  It proves
that sharp minimality is equivalent to this claim and proves the unconditional
bracket \(7 \leq \text{minimum order} \leq 21\).  No axiom is introduced for
the historical exhaustive search.

This naming style is especially honest: users can import and reason about the
reduction without accidentally receiving the open search result as a theorem.

\subsection{Deliberate nonclaims}

The Turing-degree project enumerates deep existence and classification
theorems it has not formalized.  The A198683 project withholds \(n=12\) from
the accepted OEIS data field while a proof-quality certificate is missing.
The Busy Beaver facade excludes expensive classifications, and the topic
README distinguishes a checked score theorem from a conditional root equality.
In each case, omission is an architectural status choice rather than an
absence readers must infer.

\section{Independent proofs, ports, and parity}

The Lean and Rocq directories often prove similarly named theorems, but their
relationship varies:

\begin{longtable}{L{.25\textwidth}L{.64\textwidth}}
\caption{Common cross-prover relationships.}
\label{tab:cross-prover}\\
\toprule
Relationship & Meaning \\
\midrule
\endfirsthead
\toprule
Relationship & Meaning \\
\midrule
\endhead
Independent reproof &
Both sides define comparable semantics and prove a theorem separately, as in
quantifier commutation or first-order completeness.\\
Contract-level parity &
The numeric representation differs, but both sides prove the same public
round-trip/functionality contract, as in PA list coding.\\
Complementary developments &
One side provides results absent on the other because its imported library or
constructive setting differs.  Turing degrees is the clearest example.\\
Imported specialization versus local proof &
One side wraps an upstream theorem while the other reconstructs an argument,
as with exponent-four Fermat.\\
Source port &
A pinned source module is used as the declaration inventory, while the target
proof assistant independently recreates the mathematical surface.
FoundationPort and Modal use this pattern.\\
Parallel endpoint with a parity gap &
Most layers correspond, but one endpoint remains conditional or absent in one
assistant.  Bounded PA consistency documents such a gap for the strongest
uniform sentence.\\
\bottomrule
\end{longtable}

The phrase ``independent'' should therefore be read locally.  A Rocq port can
be independent of a Lean implementation while still importing a pinned Rocq
library.  A Lean theorem can be independent of a Coq proof while relying on
Mathlib.  The relevant provenance appears in the project README and import
graph.

\section{Source transcription as a separate boundary}

Some finite formal objects originate in diagrams or tables that a kernel
cannot visually interpret.  The \project{KlarnerConstant} patterns were
manually compared against Bui's neighborhood diagrams.  The squared-square
placement comes from a Bouwkamp code.  Once coordinates are entered, all local
partitions, injections, coverage checks, and arithmetic can be formal, but
``these coordinates faithfully transcribe that page'' remains a human source
boundary unless a parser for the source artifact is itself verified and used.

The repository's strongest documentation names this boundary rather than
confusing it with mathematical proof.  A reviewer can then distinguish:

\begin{enumerate}
  \item correctness of the formal theorem relative to entered data;
  \item correctness of the data transcription;
  \item novelty or historical attribution;
  \item reproducibility of the candidate-generation process.
\end{enumerate}

Only the first is necessarily a kernel theorem.  All four can be important.

\section{A review checklist}

\begin{longtable}{L{.34\textwidth}L{.55\textwidth}}
\caption{Trust-oriented questions for a headline theorem.}
\label{tab:trust-checklist}\\
\toprule
Question & Where to answer it \\
\midrule
\endfirsthead
\toprule
Question & Where to answer it \\
\midrule
\endhead
What is the literal theorem type? &
The endpoint module and its audit, not the README paraphrase alone.\\
Is the result imported, re-proved, or computed? &
Imports, proof body, and project provenance notes.\\
Which evaluator closes finite computation? &
The proof term (\code{decide}, \code{native_decide}, \code{vm_compute}) and
trust-boundary section.\\
Can generated data be adversarial? &
The checker correctness theorem and its quantification over candidate data.\\
Does a premise survive? &
The theorem type, \code{\#check}, and status table.\\
Are all intended modules built? &
Facade imports, Lake globs/targets, or Rocq source registry.\\
Does an audit cover the endpoint? &
\code{Audit.lean}/\code{Audit.v}; check the exact declaration name.\\
Do Lean and Rocq prove the same statement? &
Coverage ledger, encoding contract, and hypothesis lists.\\
Is the source transcription formalized? &
Project README's trust-boundary or provenance section.\\
Is an omitted theorem intentionally omitted? &
``Deliberate boundary,'' ``nonclaims,'' or ``residue'' documentation.\\
\bottomrule
\end{longtable}

\begin{sourcebox}
\RepoFile{README.md};
\RepoFile{Algebra/JacobianConjecture/Lean/JacobianConjecture/Audit.lean};
\RepoFile{Combinatorics/SquaredSquare/README.md};
\RepoFile{Combinatorics/Polyominoes/KlarnerConstant/README.md};
\RepoFile{Computability/TuringDegrees/COVERAGE.md};
\RepoFile{Logic/Interpretability/ZFCinPA/README.md}.
\end{sourcebox}

\chapter{Recurring module architectures}
\label{chap:patterns}

\section{Facade, implementation, endpoint, audit}

The most common project shape is:

\[
\text{facade} \longrightarrow
\text{definitions and lemmas} \longrightarrow
\text{headline endpoint} \longrightarrow
\text{audit}.
\]

The arrow here means ``imports or exposes,'' not necessarily one file per
stage.  A facade normally has little or no theorem body.  The implementation
modules are organized by mathematical responsibility.  An endpoint module
assembles reusable lemmas into the theorem promised by the README.  The audit
imports the public surface and asks the prover to display types and
assumptions.

\begin{figure}[htbp]
\centering
\begin{tikzpicture}[
  n/.style={draw=blue!45,fill=sky,rounded corners=1.1mm,
            minimum width=31mm,minimum height=11mm,align=center,
            font=\sffamily\small},
  a/.style={draw=teal!50,fill=mint,rounded corners=1.1mm,
            minimum width=31mm,minimum height=11mm,align=center,
            font=\sffamily\small},
  arrow/.style={-{Latex[length=2mm]},thick,draw=muted},
  node distance=8mm
]
\node[n] (facade) {facade};
\node[n,right=of facade] (core) {semantic core};
\node[n,right=of core] (bridge) {bridges/lemmas};
\node[n,right=of bridge] (endpoint) {endpoint};
\node[a,below=of endpoint] (audit) {audit};
\draw[arrow] (facade) -- (core);
\draw[arrow] (core) -- (bridge);
\draw[arrow] (bridge) -- (endpoint);
\draw[arrow] (endpoint) -- (audit);
\draw[arrow,dashed] (facade) to[bend right=23] (endpoint);
\end{tikzpicture}
\caption{The ordinary hand-written project shape.  A facade may import every stage directly, but the conceptual dependency direction remains useful.}
\label{fig:ordinary-shape}
\end{figure}

Examples:

\begin{itemize}
  \item \project{FirstOrder}: \code{Fol} $\to$ \code{Calculus}
        $\to$ \code{Completeness} $\to$
        \code{ClassicalCompleteness}/\code{Compactness} $\to$ \code{Audit}.
  \item \project{PresburgerArithmetic}: \code{Syntax} $\to$
        \code{NormalForm} $\to$ \code{Cooper} $\to$
        \code{Elimination} $\to$ \code{Decision} $\to$ \code{Audit}.
  \item \project{JacobianConjecture}: shared definitions and the main map,
        collision/scaling/equivariance/simplification/stabilization modules,
        then a broad audit.
\end{itemize}

This pattern gives a reliable debugging rule: an error in a facade import is
usually downstream of the true source.  Build or elaborate the lowest module
that still reproduces the problem.

\section{The explicit semantic ladder}

Metatheory projects make levels that ordinary application code often merges:

\begin{enumerate}
  \item object-language syntax;
  \item variable operations such as shift, substitution, or renaming;
  \item interpretation in a structure under an assignment;
  \item derivation objects or a provability relation;
  \item soundness;
  \item model construction or decision;
  \item completeness/compactness/consistency endpoint.
\end{enumerate}

The separation is not ceremonial.  A theorem about syntactic consequence can
be wrong even if the semantic theorem it resembles is true.  A coded proof
predicate can inspect only conclusions even though the intended statement
bounds every formula occurrence.  A substitution theorem can hold on standard
quoted terms while failing for nonstandard codes in a nonstandard model.

The code therefore contains many ``bridge'' modules whose only purpose is to
connect adjacent levels: quotation adequacy, decoder soundness, erasure from a
data-carrying proof tree to a Prop-valued derivation, translation identity,
truth lemmas, or model-relative satisfaction specifications.

\section{Data/checker/meaning for finite certificates}

A robust finite certificate has three logically distinct components.

\begin{description}
  \item[Data] The concrete list, table, trace, placement, machine, sparse
        polynomial, or decomposition.
  \item[Checker] An executable predicate or function that computes whether the
        data satisfy a finite condition.
  \item[Meaning theorem] A proof that acceptance implies the mathematical
        property stated by the endpoint.
\end{description}

For a squared-square dissection, data are integer rectangles; the checker
establishes bounds, separation, distinct sides, and exact unit-cell coverage;
analytic lemmas lift that certificate to closed squares in the real plane.
For Busy Beaver, data are machine traces or classification tables; semantics
connect trace replay to halting score.  For Dummit coefficients, data are
sparse polynomial normal forms; substitution and normalization semantics
connect list equality to the polynomial identity.

\begin{boundarybox}
A checker that validates only a \emph{necessary} condition cannot justify an
endpoint that needs a sufficient condition.  Completeness of the checker for
the consumed claim is a theorem obligation, not a property conferred by
calling the data a certificate.
\end{boundarybox}

\section{Generated sharding}

Large exact expressions are split for both computational and software reasons.
Common names reveal the split:

\begin{itemize}
  \item \code{Data}: literal generated definitions;
  \item \code{Certificate}: equality or acceptance proofs;
  \item \code{TermNN}: one source term's contribution;
  \item \code{TailNN}: an inductive suffix of a long sum;
  \item \code{RowData}: stable shared table input;
  \item \code{Support}: helper definitions or finite support;
  \item \code{Shard}, \code{Part}, or numbered coefficient modules:
        bounded pieces of one global identity;
  \item \code{Core}, \code{Common}, \code{Bridge}: hand-written semantic
        infrastructure intended to remain small.
\end{itemize}

A typical tail certificate has the following shape:

\begin{lstlisting}[language=Lean,caption={Representative generated-certificate composition.}]
import PolynomialFormulas....Tail22Data

set_option maxRecDepth 1000000
set_option maxHeartbeats 20000000

theorem table0_tail22_merge_certificate :
    table0Tail22Candidate = table0Tail22Normal := by
  decide

theorem table0_tail22_certificate :
    SparsePolynomial.substitute table0Tail22 elementaryPolynomials =
      table0Tail22Normal := by
  change add (substitute term88 ...) (
    add (substitute term89 ...) (
    add (substitute term90 ...) (
    add (substitute term91 ...)
        (substitute table0Tail23 ...)))) = table0Tail22Normal
  rw [table0_term88_certificate,
      table0_term89_certificate,
      table0_term90_certificate,
      table0_term91_certificate,
      table0_tail23_certificate]
  exact table0_tail22_merge_certificate
\end{lstlisting}

The long \code{change} statement exposes definitional structure so rewriting
can compose small certificates.  The final \code{decide} compares concrete
normalized lists.  A reviewer should read the sparse data type and substitution
semantics first, then one representative term certificate and one tail
certificate, and only then inspect generator code.

\section{Compatibility facades}

Refactoring large developments often leaves small modules that preserve an old
import surface.  The \project{KlarnerConstant} README names
\code{GeometricPPartition.lean}, \code{GeometricQ.lean}, and
\code{GeometricVW.lean} as compatibility facades over newer multi-stage
chains.  These files are not dead weight: they decouple downstream imports from
an internal split and make git history easier to follow.

Likewise, the root \code{ProveIt.lean} is a compatibility and convenience
surface over a physical taxonomy.  Removing facades merely because they have
no declarations would make the import graph more brittle.

\section{Parallel but noncomposable semantics}

Two formalizations can both be correct and still be unsafe to compose without
a bridge.  \project{CombinatoryLogic} calls this out explicitly: its local
lambda calculus and the lambda semantics used by the vendored Undecidability
library have different operational presentations.  Theorems are not silently
chained across them.

The same warning applies to:

\begin{itemize}
  \item quotient versus setoid presentations of Turing degrees;
  \item different natural-number encodings of lists;
  \item the repository-local fixed first-order language versus Mathlib's
        arbitrary signatures;
  \item Foundation's set-theory language versus the repository-local
        \code{Form} used by the lightweight ZF project;
  \item external standard numerals versus internal elements of a nonstandard
        model.
\end{itemize}

A bridge needs a translation and a theorem preserving the relevant semantics
or derivability.  Similar notation and matching informal intent are not enough.

\section{Status documentation as code architecture}

Some README sections function almost like interfaces:

\begin{description}
  \item[Status table] Lists exact endpoints and whether each is proved,
        conditional, open, or proved only at standard indices.
  \item[Coverage ledger] Maps source claims to Lean/Rocq declaration names and
        states missing claims.
  \item[Residue] Orders remaining obligations by dependency.
  \item[Trust boundary] Separates kernel facts, computation modes, generated
        candidates, and transcription.
  \item[Deliberate boundary/nonclaims] Prevents a partial development from
        being mistaken for a theorem inventory.
\end{description}

For fast-moving formalization, this documentation is not secondary.  It is the
human-readable counterpart of theorem types and build targets.  The strongest
examples are \RepoFile{Computability/TuringDegrees/COVERAGE.md},
\RepoFile{Logic/Interpretability/ZFCinPA/README.md}, and the bounded-consistency
READMEs.

\section{Resource controls are part of module design}

Options such as \code{maxHeartbeats} and \code{maxRecDepth} are common in
generated certificate modules.  Their presence says something architectural:

\begin{itemize}
  \item the theorem is expected to close by reduction or large rewriting;
  \item the shard size was chosen against elaborator limits;
  \item moving a definition from reducible data to an opaque theorem can alter
        performance drastically;
  \item a ``cleanup'' that merges modules may make the proof infeasible even
        if the mathematical term is equivalent.
\end{itemize}

Worker settings in build instructions similarly reflect memory behavior.  One
large module per process can be safer than many concurrent modest modules.
Generated certificate topology is therefore partly a scheduling design.

\section{Naming as a navigation API}

Names generally tell a reviewer where to look:

\begin{longtable}{L{.25\textwidth}L{.64\textwidth}}
\caption{Frequent naming signals.}
\label{tab:naming-signals}\\
\toprule
Name fragment & Likely role \\
\midrule
\endfirsthead
\toprule
Name fragment & Likely role \\
\midrule
\endhead
\code{Basic}, \code{Core}, \code{Defs} & Fundamental datatypes and semantics.\\
\code{Common}, \code{Support} & Shared technical lemmas, often not a public endpoint.\\
\code{Syntax}, \code{Coding} & Object syntax or numeric encodings.\\
\code{Evaluation}, \code{Semantics}, \code{Soundness} & Meaning bridge.\\
\code{Compiler}, \code{Translation}, \code{Adapter}, \code{Bridge} & Preservation across representations.\\
\code{Completeness}, \code{Endpoint}, \code{Main} & Assembly toward a headline theorem.\\
\code{Audit} & Public declarations and assumption reporting.\\
\code{Data}, \code{Table}, \code{Matrix} & Literal or generated finite objects.\\
\code{Certificate}, \code{Check} & Acceptance proof for data.\\
\code{Research}, \code{Article} & Exposition/provenance; not automatically built theorem code.\\
\code{Generate}, \code{verify_} & External or executable support tooling.\\
\bottomrule
\end{longtable}

These conventions are strong enough that a reader can often reconstruct the
intended dataflow from a directory listing before opening a proof.

\begin{sourcebox}
\RepoFile{Logic/FirstOrder/README.md};
\RepoFile{Logic/PresburgerArithmetic/README.md};
\RepoFile{Algebra/PolynomialFormulas/Lean/PolynomialFormulas/ComputableDummitKernelCertificateTable0Tail22Certificate.lean};
\RepoFile{Algebra/PolynomialFormulas/Lean/PolynomialFormulas/ComputableDummitKernelCertificateTable0Tail22Data.lean};
\RepoFile{Combinatorics/Polyominoes/KlarnerConstant/README.md};
\RepoFile{Computability/CombinatoryLogic/README.md}.
\end{sourcebox}

\part{Subject-by-subject walkthrough}

\chapter{Algebra}
\label{chap:algebra}

\section{Domain overview}

\RepoDir{Algebra} contains only two project directories, but they occupy very
different architectural scales.

\begin{description}
  \item[\project{JacobianConjecture}] A compact explicit polynomial-map
        construction.  Its main facts are determinant and collision
        identities, followed by invariant-preserving transformations such as
        scaling, stabilization, and degree-lowering shears.
  \item[\project{PolynomialFormulas}] A broad theory and computation project:
        low-degree formulas, algebraic radical syntax, numerical
        approximation, Abel--Ruffini, per-polynomial solvability decisions,
        and the very large Lazard/Dummit invariant-theory certificate program.
\end{description}

The first project demonstrates the ordinary
definitions--calculation--endpoint--audit pattern.  The second demonstrates
how that pattern scales when one theorem consumes hundreds or thousands of
exact generated sub-certificates.

\section{JacobianConjecture}
\label{sec:jacobian}

\subsection{Public claim and mathematical object}

The central module \RepoFile{Algebra/JacobianConjecture/Lean/JacobianConjecture/Counterexample.lean}
defines an \(n\)-dimensional polynomial map over a commutative semiring as

\[
  \operatorname{PolynomialMap}(n,R)
  \;=\; \operatorname{Fin}(n)\to R[X_0,\ldots,X_{n-1}],
\]

then defines coordinatewise point evaluation, the formal Jacobian matrix
\(\partial f_i/\partial X_j\), its determinant, the Keller condition
``nonzero constant Jacobian,'' and polynomial invertibility via the induced
maps on points.

The displayed three-variable map is

\[
\begin{aligned}
P(x,y,z) &= (1+xy)^3z+y^2(1+xy)(4+3xy),\\
Q(x,y,z) &= y+3x(1+xy)^2z+3xy^2(4+3xy),\\
R(x,y,z) &= 2x-3x^2y-x^3z.
\end{aligned}
\]

The formal determinant is proved equal to \(-2\).  The module also proves that
the distinct integral points
\[
  (-1,1,5),\qquad (0,-2,-16)
\]
both evaluate to \((0,-2,0)\).  It therefore derives noninjectivity, absence
of a pointwise left inverse, absence of a polynomial two-sided inverse, and
the negation of the repository's fixed-dimensional and all-dimensional
Jacobian-conjecture predicates over characteristic-zero fields.

This was an exceptionally recent mathematical development at the snapshot:
the explicit map was publicly announced in July 2026.  The code walkthrough
does not rely on community-status language; the local formal evidence is
particularly transparent because the two decisive claims are finite
polynomial identities.

\subsection{Core definitions}

The public API in \code{Counterexample.lean} is intentionally small:

\begin{lstlisting}[language=Lean,caption={Normalized outline of the Jacobian core definitions.}]
abbrev PolynomialMap (n : Nat) (R : Type u) [CommSemiring R] :=
  Fin n -> MvPolynomial (Fin n) R

noncomputable def PolynomialMap.evalMap
    (f : PolynomialMap n R) (a : Fin n -> R) : Fin n -> R :=
  fun i => MvPolynomial.eval a (f i)

noncomputable def PolynomialMap.jacobian
    (f : PolynomialMap n R) :
    Matrix (Fin n) (Fin n) (MvPolynomial (Fin n) R) :=
  fun i j => pderiv j (f i)

noncomputable def PolynomialMap.jacobianDet
    (f : PolynomialMap n R) : MvPolynomial (Fin n) R :=
  (jacobian f).det

def PolynomialMap.HasNonzeroConstantJacobian (f : PolynomialMap n R) : Prop :=
  exists c : R, c != 0 /\ jacobianDet f = C c

def PolynomialMap.HasPolynomialInverse (f : PolynomialMap n R) : Prop :=
  exists g, LeftInverse (evalMap g) (evalMap f) /\
            RightInverse (evalMap g) (evalMap f)
\end{lstlisting}

The inverse predicate is stated extensionally on points.  This is sufficient
for the refutation: any two-sided polynomial inverse supplies a left inverse,
and a left inverse forces injectivity.  The module never needs to formalize
polynomial composition equality in order to conclude that the displayed map
has no polynomial inverse.

\subsection{Determinant path}

Rather than asking the general matrix determinant simplifier to normalize the
entire expression opaquely, the file defines the closed \(3\times3\)
determinant formula \code{det3} and proves it equals Mathlib's
\code{Matrix.det}.  It then:

\begin{enumerate}
  \item unfolds the three coordinates and formal partial derivatives;
  \item supplies a helper proving that numeric literals have zero
        \code{pderiv};
  \item normalizes the resulting polynomial identity with \code{ring};
  \item transports the result through \code{det3_eq_det};
  \item rewrites the constant polynomial into \code{C (-2)}.
\end{enumerate}

The expensive local lemma receives scoped
\code{maxHeartbeats 1600000} and \code{maxRecDepth 8000}.  This keeps the
resource concession adjacent to the calculation that needs it rather than
globally weakening safeguards.

The reusable theorem \code{det_eq_det_toBlocks11} handles a block
lower-triangular matrix with zero upper-right block and identity lower-right
block.  Later five-variable and stabilized determinant arguments reduce to
this lemma.  This is a good example of extracting the mathematical structure
from what could otherwise become repeated large \code{ring} calls.

\subsection{Collision path}

The collision proof is intentionally elementary:

\begin{enumerate}
  \item define the two points and common value as \code{Fin 3 -> R} vectors;
  \item prove each coordinate equality by function extensionality;
  \item split the finite coordinate index with \code{fin_cases};
  \item evaluate with \code{norm_num};
  \item prove the points distinct by comparing coordinate zero;
  \item contradict injectivity.
\end{enumerate}

The determinant theorem is polymorphic over every commutative ring.  The
Keller condition adds \code{Field} and \code{CharZero} only where \(-2\neq0\)
is needed.  The collision theorem needs merely a nontrivial commutative ring.
This stratification makes the algebraic content sharper than a proof written
only over \(\mathbb C\).

\subsection{Supporting modules}

\begin{longtable}{L{.29\textwidth}L{.60\textwidth}}
\caption{JacobianConjecture Lean module responsibilities.}
\label{tab:jacobian-modules}\\
\toprule
Module & Role \\
\midrule
\endfirsthead
\toprule
Module & Role \\
\midrule
\endhead
\code{Counterexample} &
Definitions, determinant, integral collision, noninjectivity, and the
dimension-three/all-dimensions endpoints.\\
\code{CollisionFamily} &
A parameterized family of collisions and useful specializations.\\
\code{Equivariance} &
Symmetries of the map and mirrored/rational collision consequences.\\
\code{Scaling} &
Source/target scalings, composition laws, and scaled collision families.\\
\code{SimplerCounterexample} &
A stably equivalent degree-lowered presentation, including total-degree
bounds and its own determinant/collision endpoint.\\
\code{Stabilization} &
Adjoins identity coordinates, transports evaluation and determinant, and
extends the counterexample to every dimension at least three.\\
\code{Audit} &
Checks and prints axioms for the public declarations across all of the above.\\
\bottomrule
\end{longtable}

The audit is unusually comprehensive.  It does not stop at one final
negation; it enumerates determinant identities, all collision witnesses,
equivariance, scaling laws, degree profiles, simplified maps, and stabilization
bridges.  That makes \RepoFile{Algebra/JacobianConjecture/Lean/JacobianConjecture/Audit.lean}
the fastest public-API index.

The Rocq side uses a small polynomial-expression syntax and interprets it over
Coquelicot's complex numbers.  Its source sequence mirrors the main themes:
\code{Common.v}, \code{Counterexample.v}, \code{Scaling.v},
\code{CollisionFamily.v}, \code{SimplerCounterexample.v}, and
\code{Audit.v}.  It is an independent proof artifact rather than an extraction
of the Lean theorem.

\subsection{Recommended reading order}

\begin{readingbox}
\begin{enumerate}
  \item \RepoFile{Algebra/JacobianConjecture/README.md} for status and build.
  \item The definitions and two decisive lemmas in
        \RepoFile{Algebra/JacobianConjecture/Lean/JacobianConjecture/Counterexample.lean}.
  \item \code{Stabilization.lean} to see how a dimension-specific witness is
        transported.
  \item \code{SimplerCounterexample.lean} for the stable degree-lowering
        construction.
  \item \code{Audit.lean} as the declaration checklist.
  \item The Rocq modules only after the Lean object and theorem surface are
        clear; their representation is intentionally different.
\end{enumerate}
\end{readingbox}

\section{PolynomialFormulas}
\label{sec:polynomial-formulas}

\subsection{Why this project needs a map before a file list}

\project{PolynomialFormulas} spans several mathematically distinct layers:

\begin{enumerate}
  \item exact formulas for degrees one through four;
  \item an algebraic language of radicals and exact root collections;
  \item executable Gaussian-rational approximation;
  \item impossibility theorems such as Abel--Ruffini;
  \item per-polynomial radical-solvability decision procedures;
  \item general resolvent and invariant-theory infrastructure;
  \item explicit quintic constructions in the style of Lazard;
  \item generated sparse-polynomial tables and kernel certificates.
\end{enumerate}

A flat alphabetical listing hides this progression.  The facade
\RepoFile{Algebra/PolynomialFormulas/Lean/PolynomialFormulas.lean} provides a
much better overview.  Its imports can be grouped as follows.

\begin{longtable}{L{.22\textwidth}L{.32\textwidth}L{.35\textwidth}}
\caption{Conceptual strata of the PolynomialFormulas facade.}
\label{tab:polynomial-strata}\\
\toprule
Stratum & Representative imports & Responsibility \\
\midrule
\endfirsthead
\toprule
Stratum & Representative imports & Responsibility \\
\midrule
\endhead
Classical formulas &
\code{Basic}, \code{Cubic}, \code{Quartic} &
Low-degree equations, algebraic branch conditions, and root exhaustiveness.\\
Exact radical execution &
\bcode{GaussianPolynomialSolver}, \bcode{GaussianRadicalBounds},
\bcode{GaussianPolynomialApproximation} &
Finite radical-expression syntax over Gaussian rationals, evaluation
relations, and certified approximations.\\
Quintic constructions &
\code{LazardQuintic}, \code{LazardQuinticFourier},
\bcode{LazardQuinticSectionFiveCombinedAction},
\bcode{GaussianQuinticSolver} &
Resolvents, Fourier/invariant coordinates, conditional solver packaging,
and executable fallback behavior.\\
Impossibility and decision &
\code{AbelRuffini}, \bcode{QuinticRadicalDecision},
\bcode{SexticRadicalDecision}, \bcode{SelmerAbelRuffini},
\bcode{GenericAbelRuffini} &
Generic and explicit nonsolvability plus decidability for concrete low-degree
integer polynomials.\\
Examples &
\code{Examples} &
Small end-to-end specializations demonstrating the public API.\\
\bottomrule
\end{longtable}

The facade does not import every internal Lazard certificate module directly;
higher-level modules aggregate them.  Consequently the import graph, not the
directory's file count, defines the current endpoint.

\subsection{Low-degree formulas: relations before branches}

The degree-two, degree-three, and degree-four developments avoid treating a
computer-language square root as if it were an algebraic branch-independent
operation.  A radical witness is characterized by an equation such as
\(s^2=\Delta\), and the theorem states what follows for any witness satisfying
that relation.  This has several benefits:

\begin{itemize}
  \item the core theorem works over fields/rings with the stated algebraic
        properties rather than only over \(\mathbb C\);
  \item branch conventions do not leak into purely algebraic identities;
  \item degeneracies and characteristic restrictions remain visible;
  \item an executable complex solver can later choose a principal branch and
        prove it instantiates the abstract relation.
\end{itemize}

Root collections are proved sound and exhaustive where the hypotheses permit
it.  This is stronger than showing that substituting one displayed formula
yields zero: it explains how multiplicities, zero discriminants, and branch
families cover every root.

\subsection{Gaussian-rational radical syntax}

For executable exact work, the project defines finite expression trees whose
leaves are Gaussian rationals and whose nodes represent arithmetic and selected
radicals.  A solver can return such a tree without pretending that its
transcendental numerical rendering is exact.  Separate modules prove:

\begin{itemize}
  \item algebraic evaluation or root correctness;
  \item bounds on radical magnitude and propagation of approximation error;
  \item a computable approximation procedure to a requested precision;
  \item compatibility between exact syntax and approximate complex values.
\end{itemize}

This separation is important for the quintic layer.  The theorem that a
returned expression denotes a root is exact; approximation is an interface
for observing or testing the expression, not the proof of roothood.

\subsection{Abel--Ruffini surfaces}

The project contains several related but nonidentical impossibility results.

\begin{description}
  \item[Explicit example] A concrete quintic, such as \(X^5-4X+2\), is shown
        not solvable by radicals through imported and locally bridged Galois
        theory.
  \item[Rootwise form] The statement is expressed in terms of which individual
        roots can lie in radical towers, not merely whether a splitting field
        is solvable.
  \item[Selmer family] \(X^n-X-1\) supplies a family across arbitrary degrees
        under the relevant hypotheses.
  \item[Generic form] The symmetric rational-function extension gives the
        usual generic obstruction and connects it to explicit resolvents.
\end{description}

On the Rocq side, the MathComp-Abel submodule provides the deep Galois-theory
base.  Local modules adapt its theorem surface to the project's polynomial
definitions and examples.  A reviewer should therefore distinguish the local
bridge proof from the upstream Abel theorem itself.

\subsection{Deciding radical solvability of a concrete quintic}

The decision project is not a universal formula for quintic roots.  It takes a
specific integer/rational quintic and decides whether its Galois group is
solvable using finite exact data.  Conceptually, the Lean path is:

\begin{enumerate}
  \item normalize the polynomial and split reducible cases by bounded
        factor search;
  \item handle low-degree factors with known radical formulas;
  \item for the irreducible quintic, compute exact resolvent data;
  \item use Frobenius/Dummit-style criteria and a Chapman correction to
        classify the relevant transitive group cases;
  \item turn the finite decision into a \code{ComputablePred};
  \item connect the algorithm to the repository's Turing-machine interface.
\end{enumerate}

The bounded factor search is primitive recursive.  The resulting quintic
decision is packaged strongly enough to obtain a verified machine-level
computability result.  ``Decidable'' here means an executable theorem about
each concrete polynomial, not that the generic quintic has a radical formula.

\subsection{The sextic decision}

The sextic project generalizes the strategy but changes the recursion
classification.  It uses degree-15 and degree-10 resolvents and a finite
classification of transitive subgroups of \(S_6\).  The repository claims a
recursive/computable decision, not the stronger primitive-recursive packaging
advertised for the quintic.  Lean connects the procedure to a verified Turing
machine; the Rocq version uses its own \(\mu\)-recursive infrastructure.

This distinction is a good example of accurate theorem granularity:
computability, primitive recursiveness, and practical runtime are three
different claims.

\subsection{The Lazard/Dummit semantic core}

Before the generated modules, a small hand-written representation defines
five-variable sparse monomials:

\begin{lstlisting}[language=Lean,caption={Core sparse representation used by Dummit certificates.}]
structure Powers where
  p0 : Nat
  p1 : Nat
  p2 : Nat
  p3 : Nat
  p4 : Nat
deriving Repr, DecidableEq

structure SparseTerm where
  coeff : Int
  powers : Powers
deriving Repr, DecidableEq

abbrev SparsePolynomial := List SparseTerm

def SparseTerm.eval [CommRing R]
    (t : SparseTerm) (x : Fin 5 -> R) : R := ...

def SparsePolynomial.eval [CommRing R] :
    SparsePolynomial -> (Fin 5 -> R) -> R
  | [], _ => 0
  | t :: q, x => t.eval x + SparsePolynomial.eval q x
\end{lstlisting}

The same core supplies addition, multiplication, substitution, normalization,
and evaluation lemmas.  A seven-entry \code{dummitTable} stores coefficient
polynomials in ascending sextic degree.  This representation is deliberately
simple: lists and integer exponent tuples have decidable equality and reduce
predictably.  Sophisticated invariant theory is proved \emph{around} a
computation-friendly certificate language.

\subsection{A concrete generated certificate chain}

At the pinned commit,
\RepoFile{Algebra/PolynomialFormulas/Lean/PolynomialFormulas/ComputableDummitKernelCertificateTable0Tail22Data.lean}
imports shared row data, the next tail certificate, and certificates for terms
88 through 91.  It defines a very large literal normalized sparse polynomial.
The adjacent
\RepoFile{Algebra/PolynomialFormulas/Lean/PolynomialFormulas/ComputableDummitKernelCertificateTable0Tail22Certificate.lean}
then proves two facts:

\begin{enumerate}
  \item the assembled candidate list equals the stored normal list, by
        kernel \code{decide};
  \item formal substitution of the source tail equals that normal list, by
        rewriting the four term certificates and the next-tail certificate.
\end{enumerate}

The full proof path is therefore:

\begin{figure}[htbp]
\centering
\begin{tikzpicture}[
  b/.style={draw=blue!45,fill=sky,rounded corners=1mm,align=center,
            minimum width=31mm,minimum height=12mm,font=\sffamily\small},
  s/.style={draw=gold!50,fill=sand,rounded corners=1mm,align=center,
            minimum width=31mm,minimum height=12mm,font=\sffamily\small},
  t/.style={draw=teal!50,fill=mint,rounded corners=1mm,align=center,
            minimum width=31mm,minimum height=12mm,font=\sffamily\small},
  arrow/.style={-{Latex[length=2mm]},thick,draw=muted},
  node distance=5mm
]
\node[s] (gen) {Python generator\\or algebra search};
\node[b,right=of gen] (lit) {literal \code{Data}\\modules};
\node[b,right=of lit] (term) {term/tail\\certificates};
\node[b,right=of term] (sem) {substitution and\\evaluation lemmas};
\node[t,below=8mm of sem] (end) {Lazard/Dummit\\mathematical endpoint};
\node[t,left=of end] (audit) {facade and\\assumption audit};
\draw[arrow] (gen) -- node[above,font=\scriptsize\sffamily]{emits candidate} (lit);
\draw[arrow] (lit) -- node[above,font=\scriptsize\sffamily]{kernel equality} (term);
\draw[arrow] (term) -- (sem);
\draw[arrow] (sem) -- (end);
\draw[arrow] (end) -- (audit);
\draw[arrow,dashed] (gen) to[bend right=32] node[below,font=\scriptsize\sffamily]{not trusted for truth} (term);
\end{tikzpicture}
\caption{Generated certificate path in PolynomialFormulas.}
\label{fig:dummit-pipeline}
\end{figure}

A bug in the generator should cause a certificate theorem to fail.  A review
of the generator is still valuable for reproducibility and shard completeness,
but the theorem's soundness is meant to rest on the formal operations and
their semantic lemmas.

\subsection{Generator programs and reproducibility}

The support directory contains substantial generators, including:

\begin{itemize}
  \item \code{generate_dummit_kernel_certificates.py};
  \item \code{generate_fin5_transitive_classification_certificates.py};
  \item \code{generate_lazard_modular_product_tail_certificates.py};
  \item \code{generate_lazard_root_fourier_numerator_p2_sparse.py}.
\end{itemize}

They do more than print constants: they choose shard boundaries, canonicalize
supports, name dependencies, and emit source whose import graph composes the
certificate.  Consequently they are build-system generators as well as
mathematical calculators.  Deterministic output and a clean regeneration diff
are important maintenance properties even though the generated truth is
rechecked.

The repository also keeps expository and crosswalk artifacts such as
\RepoFile{Algebra/PolynomialFormulas/LazardQuinticFormalization.tex} and the
paper-claim crosswalk.  These explain which formal module discharges which
step of a mathematical exposition.  For a development of this size, the
crosswalk is often a more useful index than the file tree.

\subsection{Checkpoint status at the pinned commit}

The commit message for the inspected snapshot records a concrete build
checkpoint rather than claiming the entire long-term Lazard program finished.
Among the recorded milestones were:

\begin{itemize}
  \item a bounded eight-double-coset certificate stack;
  \item a focused-kernel-green five-class \(S_5\) classifier and four-class
        corollary;
  \item the bridge from classification to resolvent criteria;
  \item Dummit table-zero certificates through tail 22 and a specified term
        frontier;
  \item a deterministic 563-module root plan with three bounded product
        shards;
  \item a regenerated 28-page formalization report.
\end{itemize}

The same checkpoint still identifies pending monolithic root-state data and
later terms.  These statements are best read as build-plan coordinates:
specific module families were green, while the overall construction retained
a known frontier.

\subsection{How to read the large project without drowning}

\begin{readingbox}
\begin{enumerate}
  \item Start with \RepoFile{Algebra/PolynomialFormulas/README.md} and the
        facade \RepoFile{Algebra/PolynomialFormulas/Lean/PolynomialFormulas.lean}.
  \item Read \code{Basic}, \code{Cubic}, and \code{Quartic} to learn the
        abstract radical convention.
  \item Choose one executable path:
        \bcode{GaussianPolynomialSolver} plus approximation modules.
  \item Read \code{AbelRuffini} and one explicit example to understand imported
        Galois-theory dependencies.
  \item For decision procedures, trace
        factor search $\to$ resolvent criterion $\to$ computability wrapper.
  \item Before any generated file, read
        \code{ComputableDummitCoefficientsCore.lean}.
  \item Inspect one \code{Term...Certificate} and one
        \code{Tail...Certificate}; use them as templates for the rest.
  \item Use the formalization report/crosswalk to locate the current
        mathematical endpoint and exact residue.
\end{enumerate}
\end{readingbox}

\subsection{Maintenance hazards}

\begin{gotchabox}
The biggest hazards in this subtree are architectural rather than algebraic.

\begin{itemize}
  \item A generated module can exist but be unreachable from the facade or
        intended build target.
  \item Reordering or merging shards can exceed recursion/heartbeat limits.
  \item A generator can regenerate a syntactically different but extensionally
        equal normal form, producing enormous noisy diffs.
  \item A small bridge theorem can carry the real mathematical meaning while
        hundreds of data modules merely support equality.
  \item ``The file compiles'' is weaker than ``the endpoint imports the file.''
  \item A fast native replay and a kernel certificate are different validation
        tiers and should not be substituted silently.
\end{itemize}
\end{gotchabox}

\begin{sourcebox}
\RepoFile{Algebra/README.md};
\RepoFile{Algebra/JacobianConjecture/README.md};
\RepoFile{Algebra/JacobianConjecture/Lean/JacobianConjecture/Counterexample.lean};
\RepoFile{Algebra/JacobianConjecture/Lean/JacobianConjecture/Audit.lean};
\RepoFile{Algebra/PolynomialFormulas/README.md};
\RepoFile{Algebra/PolynomialFormulas/Lean/PolynomialFormulas.lean};
\RepoFile{Algebra/PolynomialFormulas/Lean/PolynomialFormulas/ComputableDummitCoefficientsCore.lean};
\RepoFile{Algebra/PolynomialFormulas/Lean/PolynomialFormulas/ComputableDummitKernelCertificateTable0Tail22Data.lean};
\RepoFile{Algebra/PolynomialFormulas/Lean/PolynomialFormulas/ComputableDummitKernelCertificateTable0Tail22Certificate.lean};
\RepoFile{Algebra/PolynomialFormulas/LazardQuinticFormalization.tex}.
\end{sourcebox}

\chapter{Analysis}
\label{chap:analysis}

\section{Overview}

The analysis subtree is compact enough to read end to end.  It contains
\project{ExponentialIdentities} and \project{TrigonometricIdentities}, each
with a small facade, a handful of focused proof modules, and paired Lean/Rocq
work where appropriate.  Unlike the certificate-heavy combinatorics and
algebra projects, these developments are dominated by analytic inequalities
and exact symbolic transformations rather than generated source.

\section{ExponentialIdentities}

The standout module is
\RepoFile{Analysis/ExponentialIdentities/Lean/ExponentialIdentities/TinyExponentTower.lean}.
It proves the exact floor identity

\[
\left\lfloor
10^{\,10^{\,10^{\,10^{\,10^{-10^{10}}}}}}
\right\rfloor
-10^{10^{10}}
= 2\,811\,012\,357\,389.
\]

The expression is visually enormous but analytically close to a first-order
perturbation.  The file introduces

\[
N=10^{10},\qquad
t=10^{-N},\qquad
L=\log 10,
\]
then names the successive tower layers \(u,v,c\) and the final real value.
Its proof strategy is more instructive than a direct interval evaluation.

\subsection[Step 1: an exact rational enclosure for log 10]{Step 1: an exact rational enclosure for \(\log 10\)}

The module proves \(10^{-20}\)-scale bounds for \(\log 2\) and
\(\log(5/4)\) using Mathlib's finite logarithmic-series remainder theorem.
It combines

\[
\log 10 = 3\log 2+\log(5/4)
\]

to obtain a narrow rational interval around

\[
\frac{230258509299404568403}{10^{20}}.
\]

This interval is then raised to the fourth power with monotonicity lemmas.
The key first derivative quantity is shown to satisfy

\[
M+\frac{2}{5}
< 10^{11}L^4
< M+\frac{1}{2},
\qquad
M=2\,811\,012\,357\,389.
\]

The integer in the final answer is therefore not guessed from a floating-point
tower; it is isolated by a rational inequality for the derivative coefficient.

\subsection{Step 2: propagate the tiny increment}

For \(0\le x\le1\), the module derives

\[
0\le e^x-1-x\le x^2.
\]

A reusable \code{exp_layer} lemma says that if an increment \(w\) is bounded by
\(K t\), then one additional exponentiation has a nonnegative quadratic
remainder and a new increment bounded by \(4Kt\).  The three inner exponential
layers instantiate this lemma with moderate constants.  Since
\(t=10^{-10^{10}}\), even deliberately crude constants annihilate all
higher-order terms.

The proof never attempts to materialize \(10^{10^{10}}\) as a natural number.
It works with real exponentiation, logarithmic identities, and order
properties; the gigantic scale appears symbolically in powers whose
monotonicity is easy to prove.

\subsection{Step 3: convert a logarithmic increment to a floor}

The final layer is written as a base point \(10^N\) multiplied by an
exponential perturbation.  Lower and upper inequalities show that the excess
lies strictly between \(M\) and \(M+1\).  Standard floor lemmas then give the
integer identity.

\begin{overviewbox}
The module is a model for formalizing ``numerically obvious but astronomically
large'' identities: isolate a small parameter, prove a coarse analytic
remainder bound, reserve high precision only for one moderate constant, and
avoid evaluating the large baseline altogether.
\end{overviewbox}

The Rocq counterpart uses Coquelicot/Interval infrastructure for the same exact
claim.  The two implementations have comparable mathematical decomposition
but different interval machinery and trust profiles.

\section{TrigonometricIdentities}

The trigonometric project contains exact identities such as an arctangent
square relation and golden-ratio trigonometry.  Representative modules include
\code{ArctanSquareIdentity.lean} and \code{TrigGoldenRatio.lean}, with Rocq
counterparts.

The recurring pattern is:

\begin{enumerate}
  \item reduce the displayed identity to algebraic statements about
        \(\sin\), \(\cos\), or \(\tan\);
  \item prove domain and sign conditions needed by inverse trigonometric
        functions;
  \item use exact polynomial identities for the special constants;
  \item invoke injectivity or interval-restricted inverse lemmas to recover the
        target angle.
\end{enumerate}

The sign and interval obligations are the nontrivial formalization cost.
A paper proof may silently select a quadrant after squaring or applying
\(\tan\); the formal proof makes that branch selection explicit.

\section{How the analysis projects fit the repository}

These projects are useful reference points because they contain almost none of
the monorepo's generated-certificate complexity.  Their build graph is close
to their mathematical argument, and their source files demonstrate:

\begin{itemize}
  \item local helper lemmas for exact constants;
  \item careful control of inverse-function side conditions;
  \item rational enclosure rather than floating-point evidence;
  \item scoped resource options rather than global relaxation;
  \item paired proof-assistant results without a generator layer.
\end{itemize}

\begin{readingbox}
Read \code{TinyExponentTower.lean} in four passes: declarations and theorem
statement; logarithm enclosures; generic exponential-layer lemmas; final
floor argument.  Reading it linearly on the first pass obscures the fact that
the proof is one reusable perturbation lemma instantiated three times.
\end{readingbox}

\begin{sourcebox}
\RepoFile{Analysis/README.md};
\RepoDir{Analysis/ExponentialIdentities};
\RepoFile{Analysis/ExponentialIdentities/Lean/ExponentialIdentities/TinyExponentTower.lean};
\RepoDir{Analysis/TrigonometricIdentities};
\RepoFile{Analysis/TrigonometricIdentities/Lean/TrigonometricIdentities/ArctanSquareIdentity.lean};
\RepoFile{Analysis/TrigonometricIdentities/Lean/TrigonometricIdentities/TrigGoldenRatio.lean}.
\end{sourcebox}

\chapter{Combinatorics}
\label{chap:combinatorics}

\section{Overview}

The combinatorics subtree contains four different kinds of finite reasoning:

\begin{enumerate}
  \item recursive syntax and value-set enumeration for expression trees;
  \item exact finite geometric certificates;
  \item finite injections that prove coefficient recurrences;
  \item asymptotic arguments that convert recurrences or
        supermultiplicativity into growth bounds.
\end{enumerate}

The main projects are expression enumeration (power towers and radical
expressions), derivative-expansion sequences, perfect squared squares, and
Klarner's polyomino growth constant.

\section{Expression enumeration}

\subsection{PowerTowers}

The facade
\RepoFile{Combinatorics/ExpressionEnumeration/PowerTowers/Lean/PowerTowers.lean}
imports a shared core and sequence-specific developments for A000081, A002845,
A199812, and A198683.  It also imports selected small-value certificates and a
substantial \(n=12\) analysis for A198683.

The shared core separates:

\begin{itemize}
  \item binary parenthesization syntax;
  \item structural size and recursive decomposition;
  \item interpretation under a chosen base and exponentiation convention;
  \item finite sets or multisets of attainable values;
  \item counting quotients where distinct syntax can denote the same value.
\end{itemize}

This distinction is essential.  Counting binary trees alone is a Catalan-style
problem; counting \emph{distinct values} requires a semantic equality analysis
whose difficulty depends heavily on the base.

\subsection[A198683: principal values of towers of i]{A198683: principal values of towers of \(i\)}

A198683 counts distinct principal complex values of all binary
parenthesizations of \(i^i\cdots{}^i\).  The project defines
\code{principalPow} through the principal complex logarithm and recursively
collects value sets.  Small cases combine finite-set decomposition with
analytic proofs that candidate values are distinct.

For \(n=7\), for example, the proof constructs several explicitly separated
families.  Injectivity of maps such as \(z\mapsto i^z\) is established on
carefully bounded real strips; signs of real or imaginary parts separate image
families; finite cardinalities then supply the lower bound matching the
computed upper bound.

This architecture avoids a fragile list of pairwise ``not equal'' proofs.
A geometric invariant---strip injectivity, norm, or sign---separates whole
families at once.

\subsection[Generated data and the n=12 frontier]{Generated data and the \(n=12\) frontier}

The physical layout under \code{A198683} is informative:

\begin{itemize}
  \item \code{SmallValues} contains hand-guided and bounded certificates;
  \item \code{DataCertificates} contains imported/generated rows and replay
        lemmas;
  \item \code{N12} contains probe, magnitude, overflow, symbolic,
        certificate, endpoint, and bound modules.
\end{itemize}

The main \code{A198683.lean} comment says that the accepted OEIS data field
runs through \(n=11\) and deliberately withholds \(n=12\) while a
proof-quality final certificate is outstanding.  At the same snapshot, the
separate \(n=12\) subtree has made substantial progress: endpoint estimates
for a near-\(1\) split are proved unconditionally by rational boxing,
monotone transport, and Taylor bounds.  The endpoint module even records two
earlier false numerical constants and proves their falsity before replacing
them.

This is an excellent status pattern: a difficult intermediate can become fully
formal without silently upgrading the final sequence-value claim.

\begin{boundarybox}
For A198683, floating-point clustering is a discovery aid only.  Distinctness
must ultimately come from exact symbolic relations or certified analytic
separation, because principal complex exponentiation can bring enormous
intermediate magnitudes back near one another.
\end{boundarybox}

\subsection{A158415 and derivative expansions}

The radical-expression project for A158415 has Lean and Rocq trees plus
research/support material.  Its architecture parallels power towers:
recursive expression syntax, normalization/evaluation, finite enumeration, and
certificates for initial sequence values.

The derivative-expansion work around A290268 lives under
\RepoDir{Combinatorics/DerivativeExpansions} with an auxiliary top-level
\RepoDir{Oeis/A290268} surface.  It formalizes the combinatorics of iterated
derivatives and the resulting coefficient sequence.  These are smaller
projects, but they reuse the same ``syntax first, semantic quotient second''
discipline.

\section{SquaredSquare}
\label{sec:squared-square}

\subsection{What is proved}

The project formalizes A.~J.~W. Duijvestijn's order-21 dissection of a square
of side 112 into pairwise noncongruent integer-sided squares.  The side lengths
are

\[
2,4,6,7,8,9,11,15,16,17,18,19,24,25,27,29,33,35,37,42,50.
\]

The Lean and Rocq developments independently prove:

\begin{itemize}
  \item the listed placements cover the \(112\times112\) square exactly;
  \item open interiors are pairwise disjoint;
  \item all side lengths are positive and distinct;
  \item distinct positive side lengths imply noncongruence;
  \item scaling yields a 21-piece perfect dissection of every positive-side
        square;
  \item every perfect squared square has at least seven pieces.
\end{itemize}

Therefore the minimum possible order is unconditionally bracketed by
\[
  7 \leq m \leq 21.
\]

\subsection{Geometry reduced to integer data}

\code{Basic.lean} defines an axis-parallel closed square
\([x,x+s]\times[y,y+s]\), its open box, dissections, and congruence.  It proves
that the explicit open box is the topological interior of the closed square.
Congruence is handled through squared Euclidean distance: a positive-side
square has squared diameter \(2s^2\), so congruent squares have equal sides;
equal sides are congruent by translation.

The finite placement proof then works over integers:

\begin{enumerate}
  \item every piece lies inside the bounding grid;
  \item every pair of open rectangles is separated;
  \item the side list has no duplicates;
  \item every unit grid cell belongs to exactly one piece.
\end{enumerate}

Lean discharges the certificate with \code{decide +kernel}; Rocq uses
\code{vm_compute}.  Elementary floor and integer-part lemmas lift cell coverage
to exact coverage of the real square.

\subsection{The elementary lower bound}

\code{Minimality.lean} formalizes the classical corner/edge argument.

\begin{enumerate}
  \item Each corner of the outer square is covered by a unique corner tile.
  \item No piece can have the full outer side, and the four corner tiles are
        pairwise distinct; hence at least four pieces.
  \item Tiles meeting one edge partition that edge, so their side lengths sum
        to the outer side.  The reusable one-dimensional theorem lives in
        \code{Intervals.lean}.
  \item With at most two non-corner pieces, edge-incidence cases force either
        equal corner sides or a sum of positive lengths equal to zero.
\end{enumerate}

The result excludes orders 1 through 6.  It does not replay Duijvestijn's
historical exhaustive network search excluding orders 7 through 20.

\subsection{Sharp minimality remains a named claim}

\code{MinimalOrder.lean} defines \code{DuijvestijnSearchClaim} and proves the
sharp statement ``the least achievable order is 21'' equivalent to that
claim.  The existence half and the lower bound seven are already formal; the
complete planar-network enumeration is not.

The README spells out what a full formal search would require:
commensurability, a complete enumeration of relevant planar structures, and a
verified linear-algebra check for each structure.  This is a precise future
proof obligation, not an informal ``computer search omitted'' note.

\subsection{Module path}

\begin{longtable}{L{.28\textwidth}L{.61\textwidth}}
\caption{SquaredSquare modules.}
\label{tab:squared-square-modules}\\
\toprule
Module & Responsibility \\
\midrule
\endfirsthead
\toprule
Module & Responsibility \\
\midrule
\endhead
\code{Basic} & Geometric definitions, interiors, diameter, and congruence.\\
\code{Duijvestijn} & Concrete 21-piece placement and kernel certificate.\\
\code{Scaling} & Dilation closure for arbitrary positive side.\\
\code{Intervals} & One-dimensional interval-cover sum theorem.\\
\code{Minimality} & Exclusion of perfect dissections with at most six pieces.\\
\code{MinimalOrder} & Least-order statement and reduction to the search claim.\\
\code{Audit} / \code{Coq/Audit.v} & Assumption reporting.\\
\bcode{Support/bouwkamp_decode.py} & Bouwkamp-code decoding and data checks only.\\
\bottomrule
\end{longtable}

\section{KlarnerConstant}
\label{sec:klarner}

\subsection{Headline result}

The project proves the exact upper bound
\[
  \lambda \leq \frac{9047}{2000}=4.5235
\]
for Klarner's polyomino growth constant.  The numerical improvement changes
no geometric recurrence from the cited seventeen-neighborhood system; it
supplies a new exact rational supersolution at
\(\zeta=2000/9047\).  All seventeen coordinates share denominator \(10^7\).

The formal project goes well beyond checking those seventeen inequalities.
It defines the actual fixed-polyomino counts, proves all required geometric
recurrences by finite injections, establishes supermultiplicativity, and
connects the pointwise exponential estimate to the conventional asymptotic
growth limit.

\subsection{The five layers}

\begin{figure}[htbp]
\centering
\begin{tikzpicture}[
  b/.style={draw=blue!45,fill=sky,rounded corners=1mm,align=center,
            minimum width=31mm,minimum height=12mm,font=\sffamily\small},
  e/.style={draw=teal!50,fill=mint,rounded corners=1mm,align=center,
            minimum width=31mm,minimum height=12mm,font=\sffamily\small},
  arrow/.style={-{Latex[length=2mm]},thick,draw=muted},
  node distance=6mm
]
\node[b] (objects) {polyominoes, patterns,\\translation classes};
\node[b,right=of objects] (inj) {finite injections and\\deletion partitions};
\node[b,right=of inj] (rec) {17 coefficient\\recurrences};
\node[b,right=of rec] (cert) {exact rational\\supersolution};
\node[e,below=8mm of cert] (growth) {pointwise bound, growth supremum,\\and nth-root convergence};
\draw[arrow] (objects) -- (inj);
\draw[arrow] (inj) -- (rec);
\draw[arrow] (rec) -- (cert);
\draw[arrow] (cert) -- (growth);
\end{tikzpicture}
\caption{KlarnerConstant proof flow.  The final numeric certificate is small; proving that the actual geometric counts satisfy its recurrences is the bulk of the development.}
\label{fig:klarner-flow}
\end{figure}

\paragraph{Objects and counting.}
\code{Polyomino.lean}, \code{Patterns.lean}, and \code{Counting.lean} define
finite edge-connected lattice animals, southwest normalization, required and
forbidden offset patterns, and marked-occurrence counts.  The seventeen
coefficient sequences are therefore concrete finite cardinalities, not
uninterpreted sequences postulated to satisfy a recurrence.

\paragraph{Translation quotient.}
\code{TranslationClasses.lean} constructs translation equivalence as a
quotient, proves that each class has a unique southwest-normalized
representative, and identifies the fixed count with the quotient cardinality.

\paragraph{Connected decomposition.}
\code{SeededPartition.lean} grows disjoint nonempty seeds inside a connected
cell set into disjoint connected territories covering the original.  The
geometry modules use these territories as target factors in explicit
injections.

\paragraph{Recurrence algebra.}
\code{Convolution.lean}, \code{Recurrence.lean}, \code{BuiSystem.lean},
\code{CoefficientSystem.lean}, and \code{PublishedSystem.lean} distinguish the
paper's literal indexing convention from an all-index algebraic encoding.
They prove adapters between the two instead of silently changing base cases.

\paragraph{Exact certificate.}
\code{Certificate.lean} defines the seventeen-variable polynomial map and
checks the rational supersolution with ordinary kernel-reduced
\code{norm_num}.  Python and Wolfram scripts independently replay the same
rational inequalities, but Lean's certificate is the theorem boundary.

\subsection{The geometry module families}

The recurrence proof is split by neighborhood type and proof mechanism.

\begin{longtable}{L{.30\textwidth}L{.59\textwidth}}
\caption{Klarner geometric families.}
\label{tab:klarner-geometry}\\
\toprule
Family & Role \\
\midrule
\endfirsthead
\toprule
Family & Role \\
\midrule
\endhead
\code{GeometricLinear} &
Same-size occupancy partitions for the simplest coefficient identities.\\
\code{GeometricDeletion} &
Injective marked-leaf deletion for several one-cell inequalities.\\
\bcode{GeometricPBasics/Core/Geometry/Endpoint} &
Multi-branch \(P\) decomposition, recovery, and endpoint.\\
\bcode{GeometricQCore/Geometry/Endpoint} &
Analogous \(Q\) chain.\\
\code{GeometricFourFive...} &
The multi-stage \(S\) inequality.\\
\code{GeometricU...}, \code{GeometricV...}, \code{GeometricW...} &
Separate connected-decomposition chains for the corresponding types.\\
\code{GeometricTwoDeletion} &
One- and two-cell deletion maps for \(X,Y,Z\), retaining enough coordinates
for reconstruction.\\
\code{GeometricComplete} &
Assembles all seventeen concrete recurrences and applies the certificate
without recurrence or domination hypotheses in the final statements.\\
\bottomrule
\end{longtable}

The recurring proof obligation is injectivity.  It is not enough to map a
marked polyomino to several connected territories; the target tuple must
retain enough marks, offsets, and normalization data to reconstruct the
source uniquely.  ``Geometry'' modules define the map; ``Core'' and
``Endpoint'' splits keep reconstruction and cardinal arithmetic manageable.

\subsection{Asymptotic endpoint}

\code{Concatenation.lean} vertically stacks normalized polyominoes injectively,
proving supermultiplicativity of the fixed count.  \code{Asymptotic.lean}
applies Fekete's lemma to negative logarithms, identifies the limit with the
supremum of positive-index \(n\)th roots, and derives convergence of the
standard real roots.

\code{Growth.lean} separately proves the generic passage from a pointwise
exponential bound to the growth supremum.  \code{Main.lean} remains reusable
and accepts abstract recurrence/domination hypotheses.  The unconditional
composition is \code{GeometricComplete.lean}.

\subsection{Trust and transcription}

The formal argument has no project axioms, \code{sorry}, unsafe definitions,
or native-decision boundary according to the project README and audit.
However, the seventeen finite coordinate tables in \code{Patterns.lean} were
transcribed from source diagrams.  Lean proves the theorem relative to those
tables; the human comparison with the paper's page is a separate provenance
step.

This is exactly the right place for independent support scripts: they can
replay arithmetic and detect transcription inconsistencies without being
confused with the proof term.

\section{Reading routes}

\begin{readingbox}
\textbf{To understand finite certificate style:}
read \project{SquaredSquare} in the order
\code{Basic} $\to$ \code{Duijvestijn} $\to$ \code{Scaling}
$\to$ \code{Minimality}.

\textbf{To understand finite-injection recurrence style:}
read \project{KlarnerConstant} in the order
\code{Polyomino}/\code{Patterns} $\to$ \code{SeededPartition}
$\to$ one small geometry family $\to$ \code{CoefficientSystem}
$\to$ \code{Certificate} $\to$ \code{GeometricComplete}.

\textbf{To understand analytic distinctness around generated values:}
read \project{PowerTowers} \code{Core}, one small A198683 value proof, then the
\(n=12\) symbolic and endpoint modules.
\end{readingbox}

\begin{sourcebox}
\RepoFile{Combinatorics/README.md};
\RepoFile{Combinatorics/ExpressionEnumeration/PowerTowers/Lean/PowerTowers.lean};
\RepoFile{Combinatorics/ExpressionEnumeration/PowerTowers/Lean/PowerTowers/A198683.lean};
\RepoFile{Combinatorics/ExpressionEnumeration/PowerTowers/Lean/PowerTowers/A198683/N12/N12Endpoints.lean};
\RepoFile{Combinatorics/SquaredSquare/README.md};
\RepoFile{Combinatorics/Polyominoes/KlarnerConstant/README.md}.
\end{sourcebox}

\chapter{Computability}
\label{chap:computability}

\section{Overview}

The computability subtree contains three projects with almost disjoint
mathematical styles:

\begin{description}
  \item[\project{BusyBeaver}] Executable finite-state machine semantics,
        trace/classification certificates, and explicit control of expensive
        build targets.
  \item[\project{CombinatoryLogic}] Hand-written operational semantics,
        compilers among lambda, SK, SKI, and Iota calculi, and universality
        simulations.
  \item[\project{TuringDegrees}] Higher computability theory over oracle
        reducibility, with intentionally different Lean and Rocq
        presentations and a claim-by-claim coverage ledger.
\end{description}

Together they show why ``computability formalization'' is not one architecture:
finite execution, simulation of partial recursive functions, and abstract
degree theory require different interfaces.

\section{BusyBeaver}
\label{sec:busy-beaver}

\subsection{The facade is intentionally small}

The Lean facade \code{BusyBeaver.lean} imports only
\code{BusyBeaver.Core} and \code{BusyBeaver.KnownValues}.  A source comment
states that the expensive BB2 and BB3 certificates are opt-in.

This is a deliberate answer to a monorepo problem: users who need the machine
model or known-value API should not pay for all classification shards.
Separate libraries and executable generators remain available for focused
builds.

\subsection{Machine model}

\code{Core.lean} lives in namespace
\code{SetTheory.BusyBeaver}.  It defines:

\begin{itemize}
  \item head moves and tape actions;
  \item a machine transition table indexed by state and read symbol;
  \item typed and untyped machine presentations;
  \item a finite-support tape represented by lists around the head;
  \item configurations, one-step execution, and multi-step traces;
  \item halting and score functions;
  \item compiler/bridge interfaces used by classification code.
\end{itemize}

The tape representation is optimized for formal execution rather than for
extensional infinite sequences.  Blank tails are implicit.  Read, write, and
shift operations have executable definitions whose semantic lemmas are small
enough to reuse in trace proofs.

\subsection{Score is not runtime}

The project distinguishes at least two natural Busy Beaver quantities:

\begin{itemize}
  \item \textbf{halting time}: number of transitions before halt;
  \item \textbf{score}: number of marked symbols (or the project-specific
        output measure) at halt.
\end{itemize}

A machine can maximize one without maximizing the other.  The theorem names
and README keep these measures separate, and classification certificates must
state which optimization they establish.  This is particularly important
when importing upstream BB certificates whose machine convention or objective
may differ.

\subsection{Classification and certificate modules}

The internal Lean tree includes \code{BB2}, \code{BB3}, and \code{BB4}
families, with generated shards and audit modules.  Small-state classification
typically proceeds by:

\begin{enumerate}
  \item normalize machines modulo explicit symmetries;
  \item partition the finite machine space;
  \item associate each representative with a halting trace, a domination
        argument, or a nonhalting classification;
  \item replay the finite certificate with kernel \code{decide};
  \item lift representative results through the symmetry/normalization
        theorem;
  \item derive the score bound and exhibit a matching champion.
\end{enumerate}

Executable Lean programs such as
\code{GenerateBB4R05A14Guided.lean} and
\code{GenerateBB4R13C14Guided.lean}, plus PowerShell generation scripts,
produce candidate certificate material.  They are support tools; the checked
trace/classification theorems are the intended boundary.

The project status distinguishes completed kernel classifications from
conditional or incomplete BB4 root claims.  A facade that omits a module is
not itself evidence that the result is open; the project README and audits
state the exact boundary.

\subsection{Mathlib and Rocq bridges}

\code{Mathlib.lean} connects the repository-local machine model to relevant
Mathlib computability notions.  Separate audit modules can report assumptions
for local theorems and imported Mathlib facts independently.

The Rocq side combines repository-authored models/bridges with selected
certificates vendored under \RepoDir{lib/Coq-BB5}.  The vendored README records
the upstream commit and local kernel-hardening changes.  Repository-authored
work remains under \RepoDir{Computability/BusyBeaver}, preserving provenance.

\begin{readingbox}
Read \code{Core.lean} before any certificate.  Then inspect one smallest
classification path end to end: machine normalization, one concrete trace,
certificate checker, and final bound.  Only after that should the generated
BB3/BB4 shard tree be read; otherwise the data volume hides the machine
semantics that gives it meaning.
\end{readingbox}

\section{CombinatoryLogic}
\label{sec:combinatory}

\subsection{Semantic families}

The project independently formalizes several calculi:

\begin{itemize}
  \item a scoped untyped lambda calculus with capture-avoiding substitution;
  \item pure SK terms and reduction;
  \item SKI terms and reduction;
  \item pure Iota terms;
  \item arithmetic and coding inside SKI;
  \item external computability interfaces for partial recursive functions.
\end{itemize}

The Lean module list makes the intended proof graph visible:
\code{Lambda}, \code{StrongLambda}, \code{Reduction}, \code{SK},
\code{SKI}, \code{Iota}, compilation modules, coding/evaluation modules,
confluence modules, partial-recursive bridges, \code{Universality}, and
\code{Audit}.

\subsection{Lambda to SK: bracket abstraction}

\code{LambdaToSK.lean} implements bracket abstraction.  A free variable is
eliminated structurally:

\begin{itemize}
  \item the variable itself maps to \(I\);
  \item a term not containing the variable maps to \(K\,t\);
  \item an application maps to \(S\) applied to the abstractions of its two
        children.
\end{itemize}

The key theorem is semantic, not just syntactic: applying the compiled SK term
to an argument simulates substitution in the original lambda term.  The
compiler for closed lambda terms follows by iterating the abstraction step.

The scoped representation is important.  It allows the substitution theorem
to state exactly which variables are available at each term and prevents
capture by construction or by proved renaming discipline.

\subsection{SK, SKI, and Iota}

The standard one-combinator Iota encoding is recorded explicitly.  In one
common notation the derived combinators are:

\[
 I = \iota\iota,\qquad
 K = \iota(\iota(\iota\iota)),\qquad
 S = \iota(\iota(\iota(\iota\iota))).
\]

The project proves positive simulations among pure Iota, SK, and SKI rather
than identifying their syntax.  It also proves useful injectivity properties
of compilers, ensuring that a target encoding does not collapse distinct
source terms accidentally.

\code{IotaToLambda.lean} gives exact one-, two-, and three-step beta head
traces for the derived combinators, then lifts them through contexts and
reflexive/transitive closures.  This is stronger and more reusable than a
single theorem saying ``Iota is universal'': downstream proofs can account for
operational cost and reduction context.

\subsection{Confluence and evaluation}

The SKI subtree distinguishes raw one-step reduction, parallel or
Church-style reduction, confluence, and executable evaluation.  Confluence
justifies treating normal forms as observationally stable where they exist;
an evaluator theorem connects an implementation to the inductive reduction
relation.

Coding modules assign natural numbers to syntax and prove primitive-recursive
properties of constructors, decoders, and reduction predicates.  Arithmetic
modules construct numerals and operations inside the calculus.  These are the
ingredients needed to connect a high-level universality proof to a standard
partial-recursive/Turing-machine notion.

\subsection{Partial-recursive compilation}

\code{PartrecToSKI.lean} and \code{ExactPartrecToSKI.lean} construct SKI terms
representing partial recursive functions.  The exact version tracks the
observation relation carefully enough to state successful and divergent
behavior rather than only an extensional existence theorem.

The Rocq side also connects combinatory logic to the vendored Library of
Undecidability.  It uses that library's own lambda semantics and machine
interfaces.  The project README explicitly refuses to compose this theorem
with the local Iota embedding without an operational bridge.

\begin{gotchabox}
Do not infer a commuting square merely because all four corners are called
``lambda,'' ``SKI,'' or ``Turing complete.''  A compiler theorem preserves a
specific reduction or observation relation.  Different libraries can use
weak-head, call-by-value, full beta, fuelled evaluation, or relational
semantics.  The missing bridge can be the hardest theorem in the square.
\end{gotchabox}

\subsection{What is and is not claimed}

The compilers are not claimed to be mutual inverses, reduction-reflecting in
every direction, or fully abstract for an observational equivalence.  Positive
simulation and universality are sufficient for the public endpoint.  Stating
these nonclaims prevents a user from treating an encoding theorem as an
isomorphism of calculi.

\section{TuringDegrees}
\label{sec:turing-degrees}

\subsection{Why Lean and Rocq differ}

The Lean project builds on Mathlib's oracle-computability infrastructure and
uses a quotient of set representatives by mutual Turing reducibility.  The
Rocq project uses predicates plus a setoid/equivalence presentation because
its imported Synthetic Computability stack exposes different abstractions.

Both are legitimate degree constructions, but theorem transfer requires an
explicit representation bridge.  The repository therefore maintains a
coverage ledger instead of forcing line-by-line syntactic parity.

\subsection{Lean architecture}

The Lean development includes:

\begin{itemize}
  \item oracle reducibility and equivalence over sets/predicates;
  \item the quotient type of Turing degrees;
  \item order and lattice-like operations;
  \item zero degree and its characterization by computability;
  \item an exact even/odd join construction;
  \item complement invariance;
  \item the halting degree \(0'\) and c.e.-completeness;
  \item structural/cardinality results, including continuum many degrees and
        countable lower cones or equivalence classes where stated.
\end{itemize}

An ``exact join'' theorem matters more than a mere upper bound: it proves that
the even/odd tagged union is the least upper bound under reducibility, allowing
the quotient order to inherit familiar degree operations.

The quotient presentation makes equality pleasant for consumers but can hide
representative-level computation.  Internal lemmas therefore descend and lift
properties carefully through \code{Quot.sound}.

\subsection{Rocq architecture}

The Rocq development is deeper in several classical recursion-theory
directions.  According to its README and coverage ledger, it includes
components around:

\begin{itemize}
  \item the Turing jump and jump monotonicity;
  \item Kleene--Post constructions of incomparable degrees;
  \item a limit lemma and hierarchy machinery;
  \item versions of Post's theorem;
  \item a conditional or assumption-parameterized treatment of Post's problem;
  \item compatibility wrappers over Synthetic Computability.
\end{itemize}

Several hypotheses---Markov principles, enumeration principles, or effective
presentation assumptions---are visible in theorem statements rather than
folded into a global classical axiom.  That makes the Rocq side mathematically
richer in some directions but not a drop-in replacement for Lean's quotient
API.

\subsection{The coverage ledger}

\RepoFile{Computability/TuringDegrees/COVERAGE.md} maps informal claims to
specific Lean and Rocq declarations and lists their hypotheses.  It also
records deliberate omissions and corrects source-level shorthand.

One example is a recursive approximation: a correct formal definition needs
both the input and stage arguments, even when prose writes a one-argument
limit notation.  Another is pair coding: exact pairing/unpairing statements
are preferable to an informal ``computable coding exists'' phrase because
later reductions consume their equations.

The ledger's discipline is conservative.  A mathematically related theorem
elsewhere in the repository does not count as a port until its statement and
dependency boundary have been compared.

\subsection{Build boundary}

TuringDegrees has both a Lean audit and a Rocq setup script that builds the
pinned Synthetic Computability dependency with a repository-maintained
compatibility patch.  The patch lives beside the wrapper project, not inside
the submodule.  This preserves the submodule as an auditable upstream snapshot
while making the local toolchain reproducible.

\section{Comparative map}

\begin{longtable}{@{}L{.20\textwidth}L{.20\textwidth}L{.20\textwidth}L{.24\textwidth}@{}}
\caption{The three computability projects compared.}
\label{tab:computability-compare}\\
\toprule
Project & Primary objects & Main proof mechanism & Main boundary \\
\midrule
\endfirsthead
\toprule
Project & Primary objects & Main proof mechanism & Main boundary \\
\midrule
\endhead
BusyBeaver &
Finite transition tables and traces &
Executable replay and finite classification &
Generator output, evaluator choice, and opt-in certificate targets\\
CombinatoryLogic &
Terms, reductions, compilers, observations &
Induction, simulation, confluence, coding &
Operational semantics of each source/target calculus\\
TuringDegrees &
Oracles, reductions, equivalence classes/degrees &
Abstract reducibility, quotient/setoid reasoning, priority-style constructions &
Cross-prover representation and explicit logical principles\\
\bottomrule
\end{longtable}

\begin{sourcebox}
\RepoFile{Computability/README.md};
\RepoFile{Computability/BusyBeaver/README.md};
\RepoFile{Computability/BusyBeaver/Lean/BusyBeaver.lean};
\RepoFile{Computability/BusyBeaver/Lean/BusyBeaver/Core.lean};
\RepoFile{Computability/CombinatoryLogic/README.md};
\RepoDir{Computability/CombinatoryLogic/Lean/CombinatoryLogic};
\RepoFile{Computability/TuringDegrees/README.md};
\RepoFile{Computability/TuringDegrees/COVERAGE.md}.
\end{sourcebox}

\chapter{Logic}
\label{chap:logic}

\section{Overview}

Logic is the broadest top-level domain.  It ranges from tiny constructive
propositional lemmas to first-order completeness, modal-system libraries,
arithmetized syntax, interpretability, bounded consistency, and verified
decision procedures.

A useful internal taxonomy is shown in \cref{tab:logic-map}.

\begin{longtable}{L{.24\textwidth}L{.31\textwidth}L{.35\textwidth}}
\caption{Logic project families.}
\label{tab:logic-map}\\
\toprule
Family & Representative projects & Architectural emphasis \\
\midrule
\endfirsthead
\toprule
Family & Representative projects & Architectural emphasis \\
\midrule
\endhead
Propositional/equational microprojects &
ShefferStroke, NaturalDeduction, MonotonicityOfEntailment,
PrincipleOfExplosion, EquationalLogic, BooleanAlgebra,
FiniteMatrixNoncharacterizability &
Small syntax/semantics developments, often paired and constructive.\\
First-order metatheory &
FirstOrder, ArbitraryLanguageCompactness &
From-scratch fixed-language calculus versus a Mathlib-backed generic facade.\\
Source ports and modal logic &
FoundationPort, Modal &
Pinned source inventory, independent Rocq modules, and conservative coverage ledgers.\\
Interpretability &
PAHF, ZFCinPA &
Translations among arithmetic/set theories, model-relative coding, and
explicit conditional endpoints.\\
Peano arithmetic &
ListCoding, BoundedConsistency, NoFiniteModel,
NotFinitelyAxiomatizable/finite-basis reduction, Undecidable &
Arithmetized syntax, representability, partial truth, and proof predicates.\\
Decision procedures &
PresburgerArithmetic &
Syntax normalization, Cooper elimination, executable decision, and correctness.\\
Binder semantics &
QuantifierCommutation &
Constructive positive laws and finite countermodels for plausible false laws.\\
\bottomrule
\end{longtable}

\section{Propositional and equational microprojects}

These projects are compact laboratories for repository conventions.  Most
have a Lean facade, one or more definition/proof modules, a Rocq counterpart,
and an audit.

\subsection{Sheffer stroke and Boolean algebra}

The Sheffer-stroke project develops a functionally complete single connective
and translates ordinary Boolean operations into it.  The Boolean-algebra
project supplies algebraic semantics and equational reasoning.  Together with
\project{EquationalLogic}, they separate syntactic derivability from semantic
validity and give small soundness/completeness examples.

\subsection{Natural deduction and entailment}

\project{NaturalDeduction} defines proof objects or derivability rules for a
propositional language and proves the expected structural/metatheoretic
properties.  \project{MonotonicityOfEntailment} isolates the weakening fact
that adding hypotheses preserves entailment.  \project{PrincipleOfExplosion}
isolates derivability from contradiction.

The existence of separate microprojects for seemingly elementary facts is
useful: they avoid importing a large logic framework when testing a particular
encoding or cross-prover statement.

\subsection{Finite matrix noncharacterizability}

\project{FiniteMatrixNoncharacterizability} is a more substantial
propositional result.  Its presence beside the microprojects shows that the
top-level taxonomy is by subject, not by scale.  Readers should use its facade
and audit to find the actual public theorem rather than infer the argument from
the directory name alone.

\section{QuantifierCommutation}
\label{sec:quantifier-commutation}

This dependency-free project proves constructively that adjacent universal
quantifiers commute and adjacent existential quantifiers commute:

\[
(\forall x\,\forall y\,Rxy)\leftrightarrow
(\forall y\,\forall x\,Rxy),
\]
\[
(\exists x\,\exists y\,Rxy)\leftrightarrow
(\exists y\,\exists x\,Rxy).
\]

It then gives finite counterexamples for two superficially similar binders:
nested ``there does not exist'' and nested unique existence.

The first counterexample relies on parsing.  The nested expression

\[
\not\exists x\,(\not\exists y\,Rxy)
\]

is not the same as one outer negation of a two-variable existential block.
A two-element relation can have a witness in every row but an empty column,
separating the two binder orders constructively.

For unique existence, a three-element relation is chosen with exactly one row
containing exactly one related element, while every column contains exactly
one.  Thus the outer uniqueness count changes when binders are swapped.

The Lean core defines unique existence explicitly because dependency-free Lean
does not supply the notation in that environment.  Rocq uses its standard
\code{exists!}; the README verifies that the semantic contract matches.  No
classical axiom is needed.

This project is an excellent template for a small paired development:
statement, countermodel, audit, focused build, and \code{coqchk} all fit in one
README.

\section{FirstOrder}
\label{sec:first-order}

\subsection{A deliberately fixed language}

The repository-local first-order project uses a countable language with
equality and one distinguished binary relation, \code{fMem}.  That relation
can later serve as set membership, but the completeness proof itself treats it
abstractly.

Fixing the language keeps the development dependency-light and makes every
encoding explicit.  It is not intended to replace Mathlib's generic
first-order logic.  Instead, the root also exposes
\project{ArbitraryLanguageCompactness}, a separate audited facade over
Mathlib's arbitrary-signature theorem.

\subsection{Module ladder}

\begin{figure}[htbp]
\centering
\begin{tikzpicture}[
  b/.style={draw=blue!45,fill=sky,rounded corners=1mm,align=center,
            minimum width=30mm,minimum height=11mm,font=\sffamily\small},
  e/.style={draw=teal!50,fill=mint,rounded corners=1mm,align=center,
            minimum width=30mm,minimum height=11mm,font=\sffamily\small},
  arrow/.style={-{Latex[length=2mm]},thick,draw=muted},
  node distance=5mm
]
\node[b] (fol) {\code{Fol}\\syntax/semantics};
\node[b,right=of fol] (calc) {\code{Calculus}\\derivations};
\node[b,right=of calc] (comp) {\code{Completeness}\\Henkin construction};
\node[e,right=of comp] (class) {\code{ClassicalCompleteness}\\textbook endpoint};
\node[e,below=7mm of class] (compact) {\code{Compactness}};
\node[b,below=7mm of comp] (rel) {\code{Relativization}};
\node[e,below=7mm of calc] (audit) {\code{Audit}};
\draw[arrow] (fol) -- (calc);
\draw[arrow] (calc) -- (comp);
\draw[arrow] (comp) -- (class);
\draw[arrow] (class) -- (compact);
\draw[arrow] (fol) -- (rel);
\draw[arrow] (class) to[bend right=22] (audit);
\draw[arrow] (compact) to[bend left=18] (audit);
\end{tikzpicture}
\caption{The first-order dependency ladder.}
\label{fig:first-order-ladder}
\end{figure}

\paragraph{\code{Fol}.}
Defines terms, formulas, free-variable operations, structures, assignments,
term evaluation, and formula satisfaction.  Basic substitution and assignment
lemmas establish the semantic infrastructure consumed by soundness and the
truth lemma.

\paragraph{\code{Calculus}.}
Defines the proof system and derivability from theories.  Structural rules,
substitution, equality behavior, and soundness are proved before any model
construction begins.

\paragraph{\code{Completeness}.}
Builds a maximal consistent Henkin extension, constructs the term model,
quotients terms by provable equality, proves the truth lemma, and obtains a
model of a consistent theory.

\paragraph{\code{ClassicalCompleteness}.}
Packages the construction as the familiar semantic-entailment-implies-
derivability theorem and related consistency/model-existence equivalences.

\paragraph{\code{Compactness}.}
Derives compactness from completeness by finite inconsistency/derivation
support.

\paragraph{\code{Relativization}.}
Provides the formulas and semantic lemmas needed to restrict quantifiers to a
definable domain, a key dependency for internal model constructions in set
theory and interpretability.

\subsection{The Henkin path in more detail}

The completeness proof follows the standard mathematics but must expose
several hidden paper steps:

\begin{enumerate}
  \item enumerate formulas and add witnesses while preserving consistency;
  \item extend the resulting Henkin theory to a complete/maximal consistent
        theory;
  \item define a congruence on closed terms by provable equality;
  \item show function/relation interpretation is well defined on equivalence
        classes;
  \item prove the truth lemma by induction on formulas;
  \item derive satisfaction of the original theory;
  \item turn semantic entailment into a contradiction if the target formula
        were not derivable.
\end{enumerate}

The quotient step introduces \code{Quot.sound} and classical selection into
the usual Lean assumption footprint.  The audit makes this visible rather than
describing the project simply as ``classical.''

\subsection{Two compactness projects, two provenances}

\project{FirstOrder.Compactness} is a corollary of the repository's own
fixed-language completeness proof.  \project{ArbitraryLanguageCompactness}
re-exports or specializes Mathlib's general result and audits its assumptions.
The root facade imports both.

This is not duplication: one is an independently inspectable metatheory
development used by repository-local ZF and coding projects; the other gives
consumers a generic language interface without rebuilding Mathlib.

\section{FoundationPort and Modal}
\label{sec:foundation-port}

\subsection{Pinned source, independent target}

The submodule \RepoDir{lib/FormalizedFormalLogic-Foundation} is a pinned,
read-only Lean source reference.  \RepoDir{Logic/FoundationPort} tracks an
independent Rocq port.  ``Complete'' is defined conservatively: every source
module must be accounted for, not merely every theorem that happened to be
needed by a local endpoint.

\code{ported.tsv} stores reviewed mappings.  The script
\RepoFile{Logic/FoundationPort/coverage.sh} reads the pinned tree, validates
each mapping, and emits one row for each of the 455 Lean modules.  Status is
one of:

\begin{description}
  \item[\code{ported}] The theorem surface has an independently checked Rocq
        representation, possibly more general or idiomatic.
  \item[\code{partial}] A meaningful slice exists, but declarations or a
        dependency layer remain.
  \item[\code{unported}] No module-level parity claim has been reviewed.
\end{description}

The generated full report is deliberately not committed; the reviewed mapping
and pinned gitlink are the compact sources of truth.

\subsection{Declaration-free facades still count}

Foundation contains modules whose entire role is importing a family of modules.
The port ledger records exact or partial declaration-free Rocq facades for
those source modules.  This avoids a common coverage bug: counting only
declarations misses import surfaces that are part of a library's public API.

\subsection{The Modal port}

\RepoDir{Logic/Modal} is a large independent Rocq development corresponding to
high-value parts of Foundation's modal tree.  Its scope includes:

\begin{itemize}
  \item generic Kripke frames, models, generated submodels, and p-morphisms;
  \item modal syntax, substitutions, normal forms, and semantics;
  \item Hilbert systems and canonical-model completeness;
  \item systems such as K, T, D, B, K4, K5, S4, S5, GL, Grz, and extensions;
  \item filtration and finite-model arguments;
  \item frame correspondences and definability;
  \item algebraic semantics;
  \item boxdot and translation constructions;
  \item hierarchy and incompleteness-related components.
\end{itemize}

The README/module table names partial boundaries instead of claiming blanket
parity.  Certain systems or specialized source modules remain explicitly
outside complete coverage.  The port is not imported into the Lean source and
does not alter the Foundation submodule.

\begin{overviewbox}
FoundationPort treats coverage metadata as a theorem-library artifact.  A
build answers whether existing Rocq files check; the ledger answers whether
those files cover the pinned Lean source.  Both are necessary.
\end{overviewbox}

\section{Interpretability}
\label{sec:interpretability}

\subsection{PA/HF}

\project{PAHF} is a compact paired Lean/Rocq development relating Peano
arithmetic and hereditarily finite set theory.  It supplies translations and
model constructions showing how arithmetic data can be represented in finite
set-theoretic structure and conversely.

Its topic-level documentation is lighter than the neighboring ZFC project, so
the facade and audit are the best guides.  This asymmetry is real: not every
project in the monorepo has been normalized to one README template.

\subsection{ZFCinPA: target and status}

The ambitious target is a PA theorem asserting, internally and uniformly,
that for every complexity bound \(n\), ZFC has no proof of contradiction whose
formulas stay within that bound.  Informally:

\[
  \mathrm{PA}\vdash
  \forall n\;
  \mathrm{Prov}_{\mathrm{ZFC}}\bigl(\ulcorner
    \mathrm{Con}_{n}(\mathrm{ZFC})\urcorner\bigr).
\]

At the pinned snapshot the development does \emph{not} present this as an
unconditional theorem.  Its status can be summarized as follows.

\begin{longtable}{L{.46\textwidth}L{.35\textwidth}}
\caption{ZFCinPA status at the pinned commit.}
\label{tab:zfcinpa-status}\\
\toprule
Layer & Status \\
\midrule
\endfirsthead
\toprule
Layer & Status \\
\midrule
\endhead
For each external standard numeral \(n\), ZFC proves the corresponding
bounded-consistency sentence &
\kernelchecked\ in \project{BoundedZFCConsistency}.\\
The repository-local Foundation/ZFC layer proves the required instance &
\kernelchecked.\\
The uniform endpoint is reduced to an eight-field internal certificate &
\kernelchecked\ equivalence/reduction.\\
Successor implications at standard indices &
\kernelchecked.\\
Successor implications for arbitrary, possibly nonstandard model elements &
open premise.\\
Final PA uniform theorem &
\conditional\ on \bcode{ZFCSuccessorImplicationsInAllModels}.\\
\bottomrule
\end{longtable}

\subsection{Five architectural layers}

The README divides the project into a useful ladder.

\begin{enumerate}
  \item \textbf{Definability}: encode syntax, bounded satisfaction, proof
        predicates, and the relevant functions in arithmetic.
  \item \textbf{Statement}: define the exact PA sentence and model-relative
        semantic version.
  \item \textbf{Eight-field certificate}: package the local implications
        needed to step the construction through all indices.
  \item \textbf{Compiler/translation layer}: turn source derivations and
        semantic facts into the coded target claims.
  \item \textbf{Successor/evaluation bridge}: prove the model-internal
        successor step that must work even at nonstandard indices.
\end{enumerate}

An earlier five-field certificate was proved nontransferable to the desired
models.  The replacement adds two Tarski-style truth fields and a semantic
precheck, making the missing obligation explicit.  Recording a failed
abstraction this precisely is valuable: future work is directed at the actual
nonstandard-model bridge, not at polishing an insufficient package.

\subsection{Specification hardening}

Statement-bearing modules generally use

\begin{lstlisting}[language=Lean]
set_option autoImplicit false
\end{lstlisting}

so an out-of-scope identifier cannot silently become an additional theorem
parameter.  A small number of legacy exceptions are documented.  Companion
audit modules inspect the endpoint of each layer.

This policy arose from a real specification risk, not stylistic preference.
Metamathematical theorems with many codes, models, and translations are
especially vulnerable to a theorem that is easy only because an intended
constant became a free parameter.

\subsection{Lake target discovery gotcha}

The root \code{ZFCinPA} library has an explicit facade and no automatic glob
that turns every source file into a named Lake target.  Therefore a command
such as

\begin{lstlisting}[language=PowerShell]
lake build +ZFCinPA.NumeralOmega
\end{lstlisting}

can be a silent no-op if that module is not reachable from a registered target.
The reliable alternatives are to import it from the facade, add a target, or
elaborate directly:

\begin{lstlisting}[language=PowerShell]
lake env lean Logic/Interpretability/ZFCinPA/Lean/ZFCinPA/NumeralOmega.lean
\end{lstlisting}

This is a general monorepo lesson: a successful command that schedules zero
jobs proves nothing.  Inspect the target graph or build the source directly.

\section{Peano arithmetic}
\label{sec:peano-arithmetic}

\subsection{ListCoding}

\project{PAListCoding} independently develops Lean and Rocq formulas
representing finite lists and related arithmetic predicates.  The audited
surface contains thirty formulas: twenty-five standard list/number predicates
and five formulas used for hereditary Cantor-normal-form notation below
\(\varepsilon_0\).

Representative capabilities include:

\begin{itemize}
  \item validity of a list code, length, and element lookup;
  \item singleton, append/concatenation, and flattening;
  \item occurrence counting, permutation, and subsequence relations;
  \item sum, product, extrema, and simple statistics;
  \item primality, powers, factorization, digits, and divisors;
  \item Diophantine graphs for exponentiation, tetration, and higher
        hyperoperators;
  \item hereditary notation operations used in ordinal-analysis examples.
\end{itemize}

The public contract is representation and functionality, not equality of
numeric codes across proof assistants.  Lean uses an Ackermann/graph-of-
function approach integrated with Foundation infrastructure; Rocq uses a
nested polynomial code convenient for direct arithmetic proofs.

Two conventions are easy to miss: list positions are zero-based, while the
``\(n\)th prime'' interface is one-based.  The README and audit make those
choices explicit.

\subsection{Bounded PA consistency}

\project{BoundedConsistency} formalizes a scheme saying, for each external
complexity bound, that no PA proof of contradiction exists if
\emph{every formula occurrence in the proof} stays within that bound.

Four distinctions govern the whole project:

\begin{enumerate}
  \item \textbf{External scheme versus internal universal sentence.}
        Proving a theorem for each meta-level numeral \(n\) is weaker than one
        PA derivation of \(\forall n\,\mathrm{Con}_n\).
  \item \textbf{All occurrences versus conclusion only.}
        Intermediate formulas in a proof tree must satisfy the bound.
  \item \textbf{Standard codes versus nonstandard model elements.}
        A decoder theorem for quoted standard syntax may not control an
        internally represented nonstandard proof code.
  \item \textbf{Formula complexity versus proof length.}
        The bound is on formula rank/complexity, not the number of inference
        steps.
\end{enumerate}

\begin{figure}[htbp]
\centering
\begin{tikzpicture}[
  b/.style={draw=blue!45,fill=sky,rounded corners=1mm,align=center,
            minimum width=61mm,minimum height=16mm,text width=55mm,
            font=\sffamily\small},
  e/.style={draw=gold!50,fill=sand,rounded corners=1mm,align=center,
            minimum width=61mm,minimum height=16mm,text width=55mm,
            font=\sffamily\small},
  arrow/.style={-{Latex[length=2mm]},thick,draw=muted},
  node distance=8mm and 14mm
]
\node[b] (scheme) {For each external numeral \(n\):\\
  \(\mathrm{PA}\vdash\mathrm{Con}_n(\mathrm{PA})\)};
\node[e,right=of scheme] (uniform) {One internal theorem:\\
  \(\mathrm{PA}\vdash\forall n\,\mathrm{Con}(n)\)};
\node[b,below=of scheme] (std) {quoted standard formula and proof codes};
\node[e,below=of uniform] (nonstd) {arbitrary elements in a nonstandard PA model};
\draw[arrow] (scheme) -- node[above,font=\scriptsize\sffamily]{does not automatically yield} (uniform);
\draw[arrow] (std) -- node[above,font=\scriptsize\sffamily]{bridge required} (nonstd);
\end{tikzpicture}
\caption{Two independent uniformity gaps in bounded consistency.}
\label{fig:bounded-gaps}
\end{figure}

The implementation defines coded formulas, complexity, a seventeen-rule proof
tree, a verifier that checks every node, partial truth predicates, and
soundness of bounded derivations.  The Lean and Rocq numeralwise schemes are
complete.  The strongest uniform sentence is unconditional in Lean at this
snapshot, while the Rocq counterpart remains conditional on explicit compiler
premises; the README states that asymmetry rather than claiming parity.

\subsection{NoFiniteModel}

\project{NoFiniteModel} proves that a structure satisfying the chosen PA axiom
system cannot be finite.  The usual mathematical core is the injective but
non-surjective successor function.  Formalization requires aligning the
object-language axioms with the semantic structure and extracting the
successor properties under the exact axiom encoding.

\subsection{Finite-basis and undecidability projects}

The finite-basis reduction formalizes a route from finite axiomatizability
assumptions to the relevant contradiction/nonaxiomatizability theorem.  The
undecidability project connects PA derivability or theoremhood to a standard
computability obstruction.  Their inclusion as root default targets means the
workspace treats their audits as integration-critical even though the root
facade does not expose every internal module.

For both, read the statement module before the reduction machinery:
``undecidable'' can refer to provability, validity, consistency, or a coded
decision problem, and ``not finitely axiomatizable'' depends on whether
equivalence is semantic or deductive and over which class of models.

\section{PresburgerArithmetic}
\label{sec:presburger}

\subsection{Lean end-to-end decision procedure}

The Lean development implements Cooper elimination with the following module
flow:

\[
\code{Syntax}
\to \code{NormalForm}
\to \code{Cooper}
\to \code{Elimination}
\to \code{Decision}
\to \code{Audit}.
\]

\code{Syntax} defines the arithmetic language and semantics.  Normalization
moves formulas into the restricted linear form required by elimination.
\code{Cooper} handles one quantified variable by collecting coefficients,
least common multiples, divisibility atoms, and finite residue cases.
\code{Elimination} iterates the step.  \code{Decision} evaluates the resulting
quantifier-free formula.

Correctness is proved at each boundary, so the final Boolean procedure is
sound and complete for the input sentence.  No external solver or oracle is
called.

\subsection{Rocq scope}

The Rocq side independently verifies the decisive normalized one-variable
Cooper step and its supporting arithmetic.  Its theorem surface is not
identical to the Lean end-to-end executable decision.  This is documented as
a coverage boundary, not concealed by identical project names.

\section{How to navigate Logic safely}

\begin{readingbox}
\begin{enumerate}
  \item For a microproject, read facade $\to$ definitions $\to$ audit.
  \item For metatheory, draw the syntax--semantics--derivation--model ladder
        before opening proofs.
  \item For arithmetization, write every quantifier at the meta/object/model
        level and label standard versus nonstandard variables.
  \item For a source port, consult the coverage ledger before the module tree.
  \item For a decision procedure, identify the executable function, the
        semantic specification, and the theorem connecting them.
  \item For interpretability, locate the translation identity and the
        model-relative satisfaction theorem; theorem-name similarity is not a
        bridge.
\end{enumerate}
\end{readingbox}

\begin{sourcebox}
\RepoFile{Logic/README.md};
\RepoFile{Logic/QuantifierCommutation/README.md};
\RepoFile{Logic/FirstOrder/README.md};
\RepoFile{Logic/FoundationPort/README.md};
\RepoDir{Logic/Modal};
\RepoFile{Logic/Interpretability/ZFCinPA/README.md};
\RepoFile{Logic/PeanoArithmetic/ListCoding/README.md};
\RepoFile{Logic/PeanoArithmetic/BoundedConsistency/README.md};
\RepoFile{Logic/PresburgerArithmetic/README.md}.
\end{sourcebox}

\chapter{Number theory}
\label{chap:number-theory}

\section{Overview}

The number-theory subtree is modest in byte size but deliberately varied in
proof provenance.  Its projects are:

\begin{itemize}
  \item a formal exponent-four case of Fermat's Last Theorem;
  \item integer-sum identities;
  \item Calkin--Wilf enumeration of positive rationals;
  \item an arithmetic sentence expressing the Riemann hypothesis.
\end{itemize}

The key architectural lesson is that matching theorem names on the Lean and
Rocq sides do not imply matching proofs.

\section{DiophantineEquations and Fermat exponent four}

The Lean facade imports a very small \code{FermatFour.lean}.  Its role is to
specialize and expose the relevant Mathlib theorem under the repository's
chosen statement.  The formal evidence is kernel-checked, but the deep descent
argument belongs to the imported library.

The Rocq side contains an independent descent development.  A code reader
should therefore describe the two sides as:

\begin{itemize}
  \item \textbf{Lean}: local API and specialization over a pinned Mathlib
        theorem;
  \item \textbf{Rocq}: repository-local proof of the exponent-four descent.
\end{itemize}

This distinction matters when assessing maintainability.  The Lean wrapper is
small but tied to upstream theorem naming and statement shape.  The Rocq proof
is larger but locally inspectable and can expose intermediate number-theoretic
lemmas.

\section{IntegerSums}

The integer-sums project proves exact finite-sum identities in both proof
assistants.  It is a useful compact example of algebraic normalization,
induction over a finite range, and explicit handling of integer/natural casts.

These modules often look trivial informally but are valuable regression tests
for coercion and finite-sum conventions used in larger projects.

\section{RationalEnumeration}

The Calkin--Wilf development formalizes the binary-tree/orbit enumeration of
positive rational numbers.  Its architecture separates:

\begin{itemize}
  \item the successor map or left/right tree actions;
  \item positivity and reduced-fraction invariants;
  \item existence: every positive rational is reached;
  \item uniqueness: no positive rational is reached twice;
  \item the resulting enumeration/bijection statement.
\end{itemize}

Euclidean descent is the conceptual inverse: compare numerator and denominator,
subtract the smaller from the larger, and record the branch.  Formalization
requires a termination measure and exact reconstruction theorem.

Lean and Rocq use their native rational/integer libraries, so the internal
normal-form proofs differ while the public enumeration contract is parallel.

\section{RiemannHypothesis}

This project does not attempt analytic number theory inside PA.  It constructs
a first-order arithmetic sentence representing the Riemann hypothesis under
the selected coding and proves the expected statement-level relationships.

The main review questions are therefore syntactic:

\begin{itemize}
  \item Which external RH formulation is arithmetized?
  \item Which coding of finite sequences, rationals, or approximations is used?
  \item Is the theorem an equivalence in the standard model, a derivability
        statement, or merely a definition of the sentence?
  \item Which imported representability results connect the code to analysis?
\end{itemize}

The project belongs in number theory by intended statement and in logic by
implementation technique.  Its current location keeps the user-facing
mathematical taxonomy primary.

\begin{sourcebox}
\RepoFile{NumberTheory/README.md};
\RepoFile{NumberTheory/DiophantineEquations/README.md};
\RepoDir{NumberTheory/DiophantineEquations};
\RepoDir{NumberTheory/IntegerSums};
\RepoDir{NumberTheory/RationalEnumeration};
\RepoDir{NumberTheory/RiemannHypothesis}.
\end{sourcebox}

\chapter{Set theory}
\label{chap:set-theory}

\section{Overview}

The set-theory subtree provides a lightweight internal ZF base, an alternative
closure-schema axiomatization, and a bounded-complexity consistency theorem for
ZFC.  It also supplies dependencies consumed by the logic/interpretability
projects.

These developments use the repository-local fixed first-order language rather
than importing a large set-theory library.  This gives tight control over
formula codes, quantifier rank, relativization, and internal models---exactly
the features needed for bounded consistency.

\section{ZF}
\label{sec:zf}

\subsection{Public surface and internal scale}

The Lean facade is tiny: it imports \code{ZF.Zf}.  The implementation module is
substantial and is paired with an audit.  It defines the set-theoretic formulas
and model infrastructure needed downstream, including:

\begin{itemize}
  \item the ZF/ZFC axiom formulas in the one-relation language;
  \item models and satisfaction of those axioms;
  \item internal empty set, pairing, union, separation, power set, infinity,
        and replacement constructions as available under hypotheses;
  \item set algebra and extensionality;
  \item Kuratowski ordered pairs;
  \item a natural-number object and finite recursion;
  \item formulas and lemmas used by bounded-satisfaction developments.
\end{itemize}

The goal is not to reproduce a full general-purpose set-theory ecosystem.  It
is a compact formal base with explicit syntax and semantics.  Downstream code
can quote its formulas and reason about their complexity without crossing a
translation from another language framework.

\subsection{Why one substantial module can be reasonable}

A single large \code{Zf.lean} keeps mutually recursive internal definitions and
their notation close, but it raises navigation cost.  The facade/audit pair and
section headings inside the file therefore matter.  When debugging a
downstream theorem, search for the declaration named in the import error or
audit rather than reading the module linearly.

The Rocq side mirrors the mathematical base with its own first-order and set
representations.  The pair is a contract-level parallel development, not a
shared extracted kernel object.

\section{ClosureAxiomatization}
\label{sec:closure-axiomatization}

\subsection{Alternative axiom}

The project studies a theory retaining Extensionality, Regularity, Separation,
and Powerset while replacing Pairing, Union, Infinity, and Replacement with a
single closure schema for set-like relations.  It proves equivalence with ZF.

The module organization separates two levels of equivalence.

\begin{description}
  \item[Shallow semantic directions] \code{Forward} and \code{Reverse}
        construct the missing axioms from the closure schema and the closure
        schema from ZF inside models.
  \item[Deductive equivalence] \code{Equivalence} combines the local
        first-order calculus, translations, and model theory to state the
        theory-level equivalence in the intended formal sense.
\end{description}

This distinction is important.  Showing that every model of one axiom list
satisfies each axiom of the other is the mathematical core, but a reusable
theory equivalence may also need syntax-level consequence and translation
lemmas.

\subsection{Build and exposition}

The Lean project is Mathlib-free, depending on the repository's own
first-order base.  A paired Rocq development independently checks the
directions.  The directory also includes an expository article in
\LaTeX/PDF, useful for matching theorem modules to the mathematical narrative.

\section{Bounded ZFC consistency}
\label{sec:bounded-zfc}

\subsection{The exact theorem}

For each external natural number \(n\), the project proves in ZFC the
consistency of ZFC proofs all of whose formula occurrences have quantifier rank
at most \(n\).  Schematically:

\[
  \text{for every meta-level } n,\qquad
  \mathrm{ZFC}\vdash \mathrm{Con}_n(\mathrm{ZFC}).
\]

This is a \emph{numeralwise scheme}.  It is not one object-language theorem
\(\mathrm{ZFC}\vdash\forall n\,\mathrm{Con}(n)\).  The distinction protects the
result from conflict with Gödel's second incompleteness theorem.

The rank used here counts every quantifier in the repository's language.  It
is finer than a Levy-hierarchy measure with primitive bounded quantifiers,
because bounded quantification must itself be encoded.

\subsection{Proof route}

The completed proof uses externally indexed partial satisfaction over the
repository's own first-order/ZF stack.  At a high level:

\begin{enumerate}
  \item define formula rank and the all-occurrences proof restriction;
  \item for a fixed external \(n\), define satisfaction for formulas up to the
        required rank;
  \item prove Tarski-style recursion equations for that finite rank;
  \item prove soundness of every inference rule for restricted proof trees;
  \item show the contradiction formula is false;
  \item arithmetize/quote the resulting argument as the ZFC derivation of the
        bounded-consistency sentence.
\end{enumerate}

The external index permits the meta-theory to build one finite truth
definition for each \(n\).  No uniform truth predicate for all formulas is
claimed.

\subsection{All-occurrences restriction}

The proof predicate checks every formula in the derivation, not merely the
conclusion.  This is stronger and more stable under the soundness induction:
each inference rule receives truth for exactly the formulas it mentions.
A conclusion-only restriction could admit arbitrarily complex intermediate
lemmas and would not support the same partial-satisfaction proof.

\subsection{Discarded routes documented}

The README records two important corrections:

\begin{itemize}
  \item an earlier empty-context sentence did not express the intended theorem
        in the right theory-relative way;
  \item a reflection-style route was abandoned in favor of direct externally
        indexed partial satisfaction.
\end{itemize}

Documenting wrong turns is particularly useful in metamathematics, where a
stronger-looking statement can be either false, inconsistent with known
theorems, or merely different because a quantifier moved from meta-language to
object language.

\subsection{Relation to ZFCinPA}

\project{BoundedZFCConsistency} supplies the complete numeralwise theorem.
\project{ZFCinPA}, located under Logic/Interpretability, tries to internalize
and uniformly prove the corresponding family in PA.  The latter must handle
nonstandard model elements and therefore has a genuinely stronger bridge
obligation.  The two projects should never be described as ``the same theorem
in different theories.''

\section{Set-theory dependency flow}

\begin{figure}[htbp]
\centering
\begin{tikzpicture}[
  b/.style={draw=blue!45,fill=sky,rounded corners=1mm,align=center,
            minimum width=34mm,minimum height=12mm,font=\sffamily\small},
  e/.style={draw=teal!50,fill=mint,rounded corners=1mm,align=center,
            minimum width=34mm,minimum height=12mm,font=\sffamily\small},
  c/.style={draw=gold!50,fill=sand,rounded corners=1mm,align=center,
            minimum width=34mm,minimum height=12mm,font=\sffamily\small},
  arrow/.style={-{Latex[length=2mm]},thick,draw=muted},
  node distance=7mm
]
\node[b] (fo) {FirstOrder\\syntax/calculus};
\node[b,right=of fo] (zf) {ZF formulas, models,\\internal constructions};
\node[e,right=of zf] (bzc) {BoundedZFCConsistency\\numeralwise scheme};
\node[c,below=8mm of bzc] (zpa) {ZFCinPA\\uniform conditional endpoint};
\node[e,below=8mm of zf] (closure) {ClosureAxiomatization\\ZF equivalence};
\draw[arrow] (fo) -- (zf);
\draw[arrow] (zf) -- (bzc);
\draw[arrow] (bzc) -- (zpa);
\draw[arrow] (zf) -- (closure);
\end{tikzpicture}
\caption{The repository-local set-theory stack and its main consumers.}
\label{fig:set-flow}
\end{figure}

\section{Recommended reading order}

\begin{readingbox}
\begin{enumerate}
  \item Read the FirstOrder language and relativization overview first.
  \item Read the ZF facade and audit, then search \code{Zf.lean} by public
        declaration rather than linearly.
  \item For bounded consistency, write the exact scheme and all-occurrences
        predicate on paper before following the partial truth recursion.
  \item Read ClosureAxiomatization's shallow semantic directions before its
        deductive-equivalence wrapper.
  \item Only then compare BoundedZFCConsistency with ZFCinPA; the
        standard/nonstandard index gap will be visible.
\end{enumerate}
\end{readingbox}

\begin{sourcebox}
\RepoFile{SetTheory/README.md};
\RepoFile{SetTheory/ZF/README.md};
\RepoDir{SetTheory/ZF/Lean/ZF};
\RepoFile{SetTheory/ClosureAxiomatization/README.md};
\RepoFile{SetTheory/BoundedConsistency/README.md};
\RepoFile{Logic/Interpretability/ZFCinPA/README.md}.
\end{sourcebox}

\chapter{Ancillary directories, vendored code, and tooling}
\label{chap:ancillary}

\section{Why these directories matter}

A walkthrough that covers only theorem modules misses several reproducibility
boundaries.  ProveIt has explicit homes for vendored code, proof-assistant
compatibility work, auxiliary sequence material, generated articles, and
redirect-only projects.  Their location communicates ownership and trust.

\section{\code{lib}: the only vendored-code home}

\RepoFile{lib/README.md} is the provenance index for third-party proof code.
The directory contains both git submodules and selected copied snapshots.

The policy is simple:

\begin{itemize}
  \item unmodified or locally hardened upstream source lives under
        \RepoDir{lib};
  \item repository-authored mathematical wrappers live under the owning
        subject project;
  \item compatibility patches that must remain reviewable live outside the
        submodule or are documented as local changes;
  \item licenses and upstream commits are recorded near the vendored tree.
\end{itemize}

This avoids two common monorepo failures: silently editing a dependency inside
a subject directory, and placing original project code in a subtree users
expect to be replaceable by a fresh vendor import.

\subsection{Coq Library of Undecidability}

The pinned/forked Undecidability library supplies standard machine,
lambda-calculus, and computability infrastructure.  CombinatoryLogic uses it
for external computability equivalences, but keeps local operational semantics
separate.

\subsection{FormalizedFormalLogic/Foundation}

Foundation is a read-only Lean source dependency and the authoritative module
inventory for FoundationPort.  The root Lake workspace registers it as a
library; the Rocq port does not modify it.

\subsection{MathComp-Abel}

MathComp-Abel provides the Galois-theory machinery used by Rocq
PolynomialFormulas.  Root Rocq setup installs its opam dependencies and maps
the theory root to logical namespace \code{Abel}.

\subsection{Synthetic Computability}

The Turing-degrees Rocq development wraps a pinned Synthetic Computability
source.  A project-local build script applies or stages compatibility work
needed by the current Rocq environment without turning the submodule into an
opaque mutable dependency.

\subsection{Coq-BB5}

Selected Busy Beaver certificates are copied with upstream provenance and
kernel-hardening notes.  Their status is therefore different from a git
submodule but still explicitly third-party.  Repository-authored semantics and
bridges remain outside \code{lib}.

\section{Tools/Rocq92Compat}
\label{sec:rocq-compat}

This directory is infrastructure, not a theorem project.  Its README documents
a reproducible compatibility overlay for a Windows/Rocq 9.2 environment where
package splits and upstream build assumptions do not line up cleanly.

The overlay includes several kinds of adaptation:

\begin{itemize}
  \item empty old-package-name adapters where package names changed;
  \item findlib aliases;
  \item pinned Elpi and MathComp repairs;
  \item proof-script ordering patches for Coquelicot and Interval;
  \item Windows launchers and command shims;
  \item MinGW/remake compatibility fixes.
\end{itemize}

A compatibility overlay carries its own trust and maintenance concerns.  It
can change which library source is compiled or how commands are resolved, but
it is not a mathematical axiom.  The important review properties are:
patches are version-pinned, diffable, reproducible, and kept separate from
theorem files.

\begin{gotchabox}
Do not copy a compatibility-patched dependency tree back into \code{lib} and
commit it as though it were the upstream snapshot.  Keep the immutable source,
patch set, and generated build tree distinguishable; otherwise provenance and
regeneration become impossible to audit.
\end{gotchabox}

\section{Oeis}

\RepoDir{Oeis} contains auxiliary or legacy sequence material, notably an
A290268 surface related to the combinatorics derivative-expansion project.
New sequence-specific proof code generally lives in the relevant mathematical
domain, while the OEIS directory provides sequence-oriented artifacts or
compatibility paths.

The A198683 development demonstrates a stronger sequence-status convention:
only values backed by the current accepted certificate surface are exposed as
accepted data, even if exploratory and partial proof modules for later values
already exist.

\section{Shenanigans}

The top-level \RepoDir{Shenanigans} directory is a redirect to a separate
repository.  It has no active proof architecture to include in the root build.
Treating the redirect as documentation avoids false expectations that missing
Lean/Rocq modules were omitted from this walkthrough.

\section{Articles and reports}

Several project directories contain \LaTeX{} sources and compiled PDFs, for
example the Lazard formalization report and the closure-axiomatization article.
These documents are not theorem dependencies unless a build file explicitly
says otherwise.  They serve as:

\begin{itemize}
  \item mathematical narrative around a module graph;
  \item claim-to-module crosswalks;
  \item status snapshots for long-running proof programs;
  \item source/provenance records for finite data.
\end{itemize}

A regenerated report can therefore be a useful project milestone without
changing a kernel theorem.  Conversely, a green article build does not imply
a green Lean/Rocq build.

\section{Repository hygiene files}

\RepoFile{.gitattributes} normalizes text to LF while preserving binary
treatment for PDFs, images, and archives.  This is particularly important for
generated source: newline churn can turn a deterministic regeneration into a
repository-wide diff.

\RepoFile{.editorconfig} and \RepoFile{.gitignore} establish ordinary editor
and build-output conventions.  The pinned commit has no repository-level
\code{.github/workflows} tree, so the source does not expose a single
GitHub-Actions validation matrix.  Build reproducibility is consequently
documented primarily in root and project READMEs.

\section{Ownership map}

\begin{longtable}{L{.30\textwidth}L{.28\textwidth}L{.31\textwidth}}
\caption{Where an artifact should live.}
\label{tab:ownership-map}\\
\toprule
Artifact & Preferred home & Reason \\
\midrule
\endfirsthead
\toprule
Artifact & Preferred home & Reason \\
\midrule
\endhead
Hand-written theorem or bridge &
Owning subject project &
Keeps mathematical dependencies local.\\
Generated formal data/certificate &
Owning subject project's Lean/Rocq tree &
It is part of the checked import graph.\\
Generator/search/replay program &
Project \code{Support}, \code{Tools}, or \code{Research} subtree &
Separates candidate production from proof terms.\\
Unmodified upstream source &
\code{lib} submodule or vendored snapshot &
Makes provenance and replacement obvious.\\
Compatibility patch/wrapper &
Owning project or \code{Tools/Rocq92Compat} &
Avoids mutating the pinned source invisibly.\\
Expository article/crosswalk &
Project root or documentation subtree &
Keeps narrative adjacent without making it an import.\\
Coverage ledger &
Port project root &
Status is tied to the source inventory, not one target module.\\
\bottomrule
\end{longtable}

\begin{sourcebox}
\RepoFile{lib/README.md};
\RepoFile{.gitmodules};
\RepoFile{Tools/Rocq92Compat/README.md};
\RepoDir{Oeis};
\RepoDir{Shenanigans};
\RepoFile{.gitattributes}.
\end{sourcebox}

\part{Building, debugging, and extending ProveIt}

\chapter{Working with the repository}
\label{chap:working}

\section{Start from a pinned tree}

For reproducible work, record the commit before running a generator or changing
a theorem.  The repository contains submodules, generated modules, and
fast-moving status documents; ``current main'' is not a durable input.

\begin{lstlisting}[language=PowerShell,caption={Pin and initialize a checkout.}]
git clone https://github.com/VladimirReshetnikov/ProveIt.git
Set-Location ProveIt
git checkout 675d0a2ff0fc624fba199e006049016ba806f40a
git submodule update --init --recursive
git status --short
git submodule status
\end{lstlisting}

For a narrow Lean-only project, initializing every submodule may be
unnecessary.  For root Rocq work or projects depending on Foundation,
MathComp-Abel, Undecidability, or Synthetic Computability, initialize the
specific listed submodules before interpreting missing imports as source
errors.

\section{Lean environment}

\subsection{Toolchain and cache}

The root \code{lean-toolchain} selects Lean 4.32.0, and
\code{lake-manifest.json} pins Mathlib and transitive dependencies.  With
Elan and Lake available:

\begin{lstlisting}[language=PowerShell]
lean --version
lake --version
lake exe cache get
\end{lstlisting}

The Mathlib cache removes a large amount of baseline compilation.  It does not
contain repository-local generated certificates, so a first broad build can
still be substantial.

\subsection{Broad validation}

\begin{lstlisting}[language=PowerShell]
lake build
\end{lstlisting}

This requests the root default targets.  It is useful after focused work, but
is a poor first diagnostic because a failure may occur far from the changed
module and multiple large jobs can compete for memory.

To validate only the consumer facade:

\begin{lstlisting}[language=PowerShell]
lake build ProveIt
\end{lstlisting}

To validate an explicitly registered library or audit, use the exact target
spelling in \code{lakefile.toml} or the project's README.

\subsection{Direct module elaboration}

For a source file that is not a registered target or not reachable from a
facade, use:

\begin{lstlisting}[language=PowerShell]
lake env lean path/to/Module.lean
\end{lstlisting}

This is the most reliable antidote to a silent target no-op.  It elaborates the
literal file under the root package environment and returns an error if the
file is absent or invalid.

A project-local Lake environment may be preferable:

\begin{lstlisting}[language=PowerShell]
lake --dir Logic/QuantifierCommutation/Lean build
lake --dir Logic/FirstOrder/Lean build
\end{lstlisting}

The local manifest can intentionally omit Mathlib and provide a much smaller
dependency graph.

\subsection{Threading and memory}

Use the project README's worker recommendation.  Typical patterns are:

\begin{lstlisting}[language=PowerShell]
# Serialize heavy elaboration jobs.
$env:LEAN_NUM_THREADS = "0"
lake build SomeHeavyCertificate

# Restore ordinary worker behavior for projects that need it.
Remove-Item Env:LEAN_NUM_THREADS -ErrorAction SilentlyContinue
\end{lstlisting}

Lean's \code{maxHeartbeats} limits elaboration work inside one declaration;
Lake worker count limits concurrent processes.  Raising the former does not
solve system-wide memory pressure, and serializing jobs does not fix a single
nonterminating reduction.

\section{Focused Lean validation recipe}

For an ordinary project:

\begin{enumerate}
  \item elaborate the changed leaf file directly;
  \item build its facade;
  \item build its audit;
  \item build the root library target if the project is registered there;
  \item finish with \code{lake build ProveIt} or the root default set as
        appropriate.
\end{enumerate}

For a generated certificate project:

\begin{enumerate}
  \item regenerate into a clean temporary directory;
  \item compare generated paths and hashes with the checked-in tree;
  \item elaborate the smallest changed \code{Data} module;
  \item elaborate its immediate \code{Certificate};
  \item walk up term/tail/shard parents;
  \item build the semantic bridge;
  \item build facade and audit.
\end{enumerate}

This order localizes both generation and reduction failures.

\section{Rocq environment}

\subsection{Version and package roots}

The root documentation targets Rocq 9.2 or later for current builds, while
some MetaRocq exact-equivalence work is tied to a 9.0-era environment.  Verify
the executable names supplied by the installation:

\begin{lstlisting}[language=PowerShell]
rocq --version
coqc --version
opam switch show
opam list
\end{lstlisting}

Do not assume \code{coqc}, \code{rocq c}, and \code{rocq compile} are
interchangeable in scripts across installations.  Preserve project commands
unless deliberately modernizing them.

\subsection{Root build}

The root build requires its submodule dependencies and project-specific
Synthetic Computability staging:

\begin{lstlisting}[language=PowerShell]
git submodule update --init `
  lib/Coq-Synthetic-Computability `
  lib/MathComp-Abel

opam install --yes --deps-only `
  ./lib/MathComp-Abel/coq-mathcomp-abel.opam

pwsh -NoProfile -File `
  Computability/TuringDegrees/Coq/BuildSyntheticComputability.ps1

rocq makefile -f _CoqProject -o Makefile.coq
make -f Makefile.coq
\end{lstlisting}

On Windows, consult \RepoFile{Tools/Rocq92Compat/README.md} before patching
dependencies ad hoc.

\subsection{Focused sequential build}

A dependency-light project often documents explicit commands:

\begin{lstlisting}[language=PowerShell]
coqc -Q Logic/QuantifierCommutation/Coq QuantifierCommutation `
  Logic/QuantifierCommutation/Coq/Commutation.v

coqc -Q Logic/QuantifierCommutation/Coq QuantifierCommutation `
  Logic/QuantifierCommutation/Coq/Counterexamples.v

coqc -Q Logic/QuantifierCommutation/Coq QuantifierCommutation `
  Logic/QuantifierCommutation/Coq/Audit.v

coqchk -silent `
  -Q Logic/QuantifierCommutation/Coq QuantifierCommutation `
  QuantifierCommutation.Commutation `
  QuantifierCommutation.Counterexamples `
  QuantifierCommutation.Audit
\end{lstlisting}

This style makes logical roots and module order explicit.  For a large
dependency graph, generate a project Makefile and request the endpoint
\code{.vo} or audit target rather than compiling the root universe.

\section{Audit workflow}

\subsection{Lean}

After the theorem elaborates, run the audit and read the output, not just the
exit code.  A new axiom can appear because an imported theorem changed even
when the local proof still checks.

A useful temporary scratch module is:

\begin{lstlisting}[language=Lean]
import Project.Audit

#check Project.headlineTheorem
#print Project.headlineTheorem
#print axioms Project.headlineTheorem
\end{lstlisting}

For a polymorphic theorem, inspect the fully elaborated type.  If implicit
parameters are suspicious, temporarily enable detailed pretty printing:

\begin{lstlisting}[language=Lean]
set_option pp.all true in
#check Project.headlineTheorem
\end{lstlisting}

\subsection{Rocq}

Use both source-level and kernel-level checks where documented:

\begin{lstlisting}[language=Rocq]
Print Assumptions HeadlineTheorem.
Check HeadlineTheorem.
\end{lstlisting}

Then run \code{coqchk}/\code{rocqchk} on the public module set.  This catches
damaged or inconsistent compiled objects and independently validates the
environment assumptions accepted by the kernel checker.

\section{Generator workflow}

A generator invocation should be reproducible from a clean checkout.  Record:

\begin{itemize}
  \item exact interpreter/tool version;
  \item command line and working directory;
  \item input files and their hashes;
  \item deterministic seed or proof that no randomness is used;
  \item expected output paths;
  \item formatting/canonicalization version;
  \item the first formal certificate target consuming the output.
\end{itemize}

Prefer generation into a temporary tree followed by a directory diff:

\begin{lstlisting}[language=PowerShell]
$Out = Join-Path $env:TEMP "proveit-generated"
Remove-Item $Out -Recurse -Force -ErrorAction SilentlyContinue
New-Item $Out -ItemType Directory | Out-Null

python path/to/generator.py --output-root $Out

git diff --no-index -- `
  path/to/checked-in/generated `
  $Out
\end{lstlisting}

If the generator writes directly into the source tree, first ensure
\code{git status --short} is empty and inspect every changed path before
building.

\section{Reading build output correctly}

A zero exit code can mean:

\begin{itemize}
  \item the intended target built;
  \item the target was already up to date;
  \item the command selected no target;
  \item a wrapper skipped optional work;
  \item a noisy converter returned success after producing a partial artifact.
\end{itemize}

For Lean, verify that the requested module is registered/reachable and inspect
the \code{.olean} timestamp only as a secondary clue.  For Rocq, verify the
requested logical module appears in the generated Makefile and that its
\code{.vo} is refreshed.  For generated certificates, verify the endpoint
imports the changed shard.

\section{A clean validation record}

For a substantial theorem or generated checkpoint, preserve a text record:

\begin{lstlisting}[language=PowerShell]
$ErrorActionPreference = "Stop"
$Log = "validation-675d0a2f.txt"

"Commit: $(git rev-parse HEAD)" | Tee-Object $Log
"Lean:   $(lean --version)"     | Tee-Object $Log -Append
"Lake:   $(lake --version)"     | Tee-Object $Log -Append
"Rocq:   $(rocq --version)"     | Tee-Object $Log -Append
git submodule status            | Tee-Object $Log -Append

lake build Project Project.Audit 2>&1 |
  Tee-Object $Log -Append
\end{lstlisting}

The record is not a proof artifact, but it makes an engineering claim such as
``the focused checkpoint was green on this toolchain'' reproducible.

\begin{sourcebox}
\RepoFile{README.md};
\RepoFile{lean-toolchain};
\RepoFile{lakefile.toml};
\RepoFile{lake-manifest.json};
\RepoFile{_CoqProject};
\RepoFile{Tools/Rocq92Compat/README.md};
\RepoFile{Computability/TuringDegrees/Coq/BuildSyntheticComputability.ps1}.
\end{sourcebox}

\chapter{Debugging proof and build failures}
\label{chap:debugging}

\section{Classify the failure before changing the proof}

Most wasted time in a formal monorepo comes from treating every red build as a
mathematical failure.  ProveIt failures fall into several distinct classes.

\begin{longtable}{L{.23\textwidth}L{.31\textwidth}L{.35\textwidth}}
\caption{Failure classes and first diagnostics.}
\label{tab:failure-classes}\\
\toprule
Class & Symptom & First action \\
\midrule
\endfirsthead
\toprule
Class & Symptom & First action \\
\midrule
\endhead
Target discovery &
Command succeeds instantly or reports nothing; expected module never appears. &
Elaborate the literal file with \code{lake env lean}; inspect facade and
\code{lakefile.toml}.\\
Import resolution &
Unknown module, logical path, or package. &
Check root/local workspace, submodule initialization, and \code{-Q}/\code{-R}
mapping.\\
Elaboration/type inference &
Type mismatch, failed synthesis, unsolved metavariable, unexpected implicit. &
Minimize to the leaf declaration; enable \code{pp.all}; disable
\code{autoImplicit}.\\
Tactic proof search &
Tactic timeout, no applicable rule, normalization blow-up. &
Replace the endpoint tactic with intermediate lemmas; inspect the goal after
each transformation.\\
Kernel reduction &
\code{decide} exceeds heartbeats/recursion/memory. &
Shrink or re-shard data; expose definitional structure; avoid global option
increases first.\\
Generated dependency &
Missing shard, cyclic import, stale name, wrong tail boundary. &
Compare generator manifest/import graph; build immediate data and certificate
modules.\\
Semantic bridge &
Certificate equality checks but endpoint does not follow. &
Inspect the checker-correctness or evaluation theorem, not the data.\\
Rocq environment &
Package or plugin name mismatch, Windows launcher failure. &
Reproduce the documented compatibility overlay and exact opam switch.\\
Cross-prover parity &
One assistant proves a theorem with extra premises or different encoding. &
Compare literal theorem types and coverage ledger; do not patch names to match.\\
Specification &
The theorem is easy for a suspicious reason or audits show extra parameters. &
Print the fully elaborated statement; trace object/meta/model quantifiers.\\
\bottomrule
\end{longtable}

\section{Target-discovery failures}

The strongest symptom is a ``successful'' build that touched nothing.  This
can happen when a module exists under a registered source root but is not
reachable from a target, or when a command syntax names a target Lake does not
know.

Use three tests:

\begin{enumerate}
  \item Does the file exist at the expected physical path?
  \item Can \code{lake env lean path/to/file.lean} elaborate it?
  \item Does a registered facade import it?
\end{enumerate}

If only test 2 succeeds, the source is valid but integration is missing.
Choose deliberately whether to add it to a facade, create a target, or leave it
as direct-build research material.

For Rocq, inspect the \code{_CoqProject} source list and logical mapping.  A
\code{.v} file outside the list can compile manually yet be absent from the
root Makefile.

\section{Import-resolution failures}

\subsection{Lean logical name versus physical path}

An import such as

\begin{lstlisting}[language=Lean]
import BoundedPAConsistency.Audit
\end{lstlisting}

does not imply a directory named \code{BoundedPAConsistency} at the root.
The Lake registry maps the logical prefix to a deeper source directory.
When moving modules, update the source-root mapping or preserve the logical
prefix; do not make import paths mirror physical taxonomy accidentally.

If the same project has a local Lake file, verify which environment the editor
and command line are using.  A module available in the root workspace may be
unknown in the local project, or vice versa.

\subsection{Rocq logical roots}

For Rocq, a file path and module name are connected by \code{-Q} or \code{-R}.
Changing one side without the other can create duplicate logical names,
unbound imports, or \code{coqchk} failures.  Use the exact mapping from the
project README in focused commands.

\subsection{Submodule state}

A missing Foundation, Abel, Undecidability, or Synthetic Computability import
can simply mean the gitlink was not initialized.  Check
\code{git submodule status}; a leading minus indicates an uninitialized
submodule, while a plus indicates a checked-out commit different from the
superproject's pin.

\section{Elaboration and specification failures}

\subsection{Unexpected implicit parameters}

Lean's automatic implicit-variable creation is convenient in ordinary
mathematics and dangerous in statement-heavy metatheory.  If a theorem
unexpectedly generalizes over a model, code, or translation:

\begin{enumerate}
  \item print the theorem with \code{pp.all};
  \item place \code{set_option autoImplicit false} near the declarations;
  \item declare every universe, type, model, and coding parameter explicitly;
  \item add an audit \code{\#check} with the intended type.
\end{enumerate}

A theorem that proves too easily after a refactor should be treated as a
specification regression until its fully elaborated type is reviewed.

\subsection{Coercion and notation drift}

Large projects cross \(\mathbb N\), \(\mathbb Z\), \(\mathbb Q\),
\(\mathbb R\), polynomial constants, finite indices, and quoted object-language
numerals.  Tactic failures often come from a changed coercion rather than a
false lemma.

Normalize one boundary at a time:

\begin{itemize}
  \item expose \code{Nat.cast}, \code{Int.cast}, or polynomial \code{C};
  \item use explicit type annotations on numerals;
  \item prove a local cast lemma instead of adding more simplifier lemmas
        globally;
  \item distinguish an object-language numeral constructor from the
        meta-language number.
\end{itemize}

\subsection{Quotients and representative changes}

FirstOrder term models and Turing degrees use quotient constructions.  A
definition over representatives must prove compatibility before descending.
When a quotient theorem fails after an upstream API change, look for:

\begin{itemize}
  \item a missing congruence theorem;
  \item a representative-level equality weakened to equivalence;
  \item an accidental use of definitional equality where only quotient
        equality is available;
  \item changed simplification around \code{Quot.sound}.
\end{itemize}

\section{Tactic and normalization failures}

\subsection{Do not respond first by raising all limits}

A timeout can signal that the proof lost structure.  Before increasing
heartbeats:

\begin{enumerate}
  \item identify the last fast intermediate state;
  \item split a global \code{simp}/\code{ring}/\code{norm_num} into semantic
        rewrites followed by algebraic normalization;
  \item make finite indices explicit with \code{fin_cases};
  \item isolate one polynomial coordinate or one matrix block;
  \item mark a large reusable equality as a theorem and rewrite by it;
  \item check whether a reducible definition was made opaque or vice versa.
\end{enumerate}

The Jacobian determinant proof is a good model: reduce the general determinant
to a \(3\times3\) closed formula, unfold derivatives, then call \code{ring}.

\subsection{Simplifier loops and giant goals}

Generated algebra code can make \code{simp} expand a large table repeatedly.
Use restricted simplification sets, unfold only the top-level function, and
prefer a named semantic lemma over exposing every definition.

In a certificate module, \code{change} is often intentional: it reveals a
specific list-addition spine so previously proved term certificates can
rewrite it.  Replacing that with indiscriminate unfolding can create an
unmanageable expression.

\section{Kernel-certificate failures}

\subsection{A bisect strategy for term/tail shards}

Suppose \code{table0_tail22_certificate} no longer closes.

\begin{enumerate}
  \item Build/import \code{Tail22Data} alone.  A failure here is generated
        syntax, type, or dependency.
  \item Check the immediate \code{merge_certificate}.  If \code{decide} fails,
        the candidate and normal literal differ.
  \item Check each term certificate 88--91 and tail 23 independently.
  \item Expand only the top-level tail definition and verify the expected
        association of \code{add}.
  \item Compare generator output for the first mismatching term, not the entire
        normal polynomial.
\end{enumerate}

A binary-search or Merkle-style manifest over generated pieces can shorten
this process further; see the recommendations in \cref{chap:assessment}.

\subsection{Recursion depth versus heartbeats}

\code{maxRecDepth} failures indicate deeply nested term reduction or elaborator
recursion.  \code{maxHeartbeats} failures indicate too much total work.
Increasing the wrong option merely changes the error.  Long list literals and
right-associated additions tend to stress depth; large normalization across
many monomials stresses work and memory.

\subsection{Native replay disagrees with kernel replay}

If an external or native checker says a certificate is valid but kernel
\code{decide} does not:

\begin{itemize}
  \item confirm both evaluate the same serialized data and normalization
        convention;
  \item check integer overflow or machine-word arithmetic in the fast tool;
  \item ensure sorting and duplicate-combination rules match;
  \item verify signs and exponent tuple order;
  \item reduce the disagreement to the smallest term and compare exact output.
\end{itemize}

The kernel result is authoritative for the kernel theorem.  The fast replay is
a diagnostic oracle, not a reason to replace the theorem with an assertion.

\section{Semantic-bridge failures}

A certificate can be perfectly valid for the wrong predicate.  Typical bridge
bugs include:

\begin{itemize}
  \item a coverage checker validates every unit cell but the real-lifting
        theorem assumes boundary behavior not proved;
  \item a machine trace checker proves halting but the endpoint needs maximal
        score among all machines;
  \item a sparse polynomial list equality is proved but evaluation preservation
        under normalization is missing;
  \item a coded proof verifier checks conclusions but the partial-truth
        induction requires all formula occurrences bounded;
  \item a compiler simulates successful evaluation but the universality
        theorem needs divergence preservation.
\end{itemize}

Debug by writing the exact implication:
\[
  \operatorname{check}(d)=\mathrm{true}
  \quad\Longrightarrow\quad
  \text{headline mathematical property}.
\]
Then inspect every premise of the theorem implementing that arrow.

\section{Cross-prover mismatches}

When Lean and Rocq theorem names align but statements do not, avoid forcing a
syntactic port.  Compare:

\begin{enumerate}
  \item universe/type assumptions;
  \item constructive versus classical principles;
  \item quotient versus setoid equality;
  \item encoding and indexing conventions;
  \item standard versus nonstandard quantification;
  \item total versus partial computation relations;
  \item imported theorem provenance.
\end{enumerate}

Update the coverage ledger with the semantic relationship.  A more general
Rocq theorem or a more executable Lean theorem can both be correct without one
being ``behind.''

\section{Rocq environment and compatibility failures}

The compatibility overlay should be reproduced before editing proof scripts.
Package-name adapters and proof-order patches can make the same source parse or
resolve differently.

Useful diagnostics include:

\begin{lstlisting}[language=PowerShell]
where.exe coqc
where.exe rocq
opam var lib
ocamlfind list
rocq makefile -f _CoqProject -o Makefile.coq
make -f Makefile.coq -n Target.vo
\end{lstlisting}

The dry run reveals which compiler command and logical paths the Makefile will
use.  If a patched library is involved, record its exact source commit and
patch hash.

\section{Minimal reproductions inside the monorepo}

A good minimal reproduction remains in the same package environment but
imports as little as possible:

\begin{lstlisting}[language=Lean]
import Project.SmallestDependency

namespace Scratch

-- Copy only the failing statement and the few definitions it needs.

end Scratch
\end{lstlisting}

Place it temporarily beside the project or under a scratch path excluded from
the final commit.  Do not copy giant generated definitions manually; import
the data module and isolate one theorem.

For Rocq, create a temporary \code{Scratch.v} with the same \code{-Q}/\code{-R}
flags and only the failing \code{Require Import}s.

\section{When to suspect the statement}

The following are red flags:

\begin{itemize}
  \item a difficult theorem becomes trivial after moving a variable;
  \item a bound changes from external \(n\) to internal \(\forall n\) without a
        new uniformity argument;
  \item a theorem over standard syntax is applied to arbitrary model elements;
  \item the audit prints an unexpected parameter or no longer names a crucial
        premise;
  \item a translation theorem loses its satisfaction-preservation hypothesis;
  \item a final theorem imports a \code{Placeholder}, \code{Assumption}, or
        \code{Conditional} module not mentioned in status docs.
\end{itemize}

Stop optimizing the proof and review the specification.  In formal
metamathematics, the wrong theorem can be easier than the right one.

\begin{sourcebox}
\RepoFile{Logic/Interpretability/ZFCinPA/README.md};
\RepoFile{Logic/PeanoArithmetic/BoundedConsistency/README.md};
\RepoFile{Algebra/PolynomialFormulas/Lean/PolynomialFormulas/ComputableDummitKernelCertificateTable0Tail22Certificate.lean};
\RepoFile{Tools/Rocq92Compat/README.md};
\RepoFile{Computability/TuringDegrees/COVERAGE.md}.
\end{sourcebox}

\chapter{Extending a project without blurring its boundary}
\label{chap:extension}

\section{Choose the project shape first}

Before adding files, decide which of the repository's project species the work
belongs to:

\begin{enumerate}
  \item a direct theorem over an existing semantic API;
  \item a new semantic layer or calculus;
  \item a finite certificate for an existing checker;
  \item a new checker plus generated certificates;
  \item an independent Rocq/Lean port;
  \item a status/coverage extension without a new endpoint.
\end{enumerate}

This choice determines where definitions live, what the facade should import,
which audits are needed, and whether support tooling is inside or outside the
proof boundary.

\section{Adding a direct Lean theorem}

Suppose an existing project already has the right definitions.

\subsection{Place the theorem at the lowest reusable layer}

If the theorem is a general property of the core datatype, place it near the
core or in a focused lemma module.  If it assembles project-specific results,
place it in an endpoint module.  Avoid adding a deep theorem directly to a
facade.

A typical sequence is:

\begin{lstlisting}[language=Lean]
-- Project/NewProperty.lean
import Project.Core
import Project.ExistingBridge

namespace Project

theorem newProperty ... : ... := by
  ...

end Project
\end{lstlisting}

Then import \code{Project.NewProperty} from the facade only if the theorem is
part of the consumer surface.

\subsection{Generalize only where it lowers dependencies}

Prove determinant identities over a commutative ring and add
characteristic-zero only at nonvanishing.  Prove finite combinatorial maps over
finite types and add concrete coordinates later.  But do not generalize a
statement into a forest of typeclasses that hides its intended semantics.

A useful test is whether the generalized lemma removes an import or serves at
least two endpoint paths.  Otherwise a concrete theorem can be clearer.

\subsection{Add an audit entry}

The audit should contain the literal public theorem:

\begin{lstlisting}[language=Lean]
#check Project.newProperty
#print axioms Project.newProperty
\end{lstlisting}

If a theorem is intended to remain constructive, add an explanatory comment
about the expected empty/classical assumption footprint.  If it deliberately
uses an imported classical theorem, name that provenance.

\section{Adding a semantic layer}

A new proof system, machine model, or object-language encoding needs more than
one endpoint theorem.

\subsection{Define the contract}

Write down:

\begin{itemize}
  \item syntax and well-formedness;
  \item semantics or observation relation;
  \item executable operations, if any;
  \item equivalence/congruence needed for quotients;
  \item public theorem relating syntax and semantics;
  \item exact imported assumptions.
\end{itemize}

Then split the files along those responsibilities.  The FirstOrder and
Presburger projects are good templates.

\subsection{Prove substitution/renaming early}

Many later metatheorems depend on substitution and assignment lemmas.  If these
are postponed, every semantic induction accumulates ad hoc rewrites.  Establish
identity, composition, and semantic substitution before soundness or
completeness.

\subsection{Keep data-carrying and Prop-valued relations distinct}

An inductive proof tree carrying premises is useful for executable verification
and all-occurrences bounds.  A Prop-valued derivability relation is convenient
for mathematical statements.  Provide an erasure/soundness bridge rather than
forcing one representation to do both jobs.

\section{Adding a finite certificate}

\subsection{Start with the mathematical predicate}

Define a checker against arbitrary candidate data:

\begin{lstlisting}[language=Lean]
def checkCertificate (d : CertificateData) : Bool := ...

theorem checkCertificate_sound {d : CertificateData} :
    checkCertificate d = true -> MathematicalProperty d := by
  ...
\end{lstlisting}

Only then introduce the concrete certificate.  If the soundness theorem is
proved after data generation, there is a temptation to validate exactly the
properties the candidate happens to satisfy rather than the properties the
endpoint needs.

\subsection{Choose an evaluator deliberately}

Use kernel \code{decide} for bounded computations that fit.  Use
\code{native_decide} only with a documented trust decision.  On Rocq, document
\code{vm_compute} or other reflection machinery.  Keep a slower small test
certificate that can be checked by the strongest evaluator for regression.

\subsection{Separate placement from meaning}

For geometry:

\[
\text{integer placement check}
\to \text{real-set coverage}
\to \text{dissection theorem}.
\]

For machines:

\[
\text{trace replay}
\to \text{operational halting}
\to \text{score bound/classification}.
\]

For algebra:

\[
\text{list equality}
\to \text{evaluation equality}
\to \text{invariant/resolvent theorem}.
\]

Each arrow deserves a named theorem.  This makes it impossible for a future
refactor to silently weaken the checker while leaving the endpoint name
unchanged.

\section{Adding or changing a generator}

\subsection{Treat output format as an API}

Generated Lean/Rocq source has consumers and compatibility constraints.  Define
a stable schema:

\begin{itemize}
  \item canonical ordering of terms;
  \item duplicate-combination and zero-removal rules;
  \item integer/sign formatting;
  \item shard-size and naming policy;
  \item import direction;
  \item deterministic header containing generator version and input hash.
\end{itemize}

Changing canonicalization may be mathematically harmless but can invalidate
every downstream merge certificate and create an enormous diff.

\subsection{Emit a manifest}

For each generated module, a manifest should record:

\begin{itemize}
  \item logical module name and file path;
  \item content hash;
  \item parent shards;
  \item semantic root or source term range;
  \item expected certificate module;
  \item resource options needed to elaborate it.
\end{itemize}

Even if the repository does not yet enforce such a manifest globally, adding
one to a new generator makes shard completeness and binary-search debugging
much easier.

\subsection{Regenerate all or prove locality}

A partial regeneration is safe only if the generator can prove that unchanged
shards do not depend on changed input.  Otherwise regenerate the full family
into a temporary directory and compare.  Check both content and import graph.

\subsection{Keep generator tests outside theorem claims}

Unit tests can verify parser round trips, canonical ordering, and deterministic
output.  They improve engineering reliability but do not substitute for
certificate theorems.  README language should say ``generator test'' or
``independent replay,'' not ``formal proof.''

\section{Adding a Rocq port}

\subsection{Choose the parity level}

Before coding, classify the goal:

\begin{itemize}
  \item exact statement-level port;
  \item stronger/more general target theorem;
  \item theorem-surface port with a different encoding;
  \item partial dependency layer;
  \item declaration-free facade;
  \item complementary result, not a port.
\end{itemize}

Record that classification in the coverage ledger.  A theorem with the same
English title but an extra choice principle or standardness premise is not an
exact port.

\subsection{Map dependencies, not just declarations}

For each source module, record:

\begin{itemize}
  \item source imports;
  \item target modules providing analogous concepts;
  \item representation translations;
  \item additional logical assumptions;
  \item public declarations covered;
  \item deliberate omissions.
\end{itemize}

A port can be blocked by a missing semantic layer even if the target theorem
itself looks short.  FoundationPort's conservative ledger is the model.

\subsection{Audit and kernel-check the target}

Add \code{Print Assumptions} for headline declarations and a focused
\code{coqchk}/\code{rocqchk} command.  If the project uses a vendored library,
state its commit and logical root.

\section{Registering a new project}

A new Lean project may need some or all of:

\begin{enumerate}
  \item mathematical directory under the correct top-level subject;
  \item \code{Lean/ProjectName.lean} facade;
  \item implementation namespace and modules;
  \item \code{Audit.lean};
  \item project README with status, layout, trust, and build;
  \item local \code{lakefile.toml} if independent/focused builds matter;
  \item root \code{lean_lib} mapping;
  \item optional root facade import;
  \item optional root default target.
\end{enumerate}

These last three choices are independent.  A heavyweight project can be
registered but omitted from \code{ProveIt.lean}; a critical audit can be a
default target without exposing all internals.

A new Rocq project needs a logical-root mapping, ordered source registration
or local \code{_CoqProject}, audit, and kernel-check command.

\section{Writing status documentation}

A strong project README should answer the following in explicit headings:

\begin{description}
  \item[Headline result] Literal mathematical claim and exact endpoint name.
  \item[Status] Unconditional, conditional, partial, or exploratory.
  \item[Architecture] Module path from definitions to endpoint.
  \item[Trust boundary] Kernel, evaluator, generated candidates, source
        transcription, imported theorems.
  \item[Lean/Rocq relationship] Independent, contract-parallel, complementary,
        or incomplete port.
  \item[Build] Focused commands and resource requirements.
  \item[Deliberate nonclaims] Nearby stronger statements not established.
  \item[Residue] Ordered remaining obligations for active projects.
\end{description}

Do not write ``formalized'' when only a definition or reduction is present.
Do not write ``verified by script'' when the script is outside the proof
assistant.  Do not write ``Lean and Coq proofs'' when one side is conditional
and the condition is the central theorem.

\section{Commit and review hygiene}

For hand-written changes:

\begin{itemize}
  \item keep refactors separate from theorem changes when possible;
  \item include the focused build/audit command in the commit message;
  \item state any changed theorem type explicitly;
  \item update the project README and coverage ledger in the same commit.
\end{itemize}

For generated changes:

\begin{itemize}
  \item separate generator-source changes from generated output only if the
        intermediate commits are clearly marked and not merged independently;
  \item record generator command and deterministic manifest;
  \item avoid hand-editing generated modules;
  \item include the first and final certificate targets validated;
  \item summarize shard frontier changes rather than listing thousands of
        filenames.
\end{itemize}

\section{Pre-merge checklist}

\begin{tcolorbox}[colback=lightgray,colframe=navy,title={Extension checklist}]
\begin{enumerate}
  \item The theorem statement has been reviewed with all implicit parameters
        visible.
  \item The module is reachable from the intended facade or has a documented
        direct-build path.
  \item The smallest leaf, endpoint, facade, and audit all build.
  \item Generated data are consumed by a proved checker.
  \item The evaluator and trust profile are documented.
  \item Lean/Rocq parity claims match literal theorem types.
  \item Conditional premises and open search claims remain visible.
  \item The README status and residue are current.
  \item Submodule pins and vendored provenance are unchanged or deliberately
        updated.
  \item A clean checkout can reproduce the focused validation.
\end{enumerate}
\end{tcolorbox}

\begin{sourcebox}
\RepoFile{ProveIt.lean};
\RepoFile{lakefile.toml};
\RepoFile{_CoqProject};
\RepoFile{Computability/TuringDegrees/COVERAGE.md};
\RepoFile{Logic/FoundationPort/README.md};
\RepoFile{Algebra/PolynomialFormulas/Lean/PolynomialFormulas/ComputableDummitCoefficientsCore.lean};
\RepoFile{Combinatorics/SquaredSquare/README.md}.
\end{sourcebox}

\chapter{Status and active frontiers at the pinned commit}
\label{chap:frontiers}

\section{Three different meanings of ``unfinished''}

A large formal repository can have an unfinished component for three
independent reasons.

\begin{description}
  \item[Mathematical frontier] The desired statement is open or a historical
        exhaustive argument has not been formalized.
  \item[Formalization frontier] The mathematics is known or conjectured, but a
        bridge, translation, or certificate layer is missing.
  \item[Build/checkpoint frontier] The architecture is established and shards
        are being generated or kernel-checked incrementally.
\end{description}

These should not be collapsed into one progress percentage.  A completed
conditional reduction can be mathematically valuable while its premise remains
open.  A generated module can be physically present while its parent endpoint
has not yet imported it.  An entire theorem can be complete in Lean but only
partially ported to Rocq.

\section{Snapshot catalogue}

\begin{landscape}
\begin{longtable}{L{.20\linewidth}L{.28\linewidth}L{.28\linewidth}L{.16\linewidth}}
\caption{Headline status at commit \CommitShort.}
\label{tab:frontier-catalogue}\\
\toprule
Project & Established surface & Remaining or deliberately excluded surface & Frontier type \\
\midrule
\endfirsthead
\toprule
Project & Established surface & Remaining or deliberately excluded surface & Frontier type \\
\midrule
\endhead
JacobianConjecture &
Explicit dimension-three polynomial map; determinant \(-2\); integral/rational
collisions; no inverse; scaling, stabilization, and degree-lowered variants;
Lean and Rocq audits. &
No claim about the surviving dimension-two problem; broader consequences are
outside this project. &
Scope boundary.\\

PolynomialFormulas: classical core &
Degrees one through four; exact radical relations; executable Gaussian
radical syntax/approximation; explicit and generic Abel--Ruffini surfaces;
quintic and sextic decision frameworks. &
A general executable radical formula for arbitrary solvable quintics is not
claimed merely from the decision theorem. &
Mathematical/API boundary.\\

PolynomialFormulas: Lazard/Dummit &
Large portions of finite group classification, resolvent bridges, sparse
coefficient tables, term/tail certificates, and a deterministic module plan
are checked at the recorded checkpoint. &
Later table terms, root-state data, and remaining planned certificate parents
identified by the formalization report/commit checkpoint. &
Build/checkpoint frontier.\\

PowerTowers A198683 &
Analytic and finite certificates for accepted small values through the stated
OEIS surface; substantial \(n=12\) symbolic, overflow, magnitude, and endpoint
work; near-\(1\) endpoint bundle unconditional. &
The main accepted data field deliberately withholds the final \(n=12\) value
until the complete proof-quality certificate boundary is satisfied. &
Formalization frontier.\\

SquaredSquare &
Order-21 perfect dissection, scaling, and elementary lower bound excluding
orders at most six; exact bracket \(7\le m\le21\). &
Historical exhaustive exclusion of orders 7--20, represented by
\code{DuijvestijnSearchClaim}. &
Finite-search frontier.\\

KlarnerConstant &
All seventeen geometric recurrences, exact supersolution
\(9047/2000\), pointwise bound, and asymptotic growth endpoint are
unconditional. &
Human transcription of paper diagrams remains a provenance boundary, not an
open theorem premise. &
Source boundary only.\\

BusyBeaver &
Machine semantics, known-value API, and opt-in small-state certificate
families; explicit generators and audits. &
Some BB4 root/classification endpoint remains conditional or checkpointed as
documented; expensive certificates remain outside the facade. &
Build/formalization frontier.\\

CombinatoryLogic &
Local lambda/SK/SKI/Iota simulations, coding, confluence components, and
universality; external computability connections. &
No full-abstraction, mutual-inverse, or unproved cross-semantics composition
claim. &
Deliberate nonclaim.\\

TuringDegrees &
Lean quotient theory with core degree operations and cardinality results;
Rocq setoid theory with additional jump, hierarchy, and incomparability
directions; explicit coverage ledger. &
Known deeper existence/classification theorems and exact cross-prover parity
remain deliberately incomplete; some Rocq endpoints retain logical premises. &
Coverage frontier.\\

FoundationPort / Modal &
A large independently checked Rocq theorem surface and a 455-module source
inventory with reviewed port status. &
Modules or specialized systems marked \code{partial}/\code{unported}; green
subtrees do not imply full Foundation parity. &
Port frontier.\\

FirstOrder &
Independent fixed-language soundness, completeness, model existence,
consistency/model equivalence, and compactness. &
Generic-language work is delegated to a separate Mathlib-backed facade, not an
unfinished local generalization. &
Deliberate scope.\\

PA ListCoding &
Thirty audited formulas/predicates, including list operations and
\(\varepsilon_0\)-related arithmetic, in independent Lean/Rocq encodings. &
Numeric code equality across assistants is not claimed. &
Contract boundary.\\

Bounded PA consistency &
Numeralwise all-occurrences bounded-consistency scheme in both assistants;
stronger uniform endpoint available in Lean at this snapshot. &
Rocq strongest uniform endpoint retains explicit compiler premises; exact
parity is not claimed. &
Cross-prover frontier.\\

Bounded ZFC consistency &
Numeralwise all-occurrences bounded-consistency theorem proved over the local
FirstOrder/ZF stack with no surviving project premise. &
A single internal universal sentence is deliberately not the theorem. &
Specification boundary.\\

ZFCinPA &
Definability layers, exact target, eight-field certificate reduction,
standard-index successor implications, and multiple semantic bridges. &
Nonstandard-index successor implication
\bcode{ZFCSuccessorImplicationsInAllModels}; final uniform PA endpoint remains
conditional. &
Mathematical/formalization frontier.\\

PresburgerArithmetic &
Lean end-to-end Cooper elimination and executable decision; Rocq normalized
one-variable correctness layer. &
Full identical Rocq end-to-end executable parity. &
Port frontier.\\
\bottomrule
\end{longtable}
\end{landscape}

\section{How to read a checkpoint commit}

A checkpoint commit in a generated proof program should answer:

\begin{enumerate}
  \item Which leaf certificates are kernel-green?
  \item Which parent certificates compose them?
  \item Which semantic bridge consumes the parent?
  \item Which files are generated but not yet reachable?
  \item What is the next dependency-ordered frontier?
  \item Which report or manifest was regenerated?
\end{enumerate}

The pinned commit's Lazard message does this with term/tail numbers and a
563-module root plan.  Those numbers are not theorem names; they are
coordinates in a deterministic build DAG.  A future walkthrough of a later
commit should replace them rather than treating them as permanent
architecture.

\section{Status can regress without a theorem becoming false}

Several ordinary changes can make a previous endpoint unavailable:

\begin{itemize}
  \item an upstream Mathlib or Rocq theorem changes shape;
  \item a generated shard is renamed without updating its parent;
  \item a facade stops importing the endpoint;
  \item a toolchain upgrade makes a formerly feasible kernel reduction exceed
        resources;
  \item a coverage ledger tightens its parity criterion;
  \item a specification audit reveals that an earlier statement was weaker
        than intended.
\end{itemize}

This is why ProveIt keeps theorem status near build and trust information.
``Was proved in some earlier commit'' and ``is part of the current public
surface'' are different facts.

\section{Recommended status vocabulary}

For future modules, the following terms match existing good practice:

\begin{longtable}{L{.23\textwidth}L{.66\textwidth}}
\caption{Precise status terms.}
\label{tab:status-vocabulary}\\
\toprule
Term & Use \\
\midrule
\endfirsthead
\toprule
Term & Use \\
\midrule
\endhead
\emph{Unconditional} &
The literal endpoint has no project-specific hypothesis; imported logical
axioms are reported separately.\\
\emph{Conditional} &
The endpoint theorem has an explicit premise representing unfinished or
external mathematics.\\
\emph{Numeralwise} &
A meta-level family of proofs indexed by each standard natural number.\\
\emph{Uniform} &
One object-language or model-internal theorem quantifies over the index.\\
\emph{Kernel certificate} &
Concrete computation is reduced by the kernel or converted to a kernel proof
term.\\
\emph{VM/native checked} &
A faster evaluator is part of the operational trust path.\\
\emph{Generated candidate} &
Data were produced externally but do not become evidence until a formal
checker accepts them.\\
\emph{Focused green} &
A named subset of modules built under a stated command/toolchain.\\
\emph{Ported} &
A reviewed theorem surface exists in the target assistant, under documented
encoding and assumptions.\\
\emph{Partial} &
A meaningful theorem slice exists but module-level parity is not claimed.\\
\emph{Deliberate nonclaim} &
A nearby stronger property is explicitly outside the theorem surface.\\
\bottomrule
\end{longtable}

\begin{statusbox}
All entries in this chapter describe commit \CommitShort.  They are not
predictions about the current branch after that commit.  The commit-pinned
READMEs and audits are the authoritative sources for this walkthrough.
\end{statusbox}

\begin{sourcebox}
\RepoFile{README.md};
\RepoFile{Algebra/PolynomialFormulas/README.md};
\RepoFile{Combinatorics/ExpressionEnumeration/PowerTowers/README.md};
\RepoFile{Combinatorics/SquaredSquare/README.md};
\RepoFile{Computability/BusyBeaver/README.md};
\RepoFile{Computability/TuringDegrees/COVERAGE.md};
\RepoFile{Logic/FoundationPort/README.md};
\RepoFile{Logic/Interpretability/ZFCinPA/README.md};
\RepoFile{Logic/PeanoArithmetic/BoundedConsistency/README.md};
\RepoFile{SetTheory/BoundedConsistency/README.md}.
\end{sourcebox}

\chapter{Architectural assessment and technical recommendations}
\label{chap:assessment}

\section{What the repository already does exceptionally well}

\subsection{It names proof boundaries}

The strongest recurring practice is the explicit separation among:

\begin{itemize}
  \item a theorem proved by ordinary elaboration;
  \item a theorem closed by kernel computation;
  \item a VM/native computation boundary;
  \item an externally generated candidate;
  \item a source-transcription boundary;
  \item a retained mathematical premise.
\end{itemize}

Many formal repositories make these distinctions only discoverable by reading
proof terms.  ProveIt often puts them in the README, module name, and audit.

\subsection{It avoids false cross-prover symmetry}

TuringDegrees, PA list coding, bounded consistency, and CombinatoryLogic all
document where Lean and Rocq use different representations or prove different
surfaces.  This is more trustworthy than forcing a two-column table in which
every row is marked ``done'' despite hidden premises.

\subsection{It scales exact computation by sharding}

The Lazard/Dummit tree shows that kernel-checked exact algebra can be scaled
far beyond one monolithic reduction.  Term, row, tail, and merge certificates
localize failures and let the import graph act as a proof DAG.

\subsection{It separates reusable conditional theorems from unconditional endpoints}

Klarner's \code{Main} and \code{GeometricComplete} exemplify good layering:
the abstract recurrence theorem remains reusable, while the concrete project
also discharges its hypotheses.  SquaredSquare and ZFCinPA make a different
choice where the premise remains open and keep it in the public statement.

\subsection{It treats status ledgers as artifacts}

FoundationPort's 455-module inventory and TuringDegrees' claim ledger are
first-class maintenance tools.  They prevent local success from being
misreported as full port completion.

\section{Structural risks}

\subsection{Generated-source discoverability}

Tens of thousands of Lean files make filesystem browsing a poor index.
Module names encode local structure, but there is no single machine-readable
map from:

\[
\text{headline theorem}
\to \text{semantic bridge}
\to \text{certificate roots}
\to \text{generated leaves}
\to \text{generator command}.
\]

The formalization reports fill part of this gap for Lazard, but each large
generator currently establishes its own conventions.

\subsection{Silent build no-ops}

The ZFCinPA example demonstrates that a plausible \code{lake build +Module}
command can schedule no work.  In a repository where ``file exists'' and
``target exists'' are different, this is a serious validation hazard.

\subsection{Documentation can drift faster than theorem types}

A README may still describe an earlier checkpoint after generated modules or
endpoint estimates have advanced.  The A198683 \(n=12\) subtree, for example,
contains stronger intermediate results than a simple ``not yet accepted''
summary conveys.  Human status prose needs some generated support.

\subsection{Resource feasibility is implicit}

Scoped heartbeats and recursion limits are visible in source, but a reviewer
cannot easily answer:

\begin{itemize}
  \item expected wall-clock time;
  \item peak memory;
  \item whether parallel workers are safe;
  \item which evaluator was used in the last green build;
  \item which modules form the expensive critical path.
\end{itemize}

This matters for reproducibility on machines smaller than the original
development host.

\subsection{Root validation is not visibly automated}

At the pinned commit there is no repository-level GitHub Actions workflow
tree.  Project READMEs provide exact commands, but no checked-in CI matrix
shows which tiers run on every change, nightly, or only manually.  With
multiple toolchains and very large certificates, one all-or-nothing CI job
would be impractical; the absence of a tiered machine-readable matrix remains
a maintenance risk.

\subsection{Submodule and compatibility complexity}

Rocq builds combine submodules, opam dependencies, local wrappers, and a
Windows compatibility overlay.  This is technically justified but increases
the chance that a successful build depends on an unrecorded mutable switch or
patched working tree.

\section{Recommendation 1: a machine-readable project registry}

Add a root registry such as \code{projects.toml} generated or reviewed
alongside \code{lakefile.toml}.  One record per project could contain:

\begin{lstlisting}[language={},caption={Illustrative project-registry schema.}]
[project.SquaredSquare]
subject = "Combinatorics"
lean_facade = "SquaredSquare"
lean_audit = "SquaredSquare.Audit"
rocq_root = "SquaredSquare"
readme = "Combinatorics/SquaredSquare/README.md"
status = "unconditional-bracket; sharp-minimality-conditional"
trust = ["lean-kernel-decide", "rocq-vm-compute", "manual-transcription"]
tier = "medium"
generated = false
\end{lstlisting}

For a generated project, add generator, manifest, certificate roots, and
resource tier.  This registry need not replace Lake or \code{_CoqProject}; it
would join their build facts with human status/provenance facts.

Benefits:

\begin{itemize}
  \item generate a repository project index and status table;
  \item validate that every facade/audit path exists;
  \item detect orphan generated modules;
  \item drive tiered CI;
  \item keep README summaries synchronized.
\end{itemize}

\section{Recommendation 2: target preflight}

Provide a script that fails if a requested Lean module is not scheduled:

\begin{lstlisting}[language=PowerShell]
pwsh Tools/Build-CheckedLeanTarget.ps1 `
  -Module ZFCinPA.NumeralOmega `
  -RequireWork
\end{lstlisting}

The wrapper can query Lake's build graph if available, or fall back to direct
\code{lake env lean} elaboration.  It should print the physical source path,
selected workspace, and output artifact.

Similarly, a Rocq wrapper can verify that each requested logical module is
present in the generated Makefile before invoking \code{make}.

\section{Recommendation 3: tiered validation}

A practical validation matrix could have four tiers.

\begin{longtable}{L{.17\textwidth}L{.35\textwidth}L{.37\textwidth}}
\caption{Suggested validation tiers.}
\label{tab:validation-tiers}\\
\toprule
Tier & Contents & Trigger \\
\midrule
\endfirsthead
\toprule
Tier & Contents & Trigger \\
\midrule
\endhead
Smoke &
Root facade parse/elaboration, microproject facades, registry validation,
format/provenance checks. &
Every change.\\
Core &
All hand-written endpoints and audits, excluding heavy generated certificate
families. &
Every change or merge queue.\\
Certificate &
Changed generated leaves, their transitive certificate parents, semantic
bridges, and audits. &
Path-aware on generated/support changes.\\
Full &
Root Lean default targets, all Rocq registered sources, all audits and
kernel-checks under pinned environments. &
Nightly/manual release checkpoint.\\
\bottomrule
\end{longtable}

A path-to-target manifest from Recommendation 1 would let CI select
certificate ancestors deterministically rather than rebuilding the entire
DAG.

\section{Recommendation 4: generated artifact manifests}

Each generator should emit a JSON/TOML manifest with:

\begin{itemize}
  \item generator source hash and version;
  \item input hashes;
  \item module/file hashes;
  \item dependency edges;
  \item semantic range (term numbers, rows, group classes);
  \item expected checker declaration;
  \item canonicalization version;
  \item last focused-green toolchain and resource observations.
\end{itemize}

A checker script can compare the manifest with both source imports and the
filesystem.  The proof assistant still establishes mathematical truth; the
manifest establishes completeness and reproducibility of the generated DAG.

A Merkle root over leaf data would make it cheap to locate the first changed
shard and to state exactly which data a parent certificate was built against.

\section{Recommendation 5: generated status from literal theorem types}

For conditional projects, maintain a small Lean module whose only purpose is
to emit a structured status summary from \code{\#check} declarations.  A
script can parse or query the environment to verify that:

\begin{itemize}
  \item the endpoint still has the documented premise;
  \item the unconditional wrapper exists when claimed;
  \item an audit names the same declaration;
  \item no implicit parameter appeared.
\end{itemize}

This would have caught several classes of documentation drift without trying
to infer mathematical meaning automatically.

\section{Recommendation 6: cross-prover coverage schema}

Generalize the TuringDegrees/FoundationPort ledger into a common TSV or TOML
schema:

\begin{lstlisting}[language={}]
source_claim = "FirstOrder.classical_complete"
lean_decl = "LeanProofs.FirstOrder.classical_complete"
rocq_decl = "FirstOrder.ClassicalCompleteness.classical_complete"
relation = "independent-equivalent"
lean_assumptions = ["Classical.choice", "propext", "Quot.sound"]
rocq_assumptions = ["ClassicalChoice"]
notes = "Fixed fMem language; same derivability convention"
\end{lstlisting}

Allowed relations could include
\code{exact}, \code{stronger}, \code{weaker}, \code{contract-equivalent},
\code{complementary}, \code{conditional}, and \code{unported}.  A renderer can
produce human tables while a validator checks declaration existence.

\section{Recommendation 7: source-transcription checksums}

For projects based on paper diagrams or external codes, preserve:

\begin{itemize}
  \item image/page identifier;
  \item normalized textual transcription;
  \item checksum of the source artifact when licensing permits local storage;
  \item independent parser/replay output;
  \item reviewer initials/date for manual visual comparison.
\end{itemize}

This does not make visual transcription a kernel theorem, but it turns an
informal memory into an auditable provenance step.

\section{Recommendation 8: resource contracts}

Add a small metadata block per heavy target:

\begin{lstlisting}[language={}]
target = "PolynomialFormulas....Table0Tail22Certificate"
evaluator = "kernel-decide"
max_heartbeats = 20000000
max_rec_depth = 1000000
recommended_workers = 0
last_peak_memory_gib = 6.4
last_elapsed_seconds = 183
machine = "..."
\end{lstlisting}

The timing is not a correctness condition and will vary, but large regressions
become visible.  Contributors can select feasible targets for their hardware
and avoid accidentally launching several memory-heavy shards concurrently.

\section{Recommendation 9: stable public-facade tests}

For every root facade, keep a tiny consumer module outside the implementation
namespace that imports only the facade and checks the headline declarations.
This catches accidental reliance on implementation imports available only in
the monorepo root environment.

For example:

\begin{lstlisting}[language=Lean]
import JacobianConjecture

#check LeanProofs.JacobianCounterexample.jacobianConjecture_false_over_complex
\end{lstlisting}

The existing audits often serve this role, but a strict consumer test can avoid
importing internal audit helpers and prove that the public facade alone is
sufficient.

\section{Recommendation 10: preserve the current honesty}

The most important recommendation is conservative: do not sacrifice the
repository's explicit status language while automating it.  A generated green
badge cannot express the difference between:

\begin{itemize}
  \item an unconditional theorem;
  \item a complete reduction to an open premise;
  \item a numeralwise scheme;
  \item a kernel-green shard not yet connected to an endpoint;
  \item a theorem imported from an upstream library;
  \item an independent proof in another assistant.
\end{itemize}

Machine-readable metadata should make those distinctions harder to lose, not
replace them with one ``verified'' bit.

\section{Overall assessment}

ProveIt's architecture is unusually well suited to experimental formal
mathematics: it permits small direct proofs, large generated calculations,
independent ports, and conditional research frontiers in one repository while
usually keeping their evidence types legible.  Its main scaling problem is no
longer theorem representation; it is indexing and validating the enormous
proof DAG, build graph, and status graph as one coherent system.

The next engineering gains are therefore likely to come from manifests,
target preflight, tiered validation, and generated crosswalks rather than from
flattening projects into a more uniform source layout.

\begin{sourcebox}
\RepoFile{README.md};
\RepoFile{lakefile.toml};
\RepoFile{_CoqProject};
\RepoFile{Computability/TuringDegrees/COVERAGE.md};
\RepoFile{Logic/FoundationPort/README.md};
\RepoFile{Algebra/PolynomialFormulas/LazardQuinticFormalization.tex};
\RepoFile{Tools/Rocq92Compat/README.md}.
\end{sourcebox}

\chapter{Reading itineraries}
\label{chap:itineraries}

\section{A one-hour architectural tour}

For a fast but representative tour:

\begin{enumerate}
  \item \RepoFile{README.md}: read the trust-model and status conventions.
  \item \RepoFile{ProveIt.lean}: see the consumer surface.
  \item \RepoFile{lakefile.toml}: compare imported and registered libraries.
  \item \project{QuantifierCommutation}: read an entire microproject.
  \item \project{FirstOrder}: inspect the syntax--calculus--completeness ladder.
  \item \project{SquaredSquare}: trace data $\to$ checker $\to$ real geometry.
  \item One Dummit tail certificate: see generated sharding.
  \item \RepoFile{Computability/TuringDegrees/COVERAGE.md}: see cross-prover
        status done carefully.
\end{enumerate}

This route samples every major architectural species without entering the
largest generated files.

\section{To audit one headline theorem}

Use this route regardless of subject:

\begin{enumerate}
  \item Find the literal declaration in the audit.
  \item Print/check its fully elaborated type and assumptions.
  \item Identify its endpoint module and direct imports.
  \item Classify every direct dependency as hand-written theorem, imported
        theorem, checked computation, generated data, or explicit premise.
  \item Follow the one dependency that carries semantic meaning, not the
        largest file.
  \item Confirm the facade exposes the theorem.
  \item Run the focused build and audit; read evaluator/trust notes.
  \item Compare README paraphrase with literal theorem type.
\end{enumerate}

For the Jacobian result, this means determinant and collision lemmas.  For
Klarner, it means the concrete recurrence assembly and exact supersolution.
For ZFCinPA, it means the surviving successor-implication premise.

\section{To understand generated algebra}

Read:

\begin{enumerate}
  \item \code{ComputableDummitCoefficientsCore.lean};
  \item one small sparse term's substitution/evaluation theorem;
  \item one \code{Term...Data}/\code{Term...Certificate} pair;
  \item one \code{Tail...Data}/\code{Tail...Certificate} pair;
  \item the parent row/table certificate;
  \item the invariant-theory bridge consuming the table;
  \item the generator and formalization report last.
\end{enumerate}

This order starts with semantics and uses the generated source as an instance.
The opposite order produces a great deal of scrolling and very little
understanding.

\section{To understand formal metamathematics}

Read:

\begin{enumerate}
  \item FirstOrder \code{Fol} and \code{Calculus};
  \item FirstOrder \code{Relativization};
  \item ZF formula/model definitions;
  \item one arithmetized syntax project such as PA ListCoding;
  \item bounded PA or bounded ZFC proof-tree verifier;
  \item partial truth and rule soundness;
  \item ZFCinPA only after writing the standard/nonstandard gaps explicitly.
\end{enumerate}

Do not begin with the final long theorem statement.  Its identifiers refer to
translations and coded predicates whose quantifier levels determine the
mathematics.

\section{To learn certificate engineering}

Use three progressively richer examples:

\begin{description}
  \item[SquaredSquare] Small fixed data, Boolean grid checker, clean geometric
        lifting.
  \item[BusyBeaver] Executable operational semantics, traces, normalization
        symmetries, and global classification.
  \item[PolynomialFormulas] Generated exact algebra, multi-level sharding,
        and semantic invariant bridges.
\end{description}

The three expose different completeness obligations.  A placement checker must
cover the whole region; a trace checker must be lifted to all machines for a
maximum theorem; an algebraic checker must preserve evaluation and feed the
right resolvent criterion.

\section{To compare Lean and Rocq}

Start with a genuinely small paired project such as QuantifierCommutation.
Then inspect one contract-parallel project such as PA ListCoding.  Finally read
TuringDegrees or FoundationPort, where the coverage ledger is essential.

For every declaration pair, record:

\[
(\text{statement},\text{encoding},\text{assumptions},\text{provenance},
\text{operational semantics}).
\]

Only the first component is visible from theorem names.

\section{To contribute an independent result}

A low-risk path is:

\begin{enumerate}
  \item choose a project with a small facade and complete audit;
  \item add one theorem in a focused module over existing definitions;
  \item add a matching audit entry and README note;
  \item validate both local and root environments;
  \item only then consider a Rocq port or generated certificate.
\end{enumerate}

Good starting projects include the analysis identities, integer sums,
quantifier commutation, or small power-tower values.  Projects involving
nonstandard models or Lazard certificates have much higher architectural
entry cost even when an individual theorem appears elementary.

\section{To review mathematical status rather than code style}

Read in this order:

\begin{enumerate}
  \item status table/nonclaims;
  \item literal endpoint and audit;
  \item trust boundary;
  \item coverage ledger or residue;
  \item build command and pinned dependency versions;
  \item source transcription/provenance;
  \item implementation details.
\end{enumerate}

A beautifully factored proof of the wrong endpoint is less useful than a
messy but accurately scoped conditional theorem.  ProveIt's best project
documentation recognizes that priority.

\begin{sourcebox}
\RepoFile{README.md};
\RepoFile{ProveIt.lean};
\RepoFile{lakefile.toml};
\RepoFile{Logic/QuantifierCommutation/README.md};
\RepoFile{Logic/FirstOrder/README.md};
\RepoFile{Combinatorics/SquaredSquare/README.md};
\RepoFile{Computability/TuringDegrees/COVERAGE.md};
\RepoFile{Algebra/PolynomialFormulas/LazardQuinticFormalization.tex}.
\end{sourcebox}

\appendix
\part{Reference appendices}

\chapter{Root Lean workspace registry}
\label{app:lake-registry}

The following table is an exact transcription of the library names and source
roots in \RepoFile{lakefile.toml} at the pinned commit.  ``Root import'' means
a direct import in \RepoFile{ProveIt.lean}; it does not count transitive
imports.  ``Default'' means the library itself or named audit modules appear in
the root default-target list.

\begin{landscape}
\begin{longtable}{L{.24\linewidth}L{.49\linewidth}C{.10\linewidth}C{.10\linewidth}}
\caption[Root Lean library registry]{Root \code{lean_lib} registry.}
\label{tab:exact-lake-registry}\\
\toprule
Library & \code{srcDir} & Root import & Default \\
\midrule
\endfirsthead
\toprule
Library & \code{srcDir} & Root import & Default \\
\midrule
\endhead
\code{JacobianConjecture} & \RepoDir{Algebra/JacobianConjecture/Lean} & \yes & -- \\
\code{PolynomialFormulas} & \RepoDir{Algebra/PolynomialFormulas/Lean} & \yes & -- \\
\code{ProveIt} & repository root & -- & \yes \\
\code{TrigonometricIdentities} & \RepoDir{Analysis/TrigonometricIdentities/Lean} & \yes & -- \\
\code{ExponentialIdentities} & \RepoDir{Analysis/ExponentialIdentities/Lean} & \yes & -- \\
\code{PowerTowers} & \RepoDir{Combinatorics/ExpressionEnumeration/PowerTowers/Lean} & \yes & -- \\
\code{A158415} & \RepoDir{Combinatorics/ExpressionEnumeration/RadicalExpressions/A158415/Lean} & \yes & -- \\
\code{A290268} & \RepoDir{Combinatorics/DerivativeExpansions/A290268/Lean} & \yes & -- \\
\code{SquaredSquare} & \RepoDir{Combinatorics/SquaredSquare/Lean} & \yes & -- \\
\code{KlarnerConstant} & \RepoDir{Combinatorics/Polyominoes/KlarnerConstant/Lean} & \yes & -- \\
\code{BusyBeaver} & \RepoDir{Computability/BusyBeaver/Lean} & \yes & -- \\
\code{CombinatoryLogic} & \RepoDir{Computability/CombinatoryLogic/Lean} & \yes & -- \\
\code{TuringDegrees} & \RepoDir{Computability/TuringDegrees/Lean} & \yes & -- \\
\code{FirstOrder} & \RepoDir{Logic/FirstOrder/Lean} & \yes & -- \\
\bcode{ArbitraryLanguageCompactness} & \RepoDir{Logic/FirstOrder/Compactness/Lean} & \yes & -- \\
\code{QuantifierCommutation} & \RepoDir{Logic/QuantifierCommutation/Lean} & \yes & -- \\
\code{ShefferStroke} & \RepoDir{Logic/Propositional/ShefferStroke/Lean} & \yes & -- \\
\code{NaturalDeduction} & \RepoDir{Logic/Propositional/NaturalDeduction/Lean} & \yes & -- \\
\bcode{FiniteMatrixNoncharacterizability} & \RepoDir{Logic/Propositional/FiniteMatrixNoncharacterizability/Lean} & \yes & -- \\
\code{MonotonicityOfEntailment} & \RepoDir{Logic/Propositional/MonotonicityOfEntailment/Lean} & \yes & -- \\
\code{PrincipleOfExplosion} & \RepoDir{Logic/Propositional/PrincipleOfExplosion/Lean} & \yes & -- \\
\code{EquationalLogic} & \RepoDir{Logic/Equational/Common/Lean} & \yes & -- \\
\code{BooleanAlgebra} & \RepoDir{Logic/Equational/BooleanAlgebra/Lean} & \yes & -- \\
\code{PAHF} & \RepoDir{Logic/Interpretability/PAHF/Lean} & \yes & -- \\
\code{NoFiniteModel} & \RepoDir{Logic/PeanoArithmetic/NoFiniteModel/Lean} & \yes & -- \\
\code{BoundedPAConsistency} & \RepoDir{Logic/PeanoArithmetic/BoundedConsistency/Lean} & -- & \yes \\
\code{PAListCoding} & \RepoDir{Logic/PeanoArithmetic/ListCoding/Lean} & -- & \yes \\
\code{PAFiniteBasisReduction} & \RepoDir{Logic/PeanoArithmetic/NotFinitelyAxiomatizable/Lean} & \yes & -- \\
\code{Foundation} & \RepoDir{lib/FormalizedFormalLogic-Foundation} & -- & -- \\
\code{PAUndecidable} & \RepoDir{Logic/PeanoArithmetic/Undecidable/Lean} & -- & \yes \\
\code{PresburgerArithmetic} & \RepoDir{Logic/PresburgerArithmetic/Lean} & \yes & -- \\
\code{DiophantineEquations} & \RepoDir{NumberTheory/DiophantineEquations/Lean} & \yes & -- \\
\code{IntegerSums} & \RepoDir{NumberTheory/IntegerSums/Lean} & \yes & -- \\
\code{RationalEnumeration} & \RepoDir{NumberTheory/RationalEnumeration/Lean} & \yes & -- \\
\code{RiemannHypothesis} & \RepoDir{NumberTheory/RiemannHypothesis/PAStatement/Lean} & \yes & -- \\
\code{ZF} & \RepoDir{SetTheory/ZF/Lean} & \yes & -- \\
\code{BoundedZFCConsistency} & \RepoDir{SetTheory/BoundedConsistency/Lean} & -- & -- \\
\code{ZFCinPA} & \RepoDir{Logic/Interpretability/ZFCinPA/Lean} & -- & -- \\
\code{ClosureAxiomatization} & \RepoDir{SetTheory/ClosureAxiomatization/Lean} & \yes & -- \\
\bottomrule
\end{longtable}
\end{landscape}

The default-target list also names
\bcode{PAListCoding.Audit}, \bcode{PAUndecidable.Audit}, and
\bcode{BoundedPAConsistency.Audit}.  This is why the final column is marked for
the containing libraries even though audit modules are not separate
\code{lean_lib} records.

\chapter{Project and artifact index}
\label{app:project-index}

\begin{longtable}{L{.22\textwidth}L{.27\textwidth}L{.41\textwidth}}
\caption{Where to start for each major project.}
\label{tab:project-index}\\
\toprule
Project & Start path & First artifacts to inspect \\
\midrule
\endfirsthead
\toprule
Project & Start path & First artifacts to inspect \\
\midrule
\endhead
JacobianConjecture &
\RepoDir{Algebra/JacobianConjecture} &
\code{README.md}; Lean \code{Counterexample.lean}; \code{Audit.lean};
Rocq \code{Counterexample.v}.\\
PolynomialFormulas &
\RepoDir{Algebra/PolynomialFormulas} &
\code{README.md}; Lean facade; \code{Basic/Cubic/Quartic};
\code{ComputableDummitCoefficientsCore}; formalization report.\\
ExponentialIdentities &
\RepoDir{Analysis/ExponentialIdentities} &
Facade; \code{TinyExponentTower.lean}; Rocq counterpart.\\
TrigonometricIdentities &
\RepoDir{Analysis/TrigonometricIdentities} &
Facade; arctangent and golden-ratio modules; audit.\\
PowerTowers &
\RepoDir{Combinatorics/ExpressionEnumeration/PowerTowers} &
Lean facade; \code{Core}; \code{A198683}; one small-value proof; \(n=12\)
symbolic/endpoints.\\
A158415 &
\RepoDir{Combinatorics/ExpressionEnumeration/RadicalExpressions/A158415} &
Facade; expression syntax; accepted-value certificates; research/support
material.\\
A290268 &
\RepoDir{Combinatorics/DerivativeExpansions/A290268} &
Facade; coefficient recurrence; auxiliary OEIS material.\\
SquaredSquare &
\RepoDir{Combinatorics/SquaredSquare} &
\code{README}; \code{Basic}; \code{Duijvestijn}; \code{Minimality};
\code{MinimalOrder}; audit.\\
KlarnerConstant &
\RepoDir{Combinatorics/Polyominoes/KlarnerConstant} &
\code{README}; \code{Polyomino}; \code{Patterns};
\code{SeededPartition}; one geometry family; \code{Certificate};
\code{GeometricComplete}.\\
BusyBeaver &
\RepoDir{Computability/BusyBeaver} &
\code{README}; facade; \code{Core}; \code{KnownValues}; smallest opt-in
classification; generator only afterward.\\
CombinatoryLogic &
\RepoDir{Computability/CombinatoryLogic} &
\code{README}; \code{Lambda}; \code{LambdaToSK}; \code{SKI};
\code{IotaToLambda}; \code{Universality}; audit.\\
TuringDegrees &
\RepoDir{Computability/TuringDegrees} &
\code{README}; \code{COVERAGE.md}; Lean quotient core; Rocq jump/hierarchy
modules; build wrapper.\\
QuantifierCommutation &
\RepoDir{Logic/QuantifierCommutation} &
Read the entire project: README, two proof modules, audit, and focused build.\\
FirstOrder &
\RepoDir{Logic/FirstOrder} &
README; \code{Fol}; \code{Calculus}; \code{Completeness};
\code{ClassicalCompleteness}; \code{Compactness}; audit.\\
ArbitraryLanguageCompactness &
\RepoDir{Logic/FirstOrder/Compactness} &
Facade and audit; identify the exact Mathlib theorem specialized.\\
FoundationPort &
\RepoDir{Logic/FoundationPort} &
README; \code{ported.tsv}; \code{coverage.sh}; Rocq \code{Audit.v}.\\
Modal &
\RepoDir{Logic/Modal} &
README/module coverage table; generic Kripke layer; one named-system chain;
audit.\\
PAHF &
\RepoDir{Logic/Interpretability/PAHF} &
Lean/Rocq facades and audits; translation/model modules.\\
ZFCinPA &
\RepoDir{Logic/Interpretability/ZFCinPA} &
README status table; exact target; eight-field certificate; standard-index
bridge; conditional endpoint; companion audits.\\
PA ListCoding &
\RepoDir{Logic/PeanoArithmetic/ListCoding} &
README; code representation; formula registry; functionality theorems;
audit.\\
Bounded PA consistency &
\RepoDir{Logic/PeanoArithmetic/BoundedConsistency} &
README; proof-tree/verifier; all-occurrences bound; partial truth; soundness;
uniform endpoint/audit.\\
NoFiniteModel &
\RepoDir{Logic/PeanoArithmetic/NoFiniteModel} &
Facade; semantic PA axioms; successor injection/non-surjection; audit.\\
PA finite-basis reduction &
\RepoDir{Logic/PeanoArithmetic/NotFinitelyAxiomatizable} &
Statement and reduction modules before internal coding.\\
PA undecidability &
\RepoDir{Logic/PeanoArithmetic/Undecidable} &
Exact decision problem, reduction, computability bridge, audit.\\
PresburgerArithmetic &
\RepoDir{Logic/PresburgerArithmetic} &
README; \code{Syntax}; \code{NormalForm}; \code{Cooper};
\code{Elimination}; \code{Decision}; audit.\\
DiophantineEquations &
\RepoDir{NumberTheory/DiophantineEquations} &
README; tiny Lean wrapper; independent Rocq descent.\\
IntegerSums &
\RepoDir{NumberTheory/IntegerSums} &
Facade and paired finite-sum proofs.\\
RationalEnumeration &
\RepoDir{NumberTheory/RationalEnumeration} &
README; forward Calkin--Wilf map; Euclidean inverse; bijection; audit.\\
RiemannHypothesis &
\RepoDir{NumberTheory/RiemannHypothesis/PAStatement} &
Exact arithmetic sentence and standard-model correspondence.\\
ZF &
\RepoDir{SetTheory/ZF} &
README; Lean facade; audit; search \code{Zf.lean} by declaration.\\
BoundedZFCConsistency &
\RepoDir{SetTheory/BoundedConsistency} &
README; formula rank; restricted proof predicate; partial satisfaction;
soundness; endpoint/audit.\\
ClosureAxiomatization &
\RepoDir{SetTheory/ClosureAxiomatization} &
README; forward/reverse semantic directions; deductive equivalence; article.\\
Rocq92Compat &
\RepoDir{Tools/Rocq92Compat} &
README; version pins; patch application; launchers/shims; no theorem claims.\\
\bottomrule
\end{longtable}

\chapter{Command cookbook}
\label{app:commands}

These commands are templates.  Project READMEs remain authoritative for
project-specific versions, worker settings, and dependencies.

\section{Repository state}

\begin{lstlisting}[language=PowerShell]
git rev-parse HEAD
git status --short
git submodule status
git diff --stat
git diff --check
\end{lstlisting}

\section{Lean bootstrap and broad builds}

\begin{lstlisting}[language=PowerShell]
lean --version
lake --version
lake exe cache get

# Root default targets.
lake build

# Consumer facade only.
lake build ProveIt
\end{lstlisting}

\section{Focused Lean}

\begin{lstlisting}[language=PowerShell]
# Directly elaborate a physical source file.
lake env lean path/to/Module.lean

# Build a registered library and audit.
lake build Project Project.Audit

# Use a project-local workspace.
lake --dir path/to/Project/Lean build
\end{lstlisting}

\section{Lean inspection}

\begin{lstlisting}[language=Lean]
import Project.Facade

#check Project.headlineTheorem
#print Project.headlineTheorem
#print axioms Project.headlineTheorem

set_option pp.all true in
#check Project.headlineTheorem
\end{lstlisting}

\section{Lean resource controls}

\begin{lstlisting}[language=Lean]
set_option maxHeartbeats 20000000 in
set_option maxRecDepth 1000000 in
theorem heavyCertificate : ... := by
  decide
\end{lstlisting}

\begin{lstlisting}[language=PowerShell]
# Serialize Lake jobs where project documentation recommends it.
$env:LEAN_NUM_THREADS = "0"
lake build HeavyTarget
Remove-Item Env:LEAN_NUM_THREADS -ErrorAction SilentlyContinue
\end{lstlisting}

\section{Rocq focused build}

\begin{lstlisting}[language=PowerShell]
coqc -Q path/to/Coq LogicalRoot path/to/Coq/Core.v
coqc -Q path/to/Coq LogicalRoot path/to/Coq/Endpoint.v
coqc -Q path/to/Coq LogicalRoot path/to/Coq/Audit.v

coqchk -silent -Q path/to/Coq LogicalRoot `
  LogicalRoot.Core LogicalRoot.Endpoint LogicalRoot.Audit
\end{lstlisting}

\section{Rocq root build}

\begin{lstlisting}[language=PowerShell]
git submodule update --init `
  lib/Coq-Synthetic-Computability `
  lib/MathComp-Abel

opam install --yes --deps-only `
  ./lib/MathComp-Abel/coq-mathcomp-abel.opam

pwsh -NoProfile -File `
  Computability/TuringDegrees/Coq/BuildSyntheticComputability.ps1

rocq makefile -f _CoqProject -o Makefile.coq
make -f Makefile.coq
\end{lstlisting}

\section{Rocq inspection}

\begin{lstlisting}[language=Rocq]
Check HeadlineTheorem.
Print HeadlineTheorem.
Print Assumptions HeadlineTheorem.
\end{lstlisting}

\section{Foundation port coverage}

\begin{lstlisting}[language=PowerShell]
bash Logic/FoundationPort/coverage.sh `
  > foundation-coverage.tsv
\end{lstlisting}

The script reports totals to standard error and emits one row per pinned
Foundation source module.

\section{Generated output comparison}

\begin{lstlisting}[language=PowerShell]
$Out = Join-Path $env:TEMP "proveit-generated"
Remove-Item $Out -Recurse -Force -ErrorAction SilentlyContinue
New-Item $Out -ItemType Directory | Out-Null

python path/to/generator.py --output-root $Out

git diff --no-index -- `
  path/to/checked-in/generated `
  $Out
\end{lstlisting}

\section{Find reachability}

For a suspicious generated module:

\begin{lstlisting}[language=PowerShell]
# Find direct imports.
git grep -n "import .*ModuleName"

# Find facade and audit references.
git grep -n "ModuleName" -- "*Facade*.lean" "*Audit*.lean" "*.lean"

# Elaborate directly even if no target exists.
lake env lean path/to/ModuleName.lean
\end{lstlisting}

\section{Clean rebuilds}

Use clean operations only after preserving diagnostics; a clean can erase the
evidence that distinguishes stale output from a source error.

\begin{lstlisting}[language=PowerShell]
lake clean
lake exe cache get
lake build Project Project.Audit

Remove-Item Makefile.coq -ErrorAction SilentlyContinue
rocq makefile -f _CoqProject -o Makefile.coq
make -f Makefile.coq Target.vo
\end{lstlisting}

\chapter{Trust and status checklist}
\label{app:trust-checklist}

\section{For a direct theorem}

\begin{enumerate}
  \item Locate the literal theorem type.
  \item Print all implicit parameters.
  \item Read \code{\#print axioms} or \code{Print Assumptions}.
  \item Identify imported headline theorems versus locally proved lemmas.
  \item Confirm the facade imports the endpoint.
  \item Confirm the audit names the endpoint.
  \item Check the README paraphrase against the theorem type.
\end{enumerate}

\section{For a computed theorem}

\begin{enumerate}
  \item Identify the data type and concrete candidate.
  \item Identify the executable checker.
  \item Identify the checker-soundness theorem.
  \item Identify the evaluator: kernel, VM, or native.
  \item Check integer/arithmetic domains for overflow or coercion.
  \item Confirm the endpoint uses checker soundness, not merely
        \code{check = true}.
  \item Confirm generated data are reachable from the endpoint.
\end{enumerate}

\section{For a generated certificate}

\begin{enumerate}
  \item Record generator source and command.
  \item Confirm deterministic/canonical output.
  \item Compare leaf data and certificate modules.
  \item Follow every parent shard to a semantic root.
  \item Verify no generated theorem is admitted or asserted.
  \item Verify regeneration does not leave orphan or stale modules.
  \item Build leaf, parent, semantic bridge, facade, and audit.
\end{enumerate}

\section{For a conditional theorem}

\begin{enumerate}
  \item Name the surviving premise exactly.
  \item Determine whether it is abstraction, source assumption, or open
        mathematical residue.
  \item Check whether a concrete module discharges it.
  \item Distinguish standard-index and all-model versions.
  \item Ensure README status says \emph{conditional}.
  \item Ensure no unconditional alias hides the premise.
\end{enumerate}

\section{For a Lean/Rocq pair}

\begin{enumerate}
  \item Compare literal statements, not names.
  \item Compare encodings and indexing.
  \item Compare classical/constructive assumptions.
  \item Compare evaluator/trust boundaries.
  \item Determine whether one proof is imported or independently local.
  \item Record exact parity relation in the coverage ledger.
  \item Kernel-check both public module sets.
\end{enumerate}

\section{For a source transcription}

\begin{enumerate}
  \item Identify source page, figure, table, or code.
  \item Preserve normalized transcription and checksum where possible.
  \item Independently replay basic consistency conditions.
  \item Distinguish transcription correctness from formal theorem correctness.
  \item Document any manual visual comparison.
\end{enumerate}

\chapter{Glossary}
\label{app:glossary}

\begin{description}
  \item[Audit module]
  A Lean or Rocq module that imports public declarations and prints/checks their
  types and assumptions.  It is executable documentation of the theorem
  surface.

  \item[Certificate]
  Concrete finite data intended to witness a decidable property.  A
  certificate is evidence only through a proved checker/meaning theorem.

  \item[Checker]
  An executable formal function or predicate validating a certificate.  Its
  soundness theorem connects acceptance to mathematical meaning.

  \item[Conditional endpoint]
  A theorem whose type retains a premise representing an unproved bridge,
  open claim, or intentionally abstract hypothesis.

  \item[Contract-level parity]
  Lean and Rocq use different representations but prove comparable public
  behavior such as decoding, functionality, or semantic equivalence.

  \item[Coverage ledger]
  A reviewed map from source claims/modules to target proof-assistant
  declarations and statuses.

  \item[Default target]
  A Lake target built by \code{lake build} with no explicit target.  It need
  not be directly imported by the root facade.

  \item[Direct module elaboration]
  Running \code{lake env lean path/to/File.lean}; useful when no Lake target
  exposes the file.

  \item[Facade]
  A small module whose primary role is importing the intended public surface
  of a project.

  \item[Focused green]
  A named subset of modules that successfully built under a recorded
  toolchain/command.  It is narrower than a full repository validation.

  \item[Generated data module]
  A source file containing literal exact data emitted by a support program.
  It should be consumed by a formal certificate theorem.

  \item[Kernel computation]
  Reduction performed within the proof assistant's trusted kernel, such as
  Lean \code{decide} on a closed decidable proposition.

  \item[Kernel checker]
  Lean's or Rocq's small trusted component that validates proof terms against
  the type theory.

  \item[Meaning theorem]
  A theorem that lifts a low-level checker result to the user-facing
  mathematical property.

  \item[Native/VM computation]
  Faster evaluation through a compiled or virtual-machine path.  It can remain
  integrated with the proof assistant while having a broader operational trust
  base than direct kernel reduction.

  \item[Numeralwise scheme]
  For every external standard natural \(n\), a separate formal derivation of
  the \(n\)-instance.  This is not one internal universally quantified
  sentence.

  \item[Port]
  A theorem surface independently recreated in another proof assistant.
  Exactness, generality, assumptions, and encoding must be specified.

  \item[Proof boundary]
  The point at which data or computation becomes a theorem: kernel term,
  verified checker, imported theorem, or explicit premise.

  \item[Residue]
  The ordered set of remaining obligations in an active formalization.

  \item[Root facade]
  \RepoFile{ProveIt.lean}, the broad user-facing Lean import.  It is neither
  the complete library registry nor the root default-target list.

  \item[Setoid presentation]
  A type plus an explicit equivalence relation used instead of quotienting
  immediately.  The Rocq Turing-degree project uses this style.

  \item[Shard]
  A bounded generated piece of a large certificate, often organized by term,
  row, tail, or class.

  \item[Source-transcription boundary]
  The human claim that formal data accurately reproduce an external diagram,
  table, or formula.  Formal checking can validate consequences of the data
  without proving the visual transcription itself.

  \item[Standard versus nonstandard index]
  A standard index is an external natural/numeral.  A nonstandard index is an
  arbitrary element satisfying arithmetic inside a nonstandard model.  The
  latter requires model-internal uniformity.

  \item[Uniform theorem]
  One formal statement quantifying internally over all indices, rather than a
  meta-level family of instances.

  \item[Vendored source]
  Third-party code pinned or copied under \code{lib}, with provenance and
  licensing recorded.

  \item[Worker count]
  Lake process concurrency.  It is independent of a declaration's heartbeat
  limit and can dominate peak memory for large certificate builds.
\end{description}

\chapter{Snapshot methodology and source notes}
\label{app:methodology}

\section{Pinned repository state}

All repository descriptions and hyperlinks refer to:

\begin{center}
\begin{tabular}{ll}
\toprule
Repository & \url{\RepoWeb}\\
Commit & \texttt{\CommitFull}\\
Short identifier & \texttt{\CommitShort}\\
Snapshot date & \SnapshotDate\\
Root tree & \texttt{43b4ea9da61e61551f5b416a1b7c82ac2103897e}\\
\bottomrule
\end{tabular}
\end{center}

The main branch changed during inspection, including a reorganization into the
current mathematical top-level taxonomy.  Pinning prevents a path from
silently referring to a later layout or theorem status.

\section{Inspection method}

The repository could not be cloned through the ordinary network path in the
document-production environment.  The exact commit was instead inspected
through GitHub's repository connector and Git-object/content APIs:

\begin{itemize}
  \item root commit, branch, and recursive tree objects;
  \item root README, manifests, toolchain, facade, submodule map, and Rocq
        project;
  \item domain and project READMEs;
  \item representative implementation and audit modules;
  \item generated data/certificate pairs;
  \item support/generator inventories;
  \item recent commit metadata and checkpoint message.
\end{itemize}

Content links in the PDF use the full commit SHA, so they are equivalent to
reading a local checkout at that commit for the cited paths.

\section{What was and was not independently checked}

The preparation process did:

\begin{itemize}
  \item inspect the exact source text and module trees described above;
  \item trace representative public endpoints to definitions, checker layers,
        generated data, and audits;
  \item compare the root facade, Lake registry, and Rocq registry;
  \item symbolically spot-check the displayed Jacobian determinant and
        collision outside the proof assistant as a sanity check;
  \item compile this \LaTeX{} document with Lua\LaTeX{};
  \item render and visually inspect the resulting PDF.
\end{itemize}

It did not:

\begin{itemize}
  \item perform a full local checkout and complete root \code{lake build};
  \item compile the entire Rocq federation and all submodules;
  \item regenerate every certificate family;
  \item independently peer-review the mathematical novelty or historical
        attribution of every project;
  \item certify that current main after \CommitFull{} has the same status.
\end{itemize}

Accordingly, phrases such as ``the project proves'' mean that the pinned source
contains the stated formal declaration and documented audit path.  This
walkthrough is not a substitute for running the project's build and kernel
checks on a clean checkout.

\section{Current-event context}

The Jacobian counterexample was an unusually recent event at the snapshot.
The source code's decisive identities are explicit and small, and the local
formalization states them directly.  For public context, the later
self-contained account
\href{https://arxiv.org/abs/2608.00222}{\emph{Counterexamples to the Jacobian
conjecture in dimensions greater than two}} describes the July 2026
dimension-three refutation and subsequent generalizations.  This external
context is not a dependency of the Lean/Rocq proof.

\section{Path and typography conventions}

Clickable source paths are rendered in monospaced blue text.  They point to
\code{blob} URLs for files and \code{tree} URLs for directories at the pinned
commit.  Module names without a path are printed in monospaced type.  Project
names are bold sans serif.

Code excerpts are normalized for readability.  In particular, some Lean
Unicode symbols are rendered as ASCII pseudocode in listings.  The surrounding
text and pinned source link are authoritative where a normalized excerpt
omits namespace qualification or notation details.

\section{How to update this walkthrough}

For a later commit:

\begin{enumerate}
  \item replace the commit and root-tree identifiers;
  \item diff \code{ProveIt.lean}, \code{lakefile.toml},
        \code{_CoqProject}, and \code{.gitmodules};
  \item regenerate language/file scale figures;
  \item revisit every status table and deliberate nonclaim;
  \item update generated-certificate frontiers from the new report/commit;
  \item re-run project coverage ledgers;
  \item check whether new facades, audits, or CI workflows were added;
  \item rebuild and visually inspect both source and PDF.
\end{enumerate}

A purely path-based update is unsafe.  In this repository, a theorem can keep
its filename while its premise, evaluator, facade reachability, or
cross-prover status changes.

\section{Closing perspective}

ProveIt is best read not as a pile of formal proofs but as a network of
evidence-producing components.  Some components derive a theorem directly.
Some interpret syntax, compile one calculus into another, or build a model.
Some check finite data emitted by search.  Some record exactly where a port or
uniformity argument stops.  The repository's coherence comes from making
those transitions explicit.

The practical invariant for contributors and reviewers is given the final
page to itself.

\clearpage
\thispagestyle{empty}
\vspace*{\fill}
\begin{center}
\begin{tcolorbox}[
  width=.78\textwidth,
  colback=sky,
  colframe=blue,
  boxrule=.8pt,
  arc=1.8mm,
  left=7mm,right=7mm,top=7mm,bottom=7mm
]
\centering
{\sffamily\bfseries\Large\color{navy}A closing invariant}\par
\vspace{4mm}
{\large
Every headline claim should have a visible path from statement, through
semantics and evidence, to the kernel---and every path that stops short should
stop in a named premise or status entry.\par}
\end{tcolorbox}
\end{center}
\vspace*{\fill}

\end{document}
